% =====================================================================================================
% language_hierarchy_section.tex — Language over the belief hierarchy: culture, sequencing and imagination
%
% DRAFT (2026-11-19). A proposal section: the architecture it describes is NOT yet implemented beyond the parts
% stated as existing. Source design note: docs/LANGUAGE_OVER_BELIEF_HIERARCHY.md.
% Include AFTER inner_speech_section.tex (it refers to that section's equations and modes):
%     % =====================================================================================================
% inner_speech_section.tex — Inner speech, mind-wandering and active inference
%
% DRAFT (2026-11-18). A proposal section: it describes an architecture that is NOT yet implemented,
% grounded in the existing language_agent (robocomp/active_inference). Include from main.tex with
%     % =====================================================================================================
% inner_speech_section.tex — Inner speech, mind-wandering and active inference
%
% DRAFT (2026-11-18). A proposal section: it describes an architecture that is NOT yet implemented,
% grounded in the existing language_agent (robocomp/active_inference). Include from main.tex with
%     % =====================================================================================================
% inner_speech_section.tex — Inner speech, mind-wandering and active inference
%
% DRAFT (2026-11-18). A proposal section: it describes an architecture that is NOT yet implemented,
% grounded in the existing language_agent (robocomp/active_inference). Include from main.tex with
%     \input{inner_speech_section}
% and add inner_speech.bib as a second resource:  \addbibresource{inner_speech.bib}
% All keys are new (none collides with the keys main.tex already cites).
% Packages used: amsmath, amssymb, enumitem (already loaded by main.tex) and array, which main.tex does not
% load yet: add \usepackage{array} to its preamble.
% =====================================================================================================

\section{Inner Speech as Epistemic Action}
\label{sec:inner_speech}

The robot described so far uses language in one direction only: a person speaks, and a language agent turns
the request into commands for the controller, grounded in a textual digest of the belief graph. This section
proposes to turn that capability inward. We argue that a language model embedded in an Active Inference
architecture can play the role that inner speech plays in human cognition: a serial, compositional medium in
which the agent poses questions to itself, generates hypotheses, rehearses plans and replays its past, and that
this role follows naturally from the epistemic component of expected free energy. The proposal is stated
together with the constraints that keep such a voice from becoming a source of confabulated belief, and with the
experiments that would falsify it.

\subsection{Motivation: language for thought before language for talk}
\label{sec:inner_speech:motivation}

Two traditions disagree on the origin of inner speech but agree on its function. On the internalist view,
human language evolved first as an internal system for structuring thought, and its use for communication is a
later externalisation \cite{berwick2016why}. On the developmental view, inner speech is social speech that the
child gradually internalises, abbreviates and turns into a tool for self-regulation \cite{vygotsky1986thought}.
Contemporary reviews treat inner speech as functionally involved in planning, problem solving, self-monitoring and
the maintenance of goals in working memory \cite{aldersonday2015inner}, and language more generally as a
cognitive technology that lets a learner manipulate discrete, recombinable tokens rather than raw percepts
\cite{clark1998magic}.

Large language models give a striking instance of the developmental route. They are trained exclusively on
communicative text, yet producing intermediate reasoning in that same medium before answering --- chain-of-thought
--- measurably improves their performance \cite{wei2022chain}. A system built for talk becomes better at thinking
by talking to itself. Our language agent sits at the same point of its own development: it was built to converse
with a person, and the question we ask is what it contributes once its speech is addressed to the robot itself.
Section~\ref{sec:inner_speech:why_lm} returns to what this origin implies for the internalist view.

\subsection{Inner questions as epistemic actions}
\label{sec:inner_speech:questions}

In Active Inference, policies are scored by expected free energy, whose epistemic term rewards the expected
reduction of uncertainty about hidden states \cite{friston2017process, parr2022active}. A question is the
linguistic form of exactly this quantity: it names a hidden variable whose value is uncertain and whose resolution
is expected to matter. We therefore model an inner question as a persistent object
\begin{equation}
    \mathcal{Q} = \big\langle\, e,\; \mathcal{H},\; q(h),\; \mathcal{P},\; b,\; s \,\big\rangle ,
    \label{eq:inner_question}
\end{equation}
where $e$ is the triggering event (for example an obstacle no concept agent explains), $\mathcal{H}$ a finite set
of named hypotheses (``a chair'', ``a box'', ``unknown''), $q(h)$ the current posterior over them, $\mathcal{P}$
the set of pending predictions, $b$ a budget in time and motion, and $s \in \{\text{open},\text{resolved},
\text{abandoned}\}$ its status. Questions form a stack in the sense of a \emph{question under discussion}
\cite{roberts2012information}: ``is it a chair?'' is a sub-question whose answer partially resolves the parent
``what is this?''. The question, not the language model's context window, carries the state of the reasoning from
one step to the next.

\paragraph{One step of the loop.}
Given $\mathcal{Q}$, the language model proposes a test: an epistemic action $a$ (turn towards the anomaly,
approach it, look with the front camera) together with the outcome each hypothesis predicts. The value of the test
is its expected information gain minus its cost,
\begin{equation}
    V(a) \;=\; \underbrace{H\!\left[q(h)\right] \;-\; \mathbb{E}_{p(o \mid a)}\!\left[ H\!\left[q(h \mid o, a)\right]
    \right]}_{I(h;\,o \mid a)} \;-\; c(a),
    \label{eq:inner_value}
\end{equation}
which is the negative of the epistemic part of expected free energy with an explicit action cost. The test is
executed only if $V(a) > 0$ and the action passes the same deterministic gate that already validates every
command the language agent issues. When the observation $o$ arrives, the posterior is updated by the concept
agents' own likelihoods,
\begin{equation}
    q(h \mid o, a) \;\propto\; p(o \mid h, a)\, q(h),
    \label{eq:inner_update}
\end{equation}
and the loop continues, spawns a sub-question, or terminates when
\begin{equation}
    H\!\left[q(h)\right] < \eta \quad\lor\quad \max_a V(a) \le 0 \quad\lor\quad b \text{ exhausted}.
    \label{eq:inner_stop}
\end{equation}
The stopping rule is what separates reasoning from rumination: a loop with no expected information to gain closes,
possibly with the answer ``unknown'', and its trigger is habituated for a refractory period.

\paragraph{The grounding constraint.}
The central design rule is that \emph{the language model proposes and perception disposes}. Text produced by the
model may introduce hypotheses, predictions and tests; it may never move a posterior. Only the likelihood
$p(o \mid h, a)$ in \eqref{eq:inner_update}, computed by the perceptual agents from sensor data, does. This is the
belief-side counterpart of the rule the language agent already enforces for actions, where every proposed command
is validated in C++ against the belief graph before it reaches the controller. Without it, a self-addressed
language loop is a closed system of text conditioned on text and converges on its own confabulations.

\paragraph{The proposer as a calibrated prior.}
The hypotheses the model names, and the probabilities it attaches to them, encode world knowledge the robot has
not acquired by itself (``chairs come with tables''). We treat this as a prior whose weight is earned. Let
$\pi_{\text{LM}}(h \mid x)$ be the proposer's distribution in context $x$ and $n_x(h)$ the robot's own counts of
resolved questions in that context. The prior entering \eqref{eq:inner_update} is
\begin{equation}
    q_0(h \mid x) \;\propto\; \pi_{\text{LM}}(h \mid x)^{\beta_x}\, \big(\alpha + n_x(h)\big),
    \label{eq:inner_prior}
\end{equation}
a tempered product of the proposer and a Dirichlet over experience, where $\beta_x \in [0,1]$ is set from the
proposer's measured calibration in that context (for instance from its Brier score on past questions). Early in
the robot's life the proposer dominates; as counts accumulate, the robot's own experience replaces borrowed
knowledge. The same conjugate form already governs the language agent's decision to ask a person or to assume
that they moved an object, where a Beta posterior over the person's past answers is compared against the costs of
asking and of being wrong.

\subsection{Attention, automaticity and the wandering mind}
\label{sec:inner_speech:modes}

A human driver on an easy road steers without attending to it; the mind is free to plan, imagine or recall, and is
recaptured the instant something unexpected happens \cite{smallwood2015science, christoff2016mind}. In Active
Inference this is a statement about precision: while prediction errors are small and the habitual policy is
reliable, attention --- the precision assigned to sensory prediction errors --- can be withdrawn from the task; a
surprise raises it again \cite{feldman2010attention}. The arbitration between habitual and deliberative control by
their relative uncertainty \cite{daw2005uncertainty} has the same structure.

The architecture already separates the two systems. The controller and the concept agents form an automatic layer
that drives, avoids obstacles and executes missions without the language model; the language agent is the
deliberative voice. We propose to make the allocation of the voice explicit through a scalar \emph{demand} signal,
computed on every digest cycle without the language model,
\begin{equation}
    D_t \;=\; \sum_{k} \pi_k\, \varepsilon_{k,t}^{2} \;+\; \lambda\, \mathbb{1}\!\left[\text{person active}_t\right],
    \label{eq:inner_demand}
\end{equation}
where $\varepsilon_{k,t}$ are normalised prediction errors reported by the agents (unexplained residual occupancy,
flips of existence or door beliefs, localisation uncertainty, stalled or failed mission steps, controller refusals)
and $\pi_k$ their precisions. Three modes follow from two thresholds with asymmetric hysteresis:
\begin{equation}
    m_t =
    \begin{cases}
        \textsc{focused}    & D_t > \theta_{\text{hi}},\\[2pt]
        \textsc{wandering}  & D_\tau < \theta_{\text{lo}} \;\; \forall\, \tau \in [t - T, t],\\[2pt]
        \textsc{monitoring} & \text{otherwise.}
    \end{cases}
    \label{eq:inner_modes}
\end{equation}
Entering \textsc{wandering} requires demand to stay low for a dwell time $T$; leaving it requires a single sample
above $\theta_{\text{hi}}$. The robot is slow to drift off and instant to snap back. Safety never depends on this
arbitration: the controller's reflexes act on their own, and the deliberative voice arrives after the reflex, as
awareness follows braking in the driver.

\paragraph{Focused mode} runs the question loop of Section~\ref{sec:inner_speech:questions} on the triggering
event. \paragraph{Monitoring mode} makes no language-model calls; the voice is primed but silent.

\paragraph{Wandering mode.}
Freed from the task, the voice is seeded with one recent experience chosen by salience: an episode that ended
unresolved, a failure, an object that changed, or an upcoming goal. Each step's output, together with an
associative retrieval from episodic memory, seeds the next, so that the chain drifts as spontaneous thought does
\cite{christoff2016mind}. The voice has access only to memory queries and to the architecture's imagination tools,
which simulate observations and outcomes on a copy of the belief graph; it can neither act nor write beliefs. Its
products are of three kinds: counterfactuals over past episodes (``had the door been open, the crossing would have
taken 40\,s less''), hypotheses to be tested later, and prospective plans (``if the door is shut on the way back,
the route through the next room is shorter''). Functionally this is the default network's role in simulating
past and future \cite{buckner2008brain}, and the offline replay and model reduction that Active Inference assigns
to rest and sleep \cite{friston2017curiosity, hobson2012waking}. Replay is prioritised by the expected gain of the
memory it revisits \cite{mattar2018prioritized}.

\paragraph{Source monitoring.}
Humans occasionally misremember imagined events as perceived ones \cite{johnson1993source}. The robot must not.
Everything produced while wandering is stored with an \emph{imagined} provenance, separate from the factual
episodic record, and enters the learning loop only as a hypothesis to be tested, never as evidence.

\subsection{What the return to reality gives the automatic layer}
\label{sec:inner_speech:return}

The transition from \textsc{wandering} to \textsc{focused} cancels the chain in flight (its state is kept, so that
it may resume if demand subsides) and gives the full voice to the surprise. It serves the automatic layer at three
timescales.
\begin{enumerate}
    \item \textbf{Decision.} A surprise means the automatic layer has left its domain of competence, where its only
    responses are to retry, stop or fail. The voice adds re-interpretation of the scene, arbitration between goals,
    re-planning through remembered topology, and recourse to a person.
    \item \textbf{Attention.} The voice can raise precision where it matters: orienting sensors, increasing a
    detector's rate in a region, or passing a hypothesis to the concept agents as a spatial prior. Such a prior has
    bounded weight; it may accelerate the recognition of an object that is present but may not create one that is
    not, preserving the grounding constraint.
    \item \textbf{Teaching.} Each resolved surprise is a labelled example the automatic layer could not have
    produced alone: a residual later identified as a chair is a false negative of the chair detector; a door that
    is repeatedly shut in the evening becomes a prior; a recovery that repeatedly succeeds becomes a branch of the
    controller's behaviour tree. The domain of competence grows, and the situations that require the voice become
    rarer, which is the classical progression from effortful to autonomous skill \cite{fitts1967human}.
\end{enumerate}
Plans prepared while wandering act as a cache: when the anticipated situation occurs, the focused voice verifies a
plan instead of composing one. The return ends with an explicit release --- spatial priors cleared, safety margins
restored, the automatic layer confirmed back within its competence --- without which every surprise would leave
the robot permanently vigilant.

\subsection{Relation to chain-of-thought reasoning}
\label{sec:inner_speech:cot}

Chain-of-thought is inner speech in the functional sense, but of a specific kind, and the comparison sharpens the
design (Table~\ref{tab:inner_cot}). Plain chain-of-thought is convergent and ungrounded: tokens condition on
tokens until an answer is produced, and the trace is discarded. Interleaving reasoning with action and observation
grounds it \cite{yao2023react}; our focused loop is of this kind, with the additional property that the
observation, through \eqref{eq:inner_update}, rather than the next token, determines where the reasoning goes.
The wandering voice is closer to plain chain-of-thought in its isolation from perception, but divergent rather
than convergent, interruptible, and persistent in its products.

\begin{table}[ht]
    \centering
    \small
    \begin{tabular}{|>{\raggedright\arraybackslash}p{0.19\linewidth}|>{\raggedright\arraybackslash}p{0.22\linewidth}|>{\raggedright\arraybackslash}p{0.24\linewidth}|>{\raggedright\arraybackslash}p{0.24\linewidth}|}
        \hline
                                & \textbf{Chain-of-thought} & \textbf{Focused voice} & \textbf{Wandering voice} \\ \hline
        Goal                    & answer a given question   & resolve a given surprise & none given; seed by salience \\ \hline
        Dynamics                & convergent                & convergent             & divergent, associative \\ \hline
        Contact with the world  & none                      & thought--action--observation & none (imagination only) \\ \hline
        Persistence             & discarded                 & question objects, beliefs, lessons & imagined store, hypotheses \\ \hline
        Interruptible           & no                        & yes                    & yes, resumable \\ \hline
        Latency budget          & n/a                       & $\sim$1\,s, no extended thinking & unbounded, extended thinking \\ \hline
    \end{tabular}
    \caption{Chain-of-thought compared with the two modes of the proposed inner voice.}
    \label{tab:inner_cot}
\end{table}

Two practical consequences follow. First, \emph{think when there is time, not when there is pressure}. In the
deployed language agent extended thinking is disabled, because a thinking-mode model spent 3\,000--6\,000 tokens
(40--80\,s) before each tool call; that is incompatible with the focused mode but costs nothing in the wandering
mode, where it can be enabled with a large budget and pre-empted at any moment. Slow reasoning is done in advance
and fast reasoning uses its results. Second, chain-of-thought traces are known not always to reflect the
computation that produced the answer \cite{turpin2023language}, as human verbal reports do not always reflect the
causes of behaviour \cite{nisbett1977telling}. The faithful record of the robot's reasoning is therefore the
structured one --- hypotheses, predictions, observations and posterior updates --- not the text. Finally, wandering
chains whose hypotheses were later confirmed by perception form a corpus of reasoning verified against the world,
which could be used to fine-tune the local model in the manner of self-taught reasoning \cite{zelikman2022star},
with perception, rather than a known answer, as the filter.

\subsection{Consequences for belief and learning}
\label{sec:inner_speech:learning}

The loop changes how beliefs evolve in two ways. Within an episode, uncertainty is reduced where surprise
occurred rather than where the sensors happened to point, which is efficient but leaves the rest of the map less
examined and makes the robot's own data a biased sample; a decaying novelty term \cite{schwartenbeck2019curiosity}
and importance weighting of learned statistics are required to correct this. Across episodes, every closed
question is a labelled example, and learning proceeds at four timescales (Table~\ref{tab:inner_learning}).

\begin{table}[ht]
    \centering
    \small
    \begin{tabular}{|>{\raggedright\arraybackslash}p{0.14\linewidth}|>{\raggedright\arraybackslash}p{0.37\linewidth}|>{\raggedright\arraybackslash}p{0.37\linewidth}|}
        \hline
        \textbf{Timescale} & \textbf{What changes} & \textbf{Where, in this architecture} \\ \hline
        Seconds   & posterior over the hypotheses of one question         & the question object \eqref{eq:inner_question} \\ \hline
        Days      & priors: where objects are, which doors are shut, what moves & long-term object and topological memory, as Dirichlet counts \\ \hline
        Weeks     & likelihoods: detector and concept models              & false negatives and positives labelled by resolved questions \\ \hline
        Slow      & structure: new categories                             & recurrent, mutually similar ``unknown'' resolutions named by the voice and proposed as a new concept \\ \hline
    \end{tabular}
    \caption{Learning driven by the inner-question loop.}
    \label{tab:inner_learning}
\end{table}

The last row is where the hypothesis of a language for thought is most directly engaged: a cluster of anomalies
that resists every existing concept becomes learnable once it is named. The loop also learns about itself: the
salience weights $\pi_k$ in \eqref{eq:inner_demand} are adjusted by the information each trigger actually yielded,
asking a person is weighed against looking, and loops that repeatedly resolve the same way are compiled into
reflexive rules, so that the number of language-model calls per hour should \emph{decrease} as competence grows.

\subsection{Why a language model: generation versus evaluation}
\label{sec:inner_speech:why_lm}

None of the above requires a person, so the question arises of what a language model contributes to an
architecture that already performs inference and action selection. Our answer is a division of labour. Expected
free energy evaluates policies and Bayesian updating weighs hypotheses, but both operate over a space that has
already been enumerated: each concept agent carries a fixed generative model, policies must be listed before they
can be scored, and Bayesian model reduction can prune structure but not invent it \cite{friston2017curiosity}.
Whatever falls outside the given models remains an unexplained residual. A language model is weak where this
machinery is strong --- its confidences are not calibrated posteriors and it has no access to the sensors --- and
strong where it is weak: it is a cheap, broad generator of hypotheses, explanations and plans. We therefore use it
only as a \emph{proposal distribution}, in a generate-and-test loop whose test is active inference:
\begin{itemize}
    \item \emph{Hypotheses outside the built-in models.} A residual blob can be proposed to be a bag, a toppled
    chair or a robot vacuum, none of which any concept agent represents; the question loop then tests them.
    \item \emph{Amortised policy search.} Instead of enumerating multi-step policies, the model writes a few
    candidate missions --- small programs in the mission language --- which the gate validates and expected free
    energy scores.
    \item \emph{Relational binding across agents.} The agents' representations (signed distance fields, occupancy
    grids, existence log-odds, episodic records) are mutually incompatible; a sentence such as ``the chair that was
    not next to the table in the room behind door~3 yesterday'' binds them in one query no single agent could pose.
    \item \emph{Explanation from few examples.} After a single surprise the model can propose which variable
    mattered (``the door is shut because it is evening''), turning one episode into a testable generalisation.
    \item \emph{Abstraction and naming.} Episodes are compressed into reusable lessons, and recurrent unknowns are
    named, which is what makes them candidates for a new concept (Table~\ref{tab:inner_learning}).
\end{itemize}
Language, rather than an arbitrary learned embedding, matters for three reasons: it is discrete and compositional,
so new hypotheses and plans are assembled from known parts; it is inspectable, so that deterministic gates can parse
and check it and people can audit it; and, once the robot thinks in it, communication with a person comes at no
further cost.

\paragraph{Where the knowledge comes from.}
This raises a difficulty for the motivating thesis. If the inner voice precedes the outer one, the common-sense
knowledge on which the proposer's value rests cannot have arrived through communication --- yet a language model
acquires all of it from human text, that is, from other people's outer voices. The difficulty dissolves once the
\emph{faculty} is distinguished from its \emph{content}. The internalist claim concerns the former: a
computational system for building unbounded hierarchical structures, available for thought before it is used for
talk \cite{hauser2002faculty, berwick2016why}. It makes no claim that what a person knows was acquired without
language. Human knowledge has three sources: core systems for objects, agents, number and space, present before
learning \cite{spelke2007core}; individual experience; and, once outer language exists, cumulative cultural
transmission, which supplies most of what an adult knows \cite{tomasello1999cultural}. Even on the internalist
view, the inner voice of a modern human is filled largely by the outer voices of others.

The architecture has the same three sources, in a different arrangement. The concept agents' hand-built generative
models play the role of \emph{core knowledge}: given before experience and limited to a few domains. The robot's
own counts $n_x(h)$ in \eqref{eq:inner_prior} are \emph{individual experience}. The language model is
\emph{culture}: knowledge distilled from a community's communication and inherited, for this individual, at birth.
The robot thus resembles a child born into a culture more than an early hominin, and its inner voice is Vygotskian
in origin: its format and its content were both learned from outer speech. Two consequences follow. First, the
proposal does not demonstrate the internalist thesis and should not be read as doing so; its claim is functional
--- an inner use of language has epistemic value independently of communication --- and holds whatever the origin of
the language system. Second, the format and the content of the voice can in principle be separated, which turns the
question into an experiment (P7 in Section~\ref{sec:inner_speech:predictions}): a compositional generator over the
robot's own concepts and mission language, trained only on the robot's experience, in the manner of library
learning in program induction \cite{ellis2021dreamcoder}, would supply the format without the culture.

\subsection{Mapping onto the existing architecture}
\label{sec:inner_speech:mapping}

Several components already exist in the language agent. Robot events (mission reports, failures, completions)
already wake the model through a queue, in turns restricted to perceiving, querying memory, stopping and
proposing, never acting. Every action already passes a deterministic gate. The decision to ask or to assume is
already decision-theoretic. Unexplained occupancy is already isolated by a residual agent that acts as the null
hypothesis of the scene, and episodic and long-term object memories are already joined read-only. Missing are the
persistent question objects \eqref{eq:inner_question}, a prediction evaluator implementing
\eqref{eq:inner_update} outside the language model, the demand signal and mode arbitration
\eqref{eq:inner_demand}--\eqref{eq:inner_modes}, the imagined store with its provenance, and the channel through
which bounded spatial priors reach the concept agents.

\subsection{Testable predictions}
\label{sec:inner_speech:predictions}

Consistent with the evaluation principle adopted throughout --- policies compared against each other inside one
closed loop, with body, sensors and environment fixed --- the proposal makes seven falsifiable predictions, each
tested in simulation by moving known objects and introducing unfamiliar ones:
\begin{enumerate}[label=\textbf{P\arabic*}]
    \item The unexplained residual mass per room decreases faster with the focused loop than without it.
    \item The proposer's calibration on resolved questions improves, and $\beta_x$ in \eqref{eq:inner_prior}
    decreases, as the robot's own counts accumulate.
    \item Detector false-negative rates fall after retraining on question-labelled examples.
    \item Inner questions per hour decrease over repeated exposure to the same environment while answer accuracy
    is maintained.
    \item Plans prepared while wandering reduce the latency from surprise to a validated response, and the time to
    notice a deliberately moved object does not grow despite habituation.
    \item With the language-model proposer replaced by a fixed hypothesis set (the existing concept classes and
    ``unknown'') and template-based missions, the loop resolves fewer questions about objects and failures outside
    the built-in models, and acquires new categories more slowly.
    \item A compositional generator trained only on the robot's own experience (format without culture) recovers
    part of the advantage in P6 within a familiar environment, and less of it the further a situation departs from
    what the robot has seen.
\end{enumerate}
A failure of P1 would indicate that the epistemic value of language-mediated questions does not exceed that of
the existing exploration; a failure of P4 would indicate that the loop does not transfer competence to the
automatic layer, which is its principal claim; and a failure of P6 would indicate that the language model adds
nothing a fixed hypothesis set does not, which would remove the case for including it. P7 separates the
contribution of language as a format for thought from that of the knowledge it carries.

\subsection{Limitations}
\label{sec:inner_speech:limitations}

The demand signal \eqref{eq:inner_demand} sums heterogeneous prediction errors whose precisions must themselves be
calibrated, and a poorly scaled term will either suppress the voice or keep it permanently engaged. A single
language model serves the person, the focused loop and the wandering loop; pre-emption keeps the person first but
bounds how much offline reasoning a busy robot can do. The tempered prior \eqref{eq:inner_prior} is a heuristic
combination rather than a derived posterior. The generator of P7 is substantially harder to build than the
ablation of P6, and may have to be approximated by a small model trained on the robot's logs and mission corpus,
whose pre-training would leak some culture back in. Finally, the analogy with human inner speech is functional, not
mechanistic: the proposal is that the same computational role --- serial, compositional self-questioning in the
service of epistemic value --- is useful for a robot, not that the robot's voice resembles the human one in any
further respect.

% and add inner_speech.bib as a second resource:  \addbibresource{inner_speech.bib}
% All keys are new (none collides with the keys main.tex already cites).
% Packages used: amsmath, amssymb, enumitem (already loaded by main.tex) and array, which main.tex does not
% load yet: add \usepackage{array} to its preamble.
% =====================================================================================================

\section{Inner Speech as Epistemic Action}
\label{sec:inner_speech}

The robot described so far uses language in one direction only: a person speaks, and a language agent turns
the request into commands for the controller, grounded in a textual digest of the belief graph. This section
proposes to turn that capability inward. We argue that a language model embedded in an Active Inference
architecture can play the role that inner speech plays in human cognition: a serial, compositional medium in
which the agent poses questions to itself, generates hypotheses, rehearses plans and replays its past, and that
this role follows naturally from the epistemic component of expected free energy. The proposal is stated
together with the constraints that keep such a voice from becoming a source of confabulated belief, and with the
experiments that would falsify it.

\subsection{Motivation: language for thought before language for talk}
\label{sec:inner_speech:motivation}

Two traditions disagree on the origin of inner speech but agree on its function. On the internalist view,
human language evolved first as an internal system for structuring thought, and its use for communication is a
later externalisation \cite{berwick2016why}. On the developmental view, inner speech is social speech that the
child gradually internalises, abbreviates and turns into a tool for self-regulation \cite{vygotsky1986thought}.
Contemporary reviews treat inner speech as functionally involved in planning, problem solving, self-monitoring and
the maintenance of goals in working memory \cite{aldersonday2015inner}, and language more generally as a
cognitive technology that lets a learner manipulate discrete, recombinable tokens rather than raw percepts
\cite{clark1998magic}.

Large language models give a striking instance of the developmental route. They are trained exclusively on
communicative text, yet producing intermediate reasoning in that same medium before answering --- chain-of-thought
--- measurably improves their performance \cite{wei2022chain}. A system built for talk becomes better at thinking
by talking to itself. Our language agent sits at the same point of its own development: it was built to converse
with a person, and the question we ask is what it contributes once its speech is addressed to the robot itself.
Section~\ref{sec:inner_speech:why_lm} returns to what this origin implies for the internalist view.

\subsection{Inner questions as epistemic actions}
\label{sec:inner_speech:questions}

In Active Inference, policies are scored by expected free energy, whose epistemic term rewards the expected
reduction of uncertainty about hidden states \cite{friston2017process, parr2022active}. A question is the
linguistic form of exactly this quantity: it names a hidden variable whose value is uncertain and whose resolution
is expected to matter. We therefore model an inner question as a persistent object
\begin{equation}
    \mathcal{Q} = \big\langle\, e,\; \mathcal{H},\; q(h),\; \mathcal{P},\; b,\; s \,\big\rangle ,
    \label{eq:inner_question}
\end{equation}
where $e$ is the triggering event (for example an obstacle no concept agent explains), $\mathcal{H}$ a finite set
of named hypotheses (``a chair'', ``a box'', ``unknown''), $q(h)$ the current posterior over them, $\mathcal{P}$
the set of pending predictions, $b$ a budget in time and motion, and $s \in \{\text{open},\text{resolved},
\text{abandoned}\}$ its status. Questions form a stack in the sense of a \emph{question under discussion}
\cite{roberts2012information}: ``is it a chair?'' is a sub-question whose answer partially resolves the parent
``what is this?''. The question, not the language model's context window, carries the state of the reasoning from
one step to the next.

\paragraph{One step of the loop.}
Given $\mathcal{Q}$, the language model proposes a test: an epistemic action $a$ (turn towards the anomaly,
approach it, look with the front camera) together with the outcome each hypothesis predicts. The value of the test
is its expected information gain minus its cost,
\begin{equation}
    V(a) \;=\; \underbrace{H\!\left[q(h)\right] \;-\; \mathbb{E}_{p(o \mid a)}\!\left[ H\!\left[q(h \mid o, a)\right]
    \right]}_{I(h;\,o \mid a)} \;-\; c(a),
    \label{eq:inner_value}
\end{equation}
which is the negative of the epistemic part of expected free energy with an explicit action cost. The test is
executed only if $V(a) > 0$ and the action passes the same deterministic gate that already validates every
command the language agent issues. When the observation $o$ arrives, the posterior is updated by the concept
agents' own likelihoods,
\begin{equation}
    q(h \mid o, a) \;\propto\; p(o \mid h, a)\, q(h),
    \label{eq:inner_update}
\end{equation}
and the loop continues, spawns a sub-question, or terminates when
\begin{equation}
    H\!\left[q(h)\right] < \eta \quad\lor\quad \max_a V(a) \le 0 \quad\lor\quad b \text{ exhausted}.
    \label{eq:inner_stop}
\end{equation}
The stopping rule is what separates reasoning from rumination: a loop with no expected information to gain closes,
possibly with the answer ``unknown'', and its trigger is habituated for a refractory period.

\paragraph{The grounding constraint.}
The central design rule is that \emph{the language model proposes and perception disposes}. Text produced by the
model may introduce hypotheses, predictions and tests; it may never move a posterior. Only the likelihood
$p(o \mid h, a)$ in \eqref{eq:inner_update}, computed by the perceptual agents from sensor data, does. This is the
belief-side counterpart of the rule the language agent already enforces for actions, where every proposed command
is validated in C++ against the belief graph before it reaches the controller. Without it, a self-addressed
language loop is a closed system of text conditioned on text and converges on its own confabulations.

\paragraph{The proposer as a calibrated prior.}
The hypotheses the model names, and the probabilities it attaches to them, encode world knowledge the robot has
not acquired by itself (``chairs come with tables''). We treat this as a prior whose weight is earned. Let
$\pi_{\text{LM}}(h \mid x)$ be the proposer's distribution in context $x$ and $n_x(h)$ the robot's own counts of
resolved questions in that context. The prior entering \eqref{eq:inner_update} is
\begin{equation}
    q_0(h \mid x) \;\propto\; \pi_{\text{LM}}(h \mid x)^{\beta_x}\, \big(\alpha + n_x(h)\big),
    \label{eq:inner_prior}
\end{equation}
a tempered product of the proposer and a Dirichlet over experience, where $\beta_x \in [0,1]$ is set from the
proposer's measured calibration in that context (for instance from its Brier score on past questions). Early in
the robot's life the proposer dominates; as counts accumulate, the robot's own experience replaces borrowed
knowledge. The same conjugate form already governs the language agent's decision to ask a person or to assume
that they moved an object, where a Beta posterior over the person's past answers is compared against the costs of
asking and of being wrong.

\subsection{Attention, automaticity and the wandering mind}
\label{sec:inner_speech:modes}

A human driver on an easy road steers without attending to it; the mind is free to plan, imagine or recall, and is
recaptured the instant something unexpected happens \cite{smallwood2015science, christoff2016mind}. In Active
Inference this is a statement about precision: while prediction errors are small and the habitual policy is
reliable, attention --- the precision assigned to sensory prediction errors --- can be withdrawn from the task; a
surprise raises it again \cite{feldman2010attention}. The arbitration between habitual and deliberative control by
their relative uncertainty \cite{daw2005uncertainty} has the same structure.

The architecture already separates the two systems. The controller and the concept agents form an automatic layer
that drives, avoids obstacles and executes missions without the language model; the language agent is the
deliberative voice. We propose to make the allocation of the voice explicit through a scalar \emph{demand} signal,
computed on every digest cycle without the language model,
\begin{equation}
    D_t \;=\; \sum_{k} \pi_k\, \varepsilon_{k,t}^{2} \;+\; \lambda\, \mathbb{1}\!\left[\text{person active}_t\right],
    \label{eq:inner_demand}
\end{equation}
where $\varepsilon_{k,t}$ are normalised prediction errors reported by the agents (unexplained residual occupancy,
flips of existence or door beliefs, localisation uncertainty, stalled or failed mission steps, controller refusals)
and $\pi_k$ their precisions. Three modes follow from two thresholds with asymmetric hysteresis:
\begin{equation}
    m_t =
    \begin{cases}
        \textsc{focused}    & D_t > \theta_{\text{hi}},\\[2pt]
        \textsc{wandering}  & D_\tau < \theta_{\text{lo}} \;\; \forall\, \tau \in [t - T, t],\\[2pt]
        \textsc{monitoring} & \text{otherwise.}
    \end{cases}
    \label{eq:inner_modes}
\end{equation}
Entering \textsc{wandering} requires demand to stay low for a dwell time $T$; leaving it requires a single sample
above $\theta_{\text{hi}}$. The robot is slow to drift off and instant to snap back. Safety never depends on this
arbitration: the controller's reflexes act on their own, and the deliberative voice arrives after the reflex, as
awareness follows braking in the driver.

\paragraph{Focused mode} runs the question loop of Section~\ref{sec:inner_speech:questions} on the triggering
event. \paragraph{Monitoring mode} makes no language-model calls; the voice is primed but silent.

\paragraph{Wandering mode.}
Freed from the task, the voice is seeded with one recent experience chosen by salience: an episode that ended
unresolved, a failure, an object that changed, or an upcoming goal. Each step's output, together with an
associative retrieval from episodic memory, seeds the next, so that the chain drifts as spontaneous thought does
\cite{christoff2016mind}. The voice has access only to memory queries and to the architecture's imagination tools,
which simulate observations and outcomes on a copy of the belief graph; it can neither act nor write beliefs. Its
products are of three kinds: counterfactuals over past episodes (``had the door been open, the crossing would have
taken 40\,s less''), hypotheses to be tested later, and prospective plans (``if the door is shut on the way back,
the route through the next room is shorter''). Functionally this is the default network's role in simulating
past and future \cite{buckner2008brain}, and the offline replay and model reduction that Active Inference assigns
to rest and sleep \cite{friston2017curiosity, hobson2012waking}. Replay is prioritised by the expected gain of the
memory it revisits \cite{mattar2018prioritized}.

\paragraph{Source monitoring.}
Humans occasionally misremember imagined events as perceived ones \cite{johnson1993source}. The robot must not.
Everything produced while wandering is stored with an \emph{imagined} provenance, separate from the factual
episodic record, and enters the learning loop only as a hypothesis to be tested, never as evidence.

\subsection{What the return to reality gives the automatic layer}
\label{sec:inner_speech:return}

The transition from \textsc{wandering} to \textsc{focused} cancels the chain in flight (its state is kept, so that
it may resume if demand subsides) and gives the full voice to the surprise. It serves the automatic layer at three
timescales.
\begin{enumerate}
    \item \textbf{Decision.} A surprise means the automatic layer has left its domain of competence, where its only
    responses are to retry, stop or fail. The voice adds re-interpretation of the scene, arbitration between goals,
    re-planning through remembered topology, and recourse to a person.
    \item \textbf{Attention.} The voice can raise precision where it matters: orienting sensors, increasing a
    detector's rate in a region, or passing a hypothesis to the concept agents as a spatial prior. Such a prior has
    bounded weight; it may accelerate the recognition of an object that is present but may not create one that is
    not, preserving the grounding constraint.
    \item \textbf{Teaching.} Each resolved surprise is a labelled example the automatic layer could not have
    produced alone: a residual later identified as a chair is a false negative of the chair detector; a door that
    is repeatedly shut in the evening becomes a prior; a recovery that repeatedly succeeds becomes a branch of the
    controller's behaviour tree. The domain of competence grows, and the situations that require the voice become
    rarer, which is the classical progression from effortful to autonomous skill \cite{fitts1967human}.
\end{enumerate}
Plans prepared while wandering act as a cache: when the anticipated situation occurs, the focused voice verifies a
plan instead of composing one. The return ends with an explicit release --- spatial priors cleared, safety margins
restored, the automatic layer confirmed back within its competence --- without which every surprise would leave
the robot permanently vigilant.

\subsection{Relation to chain-of-thought reasoning}
\label{sec:inner_speech:cot}

Chain-of-thought is inner speech in the functional sense, but of a specific kind, and the comparison sharpens the
design (Table~\ref{tab:inner_cot}). Plain chain-of-thought is convergent and ungrounded: tokens condition on
tokens until an answer is produced, and the trace is discarded. Interleaving reasoning with action and observation
grounds it \cite{yao2023react}; our focused loop is of this kind, with the additional property that the
observation, through \eqref{eq:inner_update}, rather than the next token, determines where the reasoning goes.
The wandering voice is closer to plain chain-of-thought in its isolation from perception, but divergent rather
than convergent, interruptible, and persistent in its products.

\begin{table}[ht]
    \centering
    \small
    \begin{tabular}{|>{\raggedright\arraybackslash}p{0.19\linewidth}|>{\raggedright\arraybackslash}p{0.22\linewidth}|>{\raggedright\arraybackslash}p{0.24\linewidth}|>{\raggedright\arraybackslash}p{0.24\linewidth}|}
        \hline
                                & \textbf{Chain-of-thought} & \textbf{Focused voice} & \textbf{Wandering voice} \\ \hline
        Goal                    & answer a given question   & resolve a given surprise & none given; seed by salience \\ \hline
        Dynamics                & convergent                & convergent             & divergent, associative \\ \hline
        Contact with the world  & none                      & thought--action--observation & none (imagination only) \\ \hline
        Persistence             & discarded                 & question objects, beliefs, lessons & imagined store, hypotheses \\ \hline
        Interruptible           & no                        & yes                    & yes, resumable \\ \hline
        Latency budget          & n/a                       & $\sim$1\,s, no extended thinking & unbounded, extended thinking \\ \hline
    \end{tabular}
    \caption{Chain-of-thought compared with the two modes of the proposed inner voice.}
    \label{tab:inner_cot}
\end{table}

Two practical consequences follow. First, \emph{think when there is time, not when there is pressure}. In the
deployed language agent extended thinking is disabled, because a thinking-mode model spent 3\,000--6\,000 tokens
(40--80\,s) before each tool call; that is incompatible with the focused mode but costs nothing in the wandering
mode, where it can be enabled with a large budget and pre-empted at any moment. Slow reasoning is done in advance
and fast reasoning uses its results. Second, chain-of-thought traces are known not always to reflect the
computation that produced the answer \cite{turpin2023language}, as human verbal reports do not always reflect the
causes of behaviour \cite{nisbett1977telling}. The faithful record of the robot's reasoning is therefore the
structured one --- hypotheses, predictions, observations and posterior updates --- not the text. Finally, wandering
chains whose hypotheses were later confirmed by perception form a corpus of reasoning verified against the world,
which could be used to fine-tune the local model in the manner of self-taught reasoning \cite{zelikman2022star},
with perception, rather than a known answer, as the filter.

\subsection{Consequences for belief and learning}
\label{sec:inner_speech:learning}

The loop changes how beliefs evolve in two ways. Within an episode, uncertainty is reduced where surprise
occurred rather than where the sensors happened to point, which is efficient but leaves the rest of the map less
examined and makes the robot's own data a biased sample; a decaying novelty term \cite{schwartenbeck2019curiosity}
and importance weighting of learned statistics are required to correct this. Across episodes, every closed
question is a labelled example, and learning proceeds at four timescales (Table~\ref{tab:inner_learning}).

\begin{table}[ht]
    \centering
    \small
    \begin{tabular}{|>{\raggedright\arraybackslash}p{0.14\linewidth}|>{\raggedright\arraybackslash}p{0.37\linewidth}|>{\raggedright\arraybackslash}p{0.37\linewidth}|}
        \hline
        \textbf{Timescale} & \textbf{What changes} & \textbf{Where, in this architecture} \\ \hline
        Seconds   & posterior over the hypotheses of one question         & the question object \eqref{eq:inner_question} \\ \hline
        Days      & priors: where objects are, which doors are shut, what moves & long-term object and topological memory, as Dirichlet counts \\ \hline
        Weeks     & likelihoods: detector and concept models              & false negatives and positives labelled by resolved questions \\ \hline
        Slow      & structure: new categories                             & recurrent, mutually similar ``unknown'' resolutions named by the voice and proposed as a new concept \\ \hline
    \end{tabular}
    \caption{Learning driven by the inner-question loop.}
    \label{tab:inner_learning}
\end{table}

The last row is where the hypothesis of a language for thought is most directly engaged: a cluster of anomalies
that resists every existing concept becomes learnable once it is named. The loop also learns about itself: the
salience weights $\pi_k$ in \eqref{eq:inner_demand} are adjusted by the information each trigger actually yielded,
asking a person is weighed against looking, and loops that repeatedly resolve the same way are compiled into
reflexive rules, so that the number of language-model calls per hour should \emph{decrease} as competence grows.

\subsection{Why a language model: generation versus evaluation}
\label{sec:inner_speech:why_lm}

None of the above requires a person, so the question arises of what a language model contributes to an
architecture that already performs inference and action selection. Our answer is a division of labour. Expected
free energy evaluates policies and Bayesian updating weighs hypotheses, but both operate over a space that has
already been enumerated: each concept agent carries a fixed generative model, policies must be listed before they
can be scored, and Bayesian model reduction can prune structure but not invent it \cite{friston2017curiosity}.
Whatever falls outside the given models remains an unexplained residual. A language model is weak where this
machinery is strong --- its confidences are not calibrated posteriors and it has no access to the sensors --- and
strong where it is weak: it is a cheap, broad generator of hypotheses, explanations and plans. We therefore use it
only as a \emph{proposal distribution}, in a generate-and-test loop whose test is active inference:
\begin{itemize}
    \item \emph{Hypotheses outside the built-in models.} A residual blob can be proposed to be a bag, a toppled
    chair or a robot vacuum, none of which any concept agent represents; the question loop then tests them.
    \item \emph{Amortised policy search.} Instead of enumerating multi-step policies, the model writes a few
    candidate missions --- small programs in the mission language --- which the gate validates and expected free
    energy scores.
    \item \emph{Relational binding across agents.} The agents' representations (signed distance fields, occupancy
    grids, existence log-odds, episodic records) are mutually incompatible; a sentence such as ``the chair that was
    not next to the table in the room behind door~3 yesterday'' binds them in one query no single agent could pose.
    \item \emph{Explanation from few examples.} After a single surprise the model can propose which variable
    mattered (``the door is shut because it is evening''), turning one episode into a testable generalisation.
    \item \emph{Abstraction and naming.} Episodes are compressed into reusable lessons, and recurrent unknowns are
    named, which is what makes them candidates for a new concept (Table~\ref{tab:inner_learning}).
\end{itemize}
Language, rather than an arbitrary learned embedding, matters for three reasons: it is discrete and compositional,
so new hypotheses and plans are assembled from known parts; it is inspectable, so that deterministic gates can parse
and check it and people can audit it; and, once the robot thinks in it, communication with a person comes at no
further cost.

\paragraph{Where the knowledge comes from.}
This raises a difficulty for the motivating thesis. If the inner voice precedes the outer one, the common-sense
knowledge on which the proposer's value rests cannot have arrived through communication --- yet a language model
acquires all of it from human text, that is, from other people's outer voices. The difficulty dissolves once the
\emph{faculty} is distinguished from its \emph{content}. The internalist claim concerns the former: a
computational system for building unbounded hierarchical structures, available for thought before it is used for
talk \cite{hauser2002faculty, berwick2016why}. It makes no claim that what a person knows was acquired without
language. Human knowledge has three sources: core systems for objects, agents, number and space, present before
learning \cite{spelke2007core}; individual experience; and, once outer language exists, cumulative cultural
transmission, which supplies most of what an adult knows \cite{tomasello1999cultural}. Even on the internalist
view, the inner voice of a modern human is filled largely by the outer voices of others.

The architecture has the same three sources, in a different arrangement. The concept agents' hand-built generative
models play the role of \emph{core knowledge}: given before experience and limited to a few domains. The robot's
own counts $n_x(h)$ in \eqref{eq:inner_prior} are \emph{individual experience}. The language model is
\emph{culture}: knowledge distilled from a community's communication and inherited, for this individual, at birth.
The robot thus resembles a child born into a culture more than an early hominin, and its inner voice is Vygotskian
in origin: its format and its content were both learned from outer speech. Two consequences follow. First, the
proposal does not demonstrate the internalist thesis and should not be read as doing so; its claim is functional
--- an inner use of language has epistemic value independently of communication --- and holds whatever the origin of
the language system. Second, the format and the content of the voice can in principle be separated, which turns the
question into an experiment (P7 in Section~\ref{sec:inner_speech:predictions}): a compositional generator over the
robot's own concepts and mission language, trained only on the robot's experience, in the manner of library
learning in program induction \cite{ellis2021dreamcoder}, would supply the format without the culture.

\subsection{Mapping onto the existing architecture}
\label{sec:inner_speech:mapping}

Several components already exist in the language agent. Robot events (mission reports, failures, completions)
already wake the model through a queue, in turns restricted to perceiving, querying memory, stopping and
proposing, never acting. Every action already passes a deterministic gate. The decision to ask or to assume is
already decision-theoretic. Unexplained occupancy is already isolated by a residual agent that acts as the null
hypothesis of the scene, and episodic and long-term object memories are already joined read-only. Missing are the
persistent question objects \eqref{eq:inner_question}, a prediction evaluator implementing
\eqref{eq:inner_update} outside the language model, the demand signal and mode arbitration
\eqref{eq:inner_demand}--\eqref{eq:inner_modes}, the imagined store with its provenance, and the channel through
which bounded spatial priors reach the concept agents.

\subsection{Testable predictions}
\label{sec:inner_speech:predictions}

Consistent with the evaluation principle adopted throughout --- policies compared against each other inside one
closed loop, with body, sensors and environment fixed --- the proposal makes seven falsifiable predictions, each
tested in simulation by moving known objects and introducing unfamiliar ones:
\begin{enumerate}[label=\textbf{P\arabic*}]
    \item The unexplained residual mass per room decreases faster with the focused loop than without it.
    \item The proposer's calibration on resolved questions improves, and $\beta_x$ in \eqref{eq:inner_prior}
    decreases, as the robot's own counts accumulate.
    \item Detector false-negative rates fall after retraining on question-labelled examples.
    \item Inner questions per hour decrease over repeated exposure to the same environment while answer accuracy
    is maintained.
    \item Plans prepared while wandering reduce the latency from surprise to a validated response, and the time to
    notice a deliberately moved object does not grow despite habituation.
    \item With the language-model proposer replaced by a fixed hypothesis set (the existing concept classes and
    ``unknown'') and template-based missions, the loop resolves fewer questions about objects and failures outside
    the built-in models, and acquires new categories more slowly.
    \item A compositional generator trained only on the robot's own experience (format without culture) recovers
    part of the advantage in P6 within a familiar environment, and less of it the further a situation departs from
    what the robot has seen.
\end{enumerate}
A failure of P1 would indicate that the epistemic value of language-mediated questions does not exceed that of
the existing exploration; a failure of P4 would indicate that the loop does not transfer competence to the
automatic layer, which is its principal claim; and a failure of P6 would indicate that the language model adds
nothing a fixed hypothesis set does not, which would remove the case for including it. P7 separates the
contribution of language as a format for thought from that of the knowledge it carries.

\subsection{Limitations}
\label{sec:inner_speech:limitations}

The demand signal \eqref{eq:inner_demand} sums heterogeneous prediction errors whose precisions must themselves be
calibrated, and a poorly scaled term will either suppress the voice or keep it permanently engaged. A single
language model serves the person, the focused loop and the wandering loop; pre-emption keeps the person first but
bounds how much offline reasoning a busy robot can do. The tempered prior \eqref{eq:inner_prior} is a heuristic
combination rather than a derived posterior. The generator of P7 is substantially harder to build than the
ablation of P6, and may have to be approximated by a small model trained on the robot's logs and mission corpus,
whose pre-training would leak some culture back in. Finally, the analogy with human inner speech is functional, not
mechanistic: the proposal is that the same computational role --- serial, compositional self-questioning in the
service of epistemic value --- is useful for a robot, not that the robot's voice resembles the human one in any
further respect.

% and add inner_speech.bib as a second resource:  \addbibresource{inner_speech.bib}
% All keys are new (none collides with the keys main.tex already cites).
% Packages used: amsmath, amssymb, enumitem (already loaded by main.tex) and array, which main.tex does not
% load yet: add \usepackage{array} to its preamble.
% =====================================================================================================

\section{Inner Speech as Epistemic Action}
\label{sec:inner_speech}

The robot described so far uses language in one direction only: a person speaks, and a language agent turns
the request into commands for the controller, grounded in a textual digest of the belief graph. This section
proposes to turn that capability inward. We argue that a language model embedded in an Active Inference
architecture can play the role that inner speech plays in human cognition: a serial, compositional medium in
which the agent poses questions to itself, generates hypotheses, rehearses plans and replays its past, and that
this role follows naturally from the epistemic component of expected free energy. The proposal is stated
together with the constraints that keep such a voice from becoming a source of confabulated belief, and with the
experiments that would falsify it.

\subsection{Motivation: language for thought before language for talk}
\label{sec:inner_speech:motivation}

Two traditions disagree on the origin of inner speech but agree on its function. On the internalist view,
human language evolved first as an internal system for structuring thought, and its use for communication is a
later externalisation \cite{berwick2016why}. On the developmental view, inner speech is social speech that the
child gradually internalises, abbreviates and turns into a tool for self-regulation \cite{vygotsky1986thought}.
Contemporary reviews treat inner speech as functionally involved in planning, problem solving, self-monitoring and
the maintenance of goals in working memory \cite{aldersonday2015inner}, and language more generally as a
cognitive technology that lets a learner manipulate discrete, recombinable tokens rather than raw percepts
\cite{clark1998magic}.

Large language models give a striking instance of the developmental route. They are trained exclusively on
communicative text, yet producing intermediate reasoning in that same medium before answering --- chain-of-thought
--- measurably improves their performance \cite{wei2022chain}. A system built for talk becomes better at thinking
by talking to itself. Our language agent sits at the same point of its own development: it was built to converse
with a person, and the question we ask is what it contributes once its speech is addressed to the robot itself.
Section~\ref{sec:inner_speech:why_lm} returns to what this origin implies for the internalist view.

\subsection{Inner questions as epistemic actions}
\label{sec:inner_speech:questions}

In Active Inference, policies are scored by expected free energy, whose epistemic term rewards the expected
reduction of uncertainty about hidden states \cite{friston2017process, parr2022active}. A question is the
linguistic form of exactly this quantity: it names a hidden variable whose value is uncertain and whose resolution
is expected to matter. We therefore model an inner question as a persistent object
\begin{equation}
    \mathcal{Q} = \big\langle\, e,\; \mathcal{H},\; q(h),\; \mathcal{P},\; b,\; s \,\big\rangle ,
    \label{eq:inner_question}
\end{equation}
where $e$ is the triggering event (for example an obstacle no concept agent explains), $\mathcal{H}$ a finite set
of named hypotheses (``a chair'', ``a box'', ``unknown''), $q(h)$ the current posterior over them, $\mathcal{P}$
the set of pending predictions, $b$ a budget in time and motion, and $s \in \{\text{open},\text{resolved},
\text{abandoned}\}$ its status. Questions form a stack in the sense of a \emph{question under discussion}
\cite{roberts2012information}: ``is it a chair?'' is a sub-question whose answer partially resolves the parent
``what is this?''. The question, not the language model's context window, carries the state of the reasoning from
one step to the next.

\paragraph{One step of the loop.}
Given $\mathcal{Q}$, the language model proposes a test: an epistemic action $a$ (turn towards the anomaly,
approach it, look with the front camera) together with the outcome each hypothesis predicts. The value of the test
is its expected information gain minus its cost,
\begin{equation}
    V(a) \;=\; \underbrace{H\!\left[q(h)\right] \;-\; \mathbb{E}_{p(o \mid a)}\!\left[ H\!\left[q(h \mid o, a)\right]
    \right]}_{I(h;\,o \mid a)} \;-\; c(a),
    \label{eq:inner_value}
\end{equation}
which is the negative of the epistemic part of expected free energy with an explicit action cost. The test is
executed only if $V(a) > 0$ and the action passes the same deterministic gate that already validates every
command the language agent issues. When the observation $o$ arrives, the posterior is updated by the concept
agents' own likelihoods,
\begin{equation}
    q(h \mid o, a) \;\propto\; p(o \mid h, a)\, q(h),
    \label{eq:inner_update}
\end{equation}
and the loop continues, spawns a sub-question, or terminates when
\begin{equation}
    H\!\left[q(h)\right] < \eta \quad\lor\quad \max_a V(a) \le 0 \quad\lor\quad b \text{ exhausted}.
    \label{eq:inner_stop}
\end{equation}
The stopping rule is what separates reasoning from rumination: a loop with no expected information to gain closes,
possibly with the answer ``unknown'', and its trigger is habituated for a refractory period.

\paragraph{The grounding constraint.}
The central design rule is that \emph{the language model proposes and perception disposes}. Text produced by the
model may introduce hypotheses, predictions and tests; it may never move a posterior. Only the likelihood
$p(o \mid h, a)$ in \eqref{eq:inner_update}, computed by the perceptual agents from sensor data, does. This is the
belief-side counterpart of the rule the language agent already enforces for actions, where every proposed command
is validated in C++ against the belief graph before it reaches the controller. Without it, a self-addressed
language loop is a closed system of text conditioned on text and converges on its own confabulations.

\paragraph{The proposer as a calibrated prior.}
The hypotheses the model names, and the probabilities it attaches to them, encode world knowledge the robot has
not acquired by itself (``chairs come with tables''). We treat this as a prior whose weight is earned. Let
$\pi_{\text{LM}}(h \mid x)$ be the proposer's distribution in context $x$ and $n_x(h)$ the robot's own counts of
resolved questions in that context. The prior entering \eqref{eq:inner_update} is
\begin{equation}
    q_0(h \mid x) \;\propto\; \pi_{\text{LM}}(h \mid x)^{\beta_x}\, \big(\alpha + n_x(h)\big),
    \label{eq:inner_prior}
\end{equation}
a tempered product of the proposer and a Dirichlet over experience, where $\beta_x \in [0,1]$ is set from the
proposer's measured calibration in that context (for instance from its Brier score on past questions). Early in
the robot's life the proposer dominates; as counts accumulate, the robot's own experience replaces borrowed
knowledge. The same conjugate form already governs the language agent's decision to ask a person or to assume
that they moved an object, where a Beta posterior over the person's past answers is compared against the costs of
asking and of being wrong.

\subsection{Attention, automaticity and the wandering mind}
\label{sec:inner_speech:modes}

A human driver on an easy road steers without attending to it; the mind is free to plan, imagine or recall, and is
recaptured the instant something unexpected happens \cite{smallwood2015science, christoff2016mind}. In Active
Inference this is a statement about precision: while prediction errors are small and the habitual policy is
reliable, attention --- the precision assigned to sensory prediction errors --- can be withdrawn from the task; a
surprise raises it again \cite{feldman2010attention}. The arbitration between habitual and deliberative control by
their relative uncertainty \cite{daw2005uncertainty} has the same structure.

The architecture already separates the two systems. The controller and the concept agents form an automatic layer
that drives, avoids obstacles and executes missions without the language model; the language agent is the
deliberative voice. We propose to make the allocation of the voice explicit through a scalar \emph{demand} signal,
computed on every digest cycle without the language model,
\begin{equation}
    D_t \;=\; \sum_{k} \pi_k\, \varepsilon_{k,t}^{2} \;+\; \lambda\, \mathbb{1}\!\left[\text{person active}_t\right],
    \label{eq:inner_demand}
\end{equation}
where $\varepsilon_{k,t}$ are normalised prediction errors reported by the agents (unexplained residual occupancy,
flips of existence or door beliefs, localisation uncertainty, stalled or failed mission steps, controller refusals)
and $\pi_k$ their precisions. Three modes follow from two thresholds with asymmetric hysteresis:
\begin{equation}
    m_t =
    \begin{cases}
        \textsc{focused}    & D_t > \theta_{\text{hi}},\\[2pt]
        \textsc{wandering}  & D_\tau < \theta_{\text{lo}} \;\; \forall\, \tau \in [t - T, t],\\[2pt]
        \textsc{monitoring} & \text{otherwise.}
    \end{cases}
    \label{eq:inner_modes}
\end{equation}
Entering \textsc{wandering} requires demand to stay low for a dwell time $T$; leaving it requires a single sample
above $\theta_{\text{hi}}$. The robot is slow to drift off and instant to snap back. Safety never depends on this
arbitration: the controller's reflexes act on their own, and the deliberative voice arrives after the reflex, as
awareness follows braking in the driver.

\paragraph{Focused mode} runs the question loop of Section~\ref{sec:inner_speech:questions} on the triggering
event. \paragraph{Monitoring mode} makes no language-model calls; the voice is primed but silent.

\paragraph{Wandering mode.}
Freed from the task, the voice is seeded with one recent experience chosen by salience: an episode that ended
unresolved, a failure, an object that changed, or an upcoming goal. Each step's output, together with an
associative retrieval from episodic memory, seeds the next, so that the chain drifts as spontaneous thought does
\cite{christoff2016mind}. The voice has access only to memory queries and to the architecture's imagination tools,
which simulate observations and outcomes on a copy of the belief graph; it can neither act nor write beliefs. Its
products are of three kinds: counterfactuals over past episodes (``had the door been open, the crossing would have
taken 40\,s less''), hypotheses to be tested later, and prospective plans (``if the door is shut on the way back,
the route through the next room is shorter''). Functionally this is the default network's role in simulating
past and future \cite{buckner2008brain}, and the offline replay and model reduction that Active Inference assigns
to rest and sleep \cite{friston2017curiosity, hobson2012waking}. Replay is prioritised by the expected gain of the
memory it revisits \cite{mattar2018prioritized}.

\paragraph{Source monitoring.}
Humans occasionally misremember imagined events as perceived ones \cite{johnson1993source}. The robot must not.
Everything produced while wandering is stored with an \emph{imagined} provenance, separate from the factual
episodic record, and enters the learning loop only as a hypothesis to be tested, never as evidence.

\subsection{What the return to reality gives the automatic layer}
\label{sec:inner_speech:return}

The transition from \textsc{wandering} to \textsc{focused} cancels the chain in flight (its state is kept, so that
it may resume if demand subsides) and gives the full voice to the surprise. It serves the automatic layer at three
timescales.
\begin{enumerate}
    \item \textbf{Decision.} A surprise means the automatic layer has left its domain of competence, where its only
    responses are to retry, stop or fail. The voice adds re-interpretation of the scene, arbitration between goals,
    re-planning through remembered topology, and recourse to a person.
    \item \textbf{Attention.} The voice can raise precision where it matters: orienting sensors, increasing a
    detector's rate in a region, or passing a hypothesis to the concept agents as a spatial prior. Such a prior has
    bounded weight; it may accelerate the recognition of an object that is present but may not create one that is
    not, preserving the grounding constraint.
    \item \textbf{Teaching.} Each resolved surprise is a labelled example the automatic layer could not have
    produced alone: a residual later identified as a chair is a false negative of the chair detector; a door that
    is repeatedly shut in the evening becomes a prior; a recovery that repeatedly succeeds becomes a branch of the
    controller's behaviour tree. The domain of competence grows, and the situations that require the voice become
    rarer, which is the classical progression from effortful to autonomous skill \cite{fitts1967human}.
\end{enumerate}
Plans prepared while wandering act as a cache: when the anticipated situation occurs, the focused voice verifies a
plan instead of composing one. The return ends with an explicit release --- spatial priors cleared, safety margins
restored, the automatic layer confirmed back within its competence --- without which every surprise would leave
the robot permanently vigilant.

\subsection{Relation to chain-of-thought reasoning}
\label{sec:inner_speech:cot}

Chain-of-thought is inner speech in the functional sense, but of a specific kind, and the comparison sharpens the
design (Table~\ref{tab:inner_cot}). Plain chain-of-thought is convergent and ungrounded: tokens condition on
tokens until an answer is produced, and the trace is discarded. Interleaving reasoning with action and observation
grounds it \cite{yao2023react}; our focused loop is of this kind, with the additional property that the
observation, through \eqref{eq:inner_update}, rather than the next token, determines where the reasoning goes.
The wandering voice is closer to plain chain-of-thought in its isolation from perception, but divergent rather
than convergent, interruptible, and persistent in its products.

\begin{table}[ht]
    \centering
    \small
    \begin{tabular}{|>{\raggedright\arraybackslash}p{0.19\linewidth}|>{\raggedright\arraybackslash}p{0.22\linewidth}|>{\raggedright\arraybackslash}p{0.24\linewidth}|>{\raggedright\arraybackslash}p{0.24\linewidth}|}
        \hline
                                & \textbf{Chain-of-thought} & \textbf{Focused voice} & \textbf{Wandering voice} \\ \hline
        Goal                    & answer a given question   & resolve a given surprise & none given; seed by salience \\ \hline
        Dynamics                & convergent                & convergent             & divergent, associative \\ \hline
        Contact with the world  & none                      & thought--action--observation & none (imagination only) \\ \hline
        Persistence             & discarded                 & question objects, beliefs, lessons & imagined store, hypotheses \\ \hline
        Interruptible           & no                        & yes                    & yes, resumable \\ \hline
        Latency budget          & n/a                       & $\sim$1\,s, no extended thinking & unbounded, extended thinking \\ \hline
    \end{tabular}
    \caption{Chain-of-thought compared with the two modes of the proposed inner voice.}
    \label{tab:inner_cot}
\end{table}

Two practical consequences follow. First, \emph{think when there is time, not when there is pressure}. In the
deployed language agent extended thinking is disabled, because a thinking-mode model spent 3\,000--6\,000 tokens
(40--80\,s) before each tool call; that is incompatible with the focused mode but costs nothing in the wandering
mode, where it can be enabled with a large budget and pre-empted at any moment. Slow reasoning is done in advance
and fast reasoning uses its results. Second, chain-of-thought traces are known not always to reflect the
computation that produced the answer \cite{turpin2023language}, as human verbal reports do not always reflect the
causes of behaviour \cite{nisbett1977telling}. The faithful record of the robot's reasoning is therefore the
structured one --- hypotheses, predictions, observations and posterior updates --- not the text. Finally, wandering
chains whose hypotheses were later confirmed by perception form a corpus of reasoning verified against the world,
which could be used to fine-tune the local model in the manner of self-taught reasoning \cite{zelikman2022star},
with perception, rather than a known answer, as the filter.

\subsection{Consequences for belief and learning}
\label{sec:inner_speech:learning}

The loop changes how beliefs evolve in two ways. Within an episode, uncertainty is reduced where surprise
occurred rather than where the sensors happened to point, which is efficient but leaves the rest of the map less
examined and makes the robot's own data a biased sample; a decaying novelty term \cite{schwartenbeck2019curiosity}
and importance weighting of learned statistics are required to correct this. Across episodes, every closed
question is a labelled example, and learning proceeds at four timescales (Table~\ref{tab:inner_learning}).

\begin{table}[ht]
    \centering
    \small
    \begin{tabular}{|>{\raggedright\arraybackslash}p{0.14\linewidth}|>{\raggedright\arraybackslash}p{0.37\linewidth}|>{\raggedright\arraybackslash}p{0.37\linewidth}|}
        \hline
        \textbf{Timescale} & \textbf{What changes} & \textbf{Where, in this architecture} \\ \hline
        Seconds   & posterior over the hypotheses of one question         & the question object \eqref{eq:inner_question} \\ \hline
        Days      & priors: where objects are, which doors are shut, what moves & long-term object and topological memory, as Dirichlet counts \\ \hline
        Weeks     & likelihoods: detector and concept models              & false negatives and positives labelled by resolved questions \\ \hline
        Slow      & structure: new categories                             & recurrent, mutually similar ``unknown'' resolutions named by the voice and proposed as a new concept \\ \hline
    \end{tabular}
    \caption{Learning driven by the inner-question loop.}
    \label{tab:inner_learning}
\end{table}

The last row is where the hypothesis of a language for thought is most directly engaged: a cluster of anomalies
that resists every existing concept becomes learnable once it is named. The loop also learns about itself: the
salience weights $\pi_k$ in \eqref{eq:inner_demand} are adjusted by the information each trigger actually yielded,
asking a person is weighed against looking, and loops that repeatedly resolve the same way are compiled into
reflexive rules, so that the number of language-model calls per hour should \emph{decrease} as competence grows.

\subsection{Why a language model: generation versus evaluation}
\label{sec:inner_speech:why_lm}

None of the above requires a person, so the question arises of what a language model contributes to an
architecture that already performs inference and action selection. Our answer is a division of labour. Expected
free energy evaluates policies and Bayesian updating weighs hypotheses, but both operate over a space that has
already been enumerated: each concept agent carries a fixed generative model, policies must be listed before they
can be scored, and Bayesian model reduction can prune structure but not invent it \cite{friston2017curiosity}.
Whatever falls outside the given models remains an unexplained residual. A language model is weak where this
machinery is strong --- its confidences are not calibrated posteriors and it has no access to the sensors --- and
strong where it is weak: it is a cheap, broad generator of hypotheses, explanations and plans. We therefore use it
only as a \emph{proposal distribution}, in a generate-and-test loop whose test is active inference:
\begin{itemize}
    \item \emph{Hypotheses outside the built-in models.} A residual blob can be proposed to be a bag, a toppled
    chair or a robot vacuum, none of which any concept agent represents; the question loop then tests them.
    \item \emph{Amortised policy search.} Instead of enumerating multi-step policies, the model writes a few
    candidate missions --- small programs in the mission language --- which the gate validates and expected free
    energy scores.
    \item \emph{Relational binding across agents.} The agents' representations (signed distance fields, occupancy
    grids, existence log-odds, episodic records) are mutually incompatible; a sentence such as ``the chair that was
    not next to the table in the room behind door~3 yesterday'' binds them in one query no single agent could pose.
    \item \emph{Explanation from few examples.} After a single surprise the model can propose which variable
    mattered (``the door is shut because it is evening''), turning one episode into a testable generalisation.
    \item \emph{Abstraction and naming.} Episodes are compressed into reusable lessons, and recurrent unknowns are
    named, which is what makes them candidates for a new concept (Table~\ref{tab:inner_learning}).
\end{itemize}
Language, rather than an arbitrary learned embedding, matters for three reasons: it is discrete and compositional,
so new hypotheses and plans are assembled from known parts; it is inspectable, so that deterministic gates can parse
and check it and people can audit it; and, once the robot thinks in it, communication with a person comes at no
further cost.

\paragraph{Where the knowledge comes from.}
This raises a difficulty for the motivating thesis. If the inner voice precedes the outer one, the common-sense
knowledge on which the proposer's value rests cannot have arrived through communication --- yet a language model
acquires all of it from human text, that is, from other people's outer voices. The difficulty dissolves once the
\emph{faculty} is distinguished from its \emph{content}. The internalist claim concerns the former: a
computational system for building unbounded hierarchical structures, available for thought before it is used for
talk \cite{hauser2002faculty, berwick2016why}. It makes no claim that what a person knows was acquired without
language. Human knowledge has three sources: core systems for objects, agents, number and space, present before
learning \cite{spelke2007core}; individual experience; and, once outer language exists, cumulative cultural
transmission, which supplies most of what an adult knows \cite{tomasello1999cultural}. Even on the internalist
view, the inner voice of a modern human is filled largely by the outer voices of others.

The architecture has the same three sources, in a different arrangement. The concept agents' hand-built generative
models play the role of \emph{core knowledge}: given before experience and limited to a few domains. The robot's
own counts $n_x(h)$ in \eqref{eq:inner_prior} are \emph{individual experience}. The language model is
\emph{culture}: knowledge distilled from a community's communication and inherited, for this individual, at birth.
The robot thus resembles a child born into a culture more than an early hominin, and its inner voice is Vygotskian
in origin: its format and its content were both learned from outer speech. Two consequences follow. First, the
proposal does not demonstrate the internalist thesis and should not be read as doing so; its claim is functional
--- an inner use of language has epistemic value independently of communication --- and holds whatever the origin of
the language system. Second, the format and the content of the voice can in principle be separated, which turns the
question into an experiment (P7 in Section~\ref{sec:inner_speech:predictions}): a compositional generator over the
robot's own concepts and mission language, trained only on the robot's experience, in the manner of library
learning in program induction \cite{ellis2021dreamcoder}, would supply the format without the culture.

\subsection{Mapping onto the existing architecture}
\label{sec:inner_speech:mapping}

Several components already exist in the language agent. Robot events (mission reports, failures, completions)
already wake the model through a queue, in turns restricted to perceiving, querying memory, stopping and
proposing, never acting. Every action already passes a deterministic gate. The decision to ask or to assume is
already decision-theoretic. Unexplained occupancy is already isolated by a residual agent that acts as the null
hypothesis of the scene, and episodic and long-term object memories are already joined read-only. Missing are the
persistent question objects \eqref{eq:inner_question}, a prediction evaluator implementing
\eqref{eq:inner_update} outside the language model, the demand signal and mode arbitration
\eqref{eq:inner_demand}--\eqref{eq:inner_modes}, the imagined store with its provenance, and the channel through
which bounded spatial priors reach the concept agents.

\subsection{Testable predictions}
\label{sec:inner_speech:predictions}

Consistent with the evaluation principle adopted throughout --- policies compared against each other inside one
closed loop, with body, sensors and environment fixed --- the proposal makes seven falsifiable predictions, each
tested in simulation by moving known objects and introducing unfamiliar ones:
\begin{enumerate}[label=\textbf{P\arabic*}]
    \item The unexplained residual mass per room decreases faster with the focused loop than without it.
    \item The proposer's calibration on resolved questions improves, and $\beta_x$ in \eqref{eq:inner_prior}
    decreases, as the robot's own counts accumulate.
    \item Detector false-negative rates fall after retraining on question-labelled examples.
    \item Inner questions per hour decrease over repeated exposure to the same environment while answer accuracy
    is maintained.
    \item Plans prepared while wandering reduce the latency from surprise to a validated response, and the time to
    notice a deliberately moved object does not grow despite habituation.
    \item With the language-model proposer replaced by a fixed hypothesis set (the existing concept classes and
    ``unknown'') and template-based missions, the loop resolves fewer questions about objects and failures outside
    the built-in models, and acquires new categories more slowly.
    \item A compositional generator trained only on the robot's own experience (format without culture) recovers
    part of the advantage in P6 within a familiar environment, and less of it the further a situation departs from
    what the robot has seen.
\end{enumerate}
A failure of P1 would indicate that the epistemic value of language-mediated questions does not exceed that of
the existing exploration; a failure of P4 would indicate that the loop does not transfer competence to the
automatic layer, which is its principal claim; and a failure of P6 would indicate that the language model adds
nothing a fixed hypothesis set does not, which would remove the case for including it. P7 separates the
contribution of language as a format for thought from that of the knowledge it carries.

\subsection{Limitations}
\label{sec:inner_speech:limitations}

The demand signal \eqref{eq:inner_demand} sums heterogeneous prediction errors whose precisions must themselves be
calibrated, and a poorly scaled term will either suppress the voice or keep it permanently engaged. A single
language model serves the person, the focused loop and the wandering loop; pre-emption keeps the person first but
bounds how much offline reasoning a busy robot can do. The tempered prior \eqref{eq:inner_prior} is a heuristic
combination rather than a derived posterior. The generator of P7 is substantially harder to build than the
ablation of P6, and may have to be approximated by a small model trained on the robot's logs and mission corpus,
whose pre-training would leak some culture back in. Finally, the analogy with human inner speech is functional, not
mechanistic: the proposal is that the same computational role --- serial, compositional self-questioning in the
service of epistemic value --- is useful for a robot, not that the robot's voice resembles the human one in any
further respect.

%     % =====================================================================================================
% language_hierarchy_section.tex — Language over the belief hierarchy: culture, sequencing and imagination
%
% DRAFT (2026-11-19). A proposal section: the architecture it describes is NOT yet implemented beyond the parts
% stated as existing. Source design note: docs/LANGUAGE_OVER_BELIEF_HIERARCHY.md.
% Include AFTER inner_speech_section.tex (it refers to that section's equations and modes):
%     % =====================================================================================================
% inner_speech_section.tex — Inner speech, mind-wandering and active inference
%
% DRAFT (2026-11-18). A proposal section: it describes an architecture that is NOT yet implemented,
% grounded in the existing language_agent (robocomp/active_inference). Include from main.tex with
%     % =====================================================================================================
% inner_speech_section.tex — Inner speech, mind-wandering and active inference
%
% DRAFT (2026-11-18). A proposal section: it describes an architecture that is NOT yet implemented,
% grounded in the existing language_agent (robocomp/active_inference). Include from main.tex with
%     \input{inner_speech_section}
% and add inner_speech.bib as a second resource:  \addbibresource{inner_speech.bib}
% All keys are new (none collides with the keys main.tex already cites).
% Packages used: amsmath, amssymb, enumitem (already loaded by main.tex) and array, which main.tex does not
% load yet: add \usepackage{array} to its preamble.
% =====================================================================================================

\section{Inner Speech as Epistemic Action}
\label{sec:inner_speech}

The robot described so far uses language in one direction only: a person speaks, and a language agent turns
the request into commands for the controller, grounded in a textual digest of the belief graph. This section
proposes to turn that capability inward. We argue that a language model embedded in an Active Inference
architecture can play the role that inner speech plays in human cognition: a serial, compositional medium in
which the agent poses questions to itself, generates hypotheses, rehearses plans and replays its past, and that
this role follows naturally from the epistemic component of expected free energy. The proposal is stated
together with the constraints that keep such a voice from becoming a source of confabulated belief, and with the
experiments that would falsify it.

\subsection{Motivation: language for thought before language for talk}
\label{sec:inner_speech:motivation}

Two traditions disagree on the origin of inner speech but agree on its function. On the internalist view,
human language evolved first as an internal system for structuring thought, and its use for communication is a
later externalisation \cite{berwick2016why}. On the developmental view, inner speech is social speech that the
child gradually internalises, abbreviates and turns into a tool for self-regulation \cite{vygotsky1986thought}.
Contemporary reviews treat inner speech as functionally involved in planning, problem solving, self-monitoring and
the maintenance of goals in working memory \cite{aldersonday2015inner}, and language more generally as a
cognitive technology that lets a learner manipulate discrete, recombinable tokens rather than raw percepts
\cite{clark1998magic}.

Large language models give a striking instance of the developmental route. They are trained exclusively on
communicative text, yet producing intermediate reasoning in that same medium before answering --- chain-of-thought
--- measurably improves their performance \cite{wei2022chain}. A system built for talk becomes better at thinking
by talking to itself. Our language agent sits at the same point of its own development: it was built to converse
with a person, and the question we ask is what it contributes once its speech is addressed to the robot itself.
Section~\ref{sec:inner_speech:why_lm} returns to what this origin implies for the internalist view.

\subsection{Inner questions as epistemic actions}
\label{sec:inner_speech:questions}

In Active Inference, policies are scored by expected free energy, whose epistemic term rewards the expected
reduction of uncertainty about hidden states \cite{friston2017process, parr2022active}. A question is the
linguistic form of exactly this quantity: it names a hidden variable whose value is uncertain and whose resolution
is expected to matter. We therefore model an inner question as a persistent object
\begin{equation}
    \mathcal{Q} = \big\langle\, e,\; \mathcal{H},\; q(h),\; \mathcal{P},\; b,\; s \,\big\rangle ,
    \label{eq:inner_question}
\end{equation}
where $e$ is the triggering event (for example an obstacle no concept agent explains), $\mathcal{H}$ a finite set
of named hypotheses (``a chair'', ``a box'', ``unknown''), $q(h)$ the current posterior over them, $\mathcal{P}$
the set of pending predictions, $b$ a budget in time and motion, and $s \in \{\text{open},\text{resolved},
\text{abandoned}\}$ its status. Questions form a stack in the sense of a \emph{question under discussion}
\cite{roberts2012information}: ``is it a chair?'' is a sub-question whose answer partially resolves the parent
``what is this?''. The question, not the language model's context window, carries the state of the reasoning from
one step to the next.

\paragraph{One step of the loop.}
Given $\mathcal{Q}$, the language model proposes a test: an epistemic action $a$ (turn towards the anomaly,
approach it, look with the front camera) together with the outcome each hypothesis predicts. The value of the test
is its expected information gain minus its cost,
\begin{equation}
    V(a) \;=\; \underbrace{H\!\left[q(h)\right] \;-\; \mathbb{E}_{p(o \mid a)}\!\left[ H\!\left[q(h \mid o, a)\right]
    \right]}_{I(h;\,o \mid a)} \;-\; c(a),
    \label{eq:inner_value}
\end{equation}
which is the negative of the epistemic part of expected free energy with an explicit action cost. The test is
executed only if $V(a) > 0$ and the action passes the same deterministic gate that already validates every
command the language agent issues. When the observation $o$ arrives, the posterior is updated by the concept
agents' own likelihoods,
\begin{equation}
    q(h \mid o, a) \;\propto\; p(o \mid h, a)\, q(h),
    \label{eq:inner_update}
\end{equation}
and the loop continues, spawns a sub-question, or terminates when
\begin{equation}
    H\!\left[q(h)\right] < \eta \quad\lor\quad \max_a V(a) \le 0 \quad\lor\quad b \text{ exhausted}.
    \label{eq:inner_stop}
\end{equation}
The stopping rule is what separates reasoning from rumination: a loop with no expected information to gain closes,
possibly with the answer ``unknown'', and its trigger is habituated for a refractory period.

\paragraph{The grounding constraint.}
The central design rule is that \emph{the language model proposes and perception disposes}. Text produced by the
model may introduce hypotheses, predictions and tests; it may never move a posterior. Only the likelihood
$p(o \mid h, a)$ in \eqref{eq:inner_update}, computed by the perceptual agents from sensor data, does. This is the
belief-side counterpart of the rule the language agent already enforces for actions, where every proposed command
is validated in C++ against the belief graph before it reaches the controller. Without it, a self-addressed
language loop is a closed system of text conditioned on text and converges on its own confabulations.

\paragraph{The proposer as a calibrated prior.}
The hypotheses the model names, and the probabilities it attaches to them, encode world knowledge the robot has
not acquired by itself (``chairs come with tables''). We treat this as a prior whose weight is earned. Let
$\pi_{\text{LM}}(h \mid x)$ be the proposer's distribution in context $x$ and $n_x(h)$ the robot's own counts of
resolved questions in that context. The prior entering \eqref{eq:inner_update} is
\begin{equation}
    q_0(h \mid x) \;\propto\; \pi_{\text{LM}}(h \mid x)^{\beta_x}\, \big(\alpha + n_x(h)\big),
    \label{eq:inner_prior}
\end{equation}
a tempered product of the proposer and a Dirichlet over experience, where $\beta_x \in [0,1]$ is set from the
proposer's measured calibration in that context (for instance from its Brier score on past questions). Early in
the robot's life the proposer dominates; as counts accumulate, the robot's own experience replaces borrowed
knowledge. The same conjugate form already governs the language agent's decision to ask a person or to assume
that they moved an object, where a Beta posterior over the person's past answers is compared against the costs of
asking and of being wrong.

\subsection{Attention, automaticity and the wandering mind}
\label{sec:inner_speech:modes}

A human driver on an easy road steers without attending to it; the mind is free to plan, imagine or recall, and is
recaptured the instant something unexpected happens \cite{smallwood2015science, christoff2016mind}. In Active
Inference this is a statement about precision: while prediction errors are small and the habitual policy is
reliable, attention --- the precision assigned to sensory prediction errors --- can be withdrawn from the task; a
surprise raises it again \cite{feldman2010attention}. The arbitration between habitual and deliberative control by
their relative uncertainty \cite{daw2005uncertainty} has the same structure.

The architecture already separates the two systems. The controller and the concept agents form an automatic layer
that drives, avoids obstacles and executes missions without the language model; the language agent is the
deliberative voice. We propose to make the allocation of the voice explicit through a scalar \emph{demand} signal,
computed on every digest cycle without the language model,
\begin{equation}
    D_t \;=\; \sum_{k} \pi_k\, \varepsilon_{k,t}^{2} \;+\; \lambda\, \mathbb{1}\!\left[\text{person active}_t\right],
    \label{eq:inner_demand}
\end{equation}
where $\varepsilon_{k,t}$ are normalised prediction errors reported by the agents (unexplained residual occupancy,
flips of existence or door beliefs, localisation uncertainty, stalled or failed mission steps, controller refusals)
and $\pi_k$ their precisions. Three modes follow from two thresholds with asymmetric hysteresis:
\begin{equation}
    m_t =
    \begin{cases}
        \textsc{focused}    & D_t > \theta_{\text{hi}},\\[2pt]
        \textsc{wandering}  & D_\tau < \theta_{\text{lo}} \;\; \forall\, \tau \in [t - T, t],\\[2pt]
        \textsc{monitoring} & \text{otherwise.}
    \end{cases}
    \label{eq:inner_modes}
\end{equation}
Entering \textsc{wandering} requires demand to stay low for a dwell time $T$; leaving it requires a single sample
above $\theta_{\text{hi}}$. The robot is slow to drift off and instant to snap back. Safety never depends on this
arbitration: the controller's reflexes act on their own, and the deliberative voice arrives after the reflex, as
awareness follows braking in the driver.

\paragraph{Focused mode} runs the question loop of Section~\ref{sec:inner_speech:questions} on the triggering
event. \paragraph{Monitoring mode} makes no language-model calls; the voice is primed but silent.

\paragraph{Wandering mode.}
Freed from the task, the voice is seeded with one recent experience chosen by salience: an episode that ended
unresolved, a failure, an object that changed, or an upcoming goal. Each step's output, together with an
associative retrieval from episodic memory, seeds the next, so that the chain drifts as spontaneous thought does
\cite{christoff2016mind}. The voice has access only to memory queries and to the architecture's imagination tools,
which simulate observations and outcomes on a copy of the belief graph; it can neither act nor write beliefs. Its
products are of three kinds: counterfactuals over past episodes (``had the door been open, the crossing would have
taken 40\,s less''), hypotheses to be tested later, and prospective plans (``if the door is shut on the way back,
the route through the next room is shorter''). Functionally this is the default network's role in simulating
past and future \cite{buckner2008brain}, and the offline replay and model reduction that Active Inference assigns
to rest and sleep \cite{friston2017curiosity, hobson2012waking}. Replay is prioritised by the expected gain of the
memory it revisits \cite{mattar2018prioritized}.

\paragraph{Source monitoring.}
Humans occasionally misremember imagined events as perceived ones \cite{johnson1993source}. The robot must not.
Everything produced while wandering is stored with an \emph{imagined} provenance, separate from the factual
episodic record, and enters the learning loop only as a hypothesis to be tested, never as evidence.

\subsection{What the return to reality gives the automatic layer}
\label{sec:inner_speech:return}

The transition from \textsc{wandering} to \textsc{focused} cancels the chain in flight (its state is kept, so that
it may resume if demand subsides) and gives the full voice to the surprise. It serves the automatic layer at three
timescales.
\begin{enumerate}
    \item \textbf{Decision.} A surprise means the automatic layer has left its domain of competence, where its only
    responses are to retry, stop or fail. The voice adds re-interpretation of the scene, arbitration between goals,
    re-planning through remembered topology, and recourse to a person.
    \item \textbf{Attention.} The voice can raise precision where it matters: orienting sensors, increasing a
    detector's rate in a region, or passing a hypothesis to the concept agents as a spatial prior. Such a prior has
    bounded weight; it may accelerate the recognition of an object that is present but may not create one that is
    not, preserving the grounding constraint.
    \item \textbf{Teaching.} Each resolved surprise is a labelled example the automatic layer could not have
    produced alone: a residual later identified as a chair is a false negative of the chair detector; a door that
    is repeatedly shut in the evening becomes a prior; a recovery that repeatedly succeeds becomes a branch of the
    controller's behaviour tree. The domain of competence grows, and the situations that require the voice become
    rarer, which is the classical progression from effortful to autonomous skill \cite{fitts1967human}.
\end{enumerate}
Plans prepared while wandering act as a cache: when the anticipated situation occurs, the focused voice verifies a
plan instead of composing one. The return ends with an explicit release --- spatial priors cleared, safety margins
restored, the automatic layer confirmed back within its competence --- without which every surprise would leave
the robot permanently vigilant.

\subsection{Relation to chain-of-thought reasoning}
\label{sec:inner_speech:cot}

Chain-of-thought is inner speech in the functional sense, but of a specific kind, and the comparison sharpens the
design (Table~\ref{tab:inner_cot}). Plain chain-of-thought is convergent and ungrounded: tokens condition on
tokens until an answer is produced, and the trace is discarded. Interleaving reasoning with action and observation
grounds it \cite{yao2023react}; our focused loop is of this kind, with the additional property that the
observation, through \eqref{eq:inner_update}, rather than the next token, determines where the reasoning goes.
The wandering voice is closer to plain chain-of-thought in its isolation from perception, but divergent rather
than convergent, interruptible, and persistent in its products.

\begin{table}[ht]
    \centering
    \small
    \begin{tabular}{|>{\raggedright\arraybackslash}p{0.19\linewidth}|>{\raggedright\arraybackslash}p{0.22\linewidth}|>{\raggedright\arraybackslash}p{0.24\linewidth}|>{\raggedright\arraybackslash}p{0.24\linewidth}|}
        \hline
                                & \textbf{Chain-of-thought} & \textbf{Focused voice} & \textbf{Wandering voice} \\ \hline
        Goal                    & answer a given question   & resolve a given surprise & none given; seed by salience \\ \hline
        Dynamics                & convergent                & convergent             & divergent, associative \\ \hline
        Contact with the world  & none                      & thought--action--observation & none (imagination only) \\ \hline
        Persistence             & discarded                 & question objects, beliefs, lessons & imagined store, hypotheses \\ \hline
        Interruptible           & no                        & yes                    & yes, resumable \\ \hline
        Latency budget          & n/a                       & $\sim$1\,s, no extended thinking & unbounded, extended thinking \\ \hline
    \end{tabular}
    \caption{Chain-of-thought compared with the two modes of the proposed inner voice.}
    \label{tab:inner_cot}
\end{table}

Two practical consequences follow. First, \emph{think when there is time, not when there is pressure}. In the
deployed language agent extended thinking is disabled, because a thinking-mode model spent 3\,000--6\,000 tokens
(40--80\,s) before each tool call; that is incompatible with the focused mode but costs nothing in the wandering
mode, where it can be enabled with a large budget and pre-empted at any moment. Slow reasoning is done in advance
and fast reasoning uses its results. Second, chain-of-thought traces are known not always to reflect the
computation that produced the answer \cite{turpin2023language}, as human verbal reports do not always reflect the
causes of behaviour \cite{nisbett1977telling}. The faithful record of the robot's reasoning is therefore the
structured one --- hypotheses, predictions, observations and posterior updates --- not the text. Finally, wandering
chains whose hypotheses were later confirmed by perception form a corpus of reasoning verified against the world,
which could be used to fine-tune the local model in the manner of self-taught reasoning \cite{zelikman2022star},
with perception, rather than a known answer, as the filter.

\subsection{Consequences for belief and learning}
\label{sec:inner_speech:learning}

The loop changes how beliefs evolve in two ways. Within an episode, uncertainty is reduced where surprise
occurred rather than where the sensors happened to point, which is efficient but leaves the rest of the map less
examined and makes the robot's own data a biased sample; a decaying novelty term \cite{schwartenbeck2019curiosity}
and importance weighting of learned statistics are required to correct this. Across episodes, every closed
question is a labelled example, and learning proceeds at four timescales (Table~\ref{tab:inner_learning}).

\begin{table}[ht]
    \centering
    \small
    \begin{tabular}{|>{\raggedright\arraybackslash}p{0.14\linewidth}|>{\raggedright\arraybackslash}p{0.37\linewidth}|>{\raggedright\arraybackslash}p{0.37\linewidth}|}
        \hline
        \textbf{Timescale} & \textbf{What changes} & \textbf{Where, in this architecture} \\ \hline
        Seconds   & posterior over the hypotheses of one question         & the question object \eqref{eq:inner_question} \\ \hline
        Days      & priors: where objects are, which doors are shut, what moves & long-term object and topological memory, as Dirichlet counts \\ \hline
        Weeks     & likelihoods: detector and concept models              & false negatives and positives labelled by resolved questions \\ \hline
        Slow      & structure: new categories                             & recurrent, mutually similar ``unknown'' resolutions named by the voice and proposed as a new concept \\ \hline
    \end{tabular}
    \caption{Learning driven by the inner-question loop.}
    \label{tab:inner_learning}
\end{table}

The last row is where the hypothesis of a language for thought is most directly engaged: a cluster of anomalies
that resists every existing concept becomes learnable once it is named. The loop also learns about itself: the
salience weights $\pi_k$ in \eqref{eq:inner_demand} are adjusted by the information each trigger actually yielded,
asking a person is weighed against looking, and loops that repeatedly resolve the same way are compiled into
reflexive rules, so that the number of language-model calls per hour should \emph{decrease} as competence grows.

\subsection{Why a language model: generation versus evaluation}
\label{sec:inner_speech:why_lm}

None of the above requires a person, so the question arises of what a language model contributes to an
architecture that already performs inference and action selection. Our answer is a division of labour. Expected
free energy evaluates policies and Bayesian updating weighs hypotheses, but both operate over a space that has
already been enumerated: each concept agent carries a fixed generative model, policies must be listed before they
can be scored, and Bayesian model reduction can prune structure but not invent it \cite{friston2017curiosity}.
Whatever falls outside the given models remains an unexplained residual. A language model is weak where this
machinery is strong --- its confidences are not calibrated posteriors and it has no access to the sensors --- and
strong where it is weak: it is a cheap, broad generator of hypotheses, explanations and plans. We therefore use it
only as a \emph{proposal distribution}, in a generate-and-test loop whose test is active inference:
\begin{itemize}
    \item \emph{Hypotheses outside the built-in models.} A residual blob can be proposed to be a bag, a toppled
    chair or a robot vacuum, none of which any concept agent represents; the question loop then tests them.
    \item \emph{Amortised policy search.} Instead of enumerating multi-step policies, the model writes a few
    candidate missions --- small programs in the mission language --- which the gate validates and expected free
    energy scores.
    \item \emph{Relational binding across agents.} The agents' representations (signed distance fields, occupancy
    grids, existence log-odds, episodic records) are mutually incompatible; a sentence such as ``the chair that was
    not next to the table in the room behind door~3 yesterday'' binds them in one query no single agent could pose.
    \item \emph{Explanation from few examples.} After a single surprise the model can propose which variable
    mattered (``the door is shut because it is evening''), turning one episode into a testable generalisation.
    \item \emph{Abstraction and naming.} Episodes are compressed into reusable lessons, and recurrent unknowns are
    named, which is what makes them candidates for a new concept (Table~\ref{tab:inner_learning}).
\end{itemize}
Language, rather than an arbitrary learned embedding, matters for three reasons: it is discrete and compositional,
so new hypotheses and plans are assembled from known parts; it is inspectable, so that deterministic gates can parse
and check it and people can audit it; and, once the robot thinks in it, communication with a person comes at no
further cost.

\paragraph{Where the knowledge comes from.}
This raises a difficulty for the motivating thesis. If the inner voice precedes the outer one, the common-sense
knowledge on which the proposer's value rests cannot have arrived through communication --- yet a language model
acquires all of it from human text, that is, from other people's outer voices. The difficulty dissolves once the
\emph{faculty} is distinguished from its \emph{content}. The internalist claim concerns the former: a
computational system for building unbounded hierarchical structures, available for thought before it is used for
talk \cite{hauser2002faculty, berwick2016why}. It makes no claim that what a person knows was acquired without
language. Human knowledge has three sources: core systems for objects, agents, number and space, present before
learning \cite{spelke2007core}; individual experience; and, once outer language exists, cumulative cultural
transmission, which supplies most of what an adult knows \cite{tomasello1999cultural}. Even on the internalist
view, the inner voice of a modern human is filled largely by the outer voices of others.

The architecture has the same three sources, in a different arrangement. The concept agents' hand-built generative
models play the role of \emph{core knowledge}: given before experience and limited to a few domains. The robot's
own counts $n_x(h)$ in \eqref{eq:inner_prior} are \emph{individual experience}. The language model is
\emph{culture}: knowledge distilled from a community's communication and inherited, for this individual, at birth.
The robot thus resembles a child born into a culture more than an early hominin, and its inner voice is Vygotskian
in origin: its format and its content were both learned from outer speech. Two consequences follow. First, the
proposal does not demonstrate the internalist thesis and should not be read as doing so; its claim is functional
--- an inner use of language has epistemic value independently of communication --- and holds whatever the origin of
the language system. Second, the format and the content of the voice can in principle be separated, which turns the
question into an experiment (P7 in Section~\ref{sec:inner_speech:predictions}): a compositional generator over the
robot's own concepts and mission language, trained only on the robot's experience, in the manner of library
learning in program induction \cite{ellis2021dreamcoder}, would supply the format without the culture.

\subsection{Mapping onto the existing architecture}
\label{sec:inner_speech:mapping}

Several components already exist in the language agent. Robot events (mission reports, failures, completions)
already wake the model through a queue, in turns restricted to perceiving, querying memory, stopping and
proposing, never acting. Every action already passes a deterministic gate. The decision to ask or to assume is
already decision-theoretic. Unexplained occupancy is already isolated by a residual agent that acts as the null
hypothesis of the scene, and episodic and long-term object memories are already joined read-only. Missing are the
persistent question objects \eqref{eq:inner_question}, a prediction evaluator implementing
\eqref{eq:inner_update} outside the language model, the demand signal and mode arbitration
\eqref{eq:inner_demand}--\eqref{eq:inner_modes}, the imagined store with its provenance, and the channel through
which bounded spatial priors reach the concept agents.

\subsection{Testable predictions}
\label{sec:inner_speech:predictions}

Consistent with the evaluation principle adopted throughout --- policies compared against each other inside one
closed loop, with body, sensors and environment fixed --- the proposal makes seven falsifiable predictions, each
tested in simulation by moving known objects and introducing unfamiliar ones:
\begin{enumerate}[label=\textbf{P\arabic*}]
    \item The unexplained residual mass per room decreases faster with the focused loop than without it.
    \item The proposer's calibration on resolved questions improves, and $\beta_x$ in \eqref{eq:inner_prior}
    decreases, as the robot's own counts accumulate.
    \item Detector false-negative rates fall after retraining on question-labelled examples.
    \item Inner questions per hour decrease over repeated exposure to the same environment while answer accuracy
    is maintained.
    \item Plans prepared while wandering reduce the latency from surprise to a validated response, and the time to
    notice a deliberately moved object does not grow despite habituation.
    \item With the language-model proposer replaced by a fixed hypothesis set (the existing concept classes and
    ``unknown'') and template-based missions, the loop resolves fewer questions about objects and failures outside
    the built-in models, and acquires new categories more slowly.
    \item A compositional generator trained only on the robot's own experience (format without culture) recovers
    part of the advantage in P6 within a familiar environment, and less of it the further a situation departs from
    what the robot has seen.
\end{enumerate}
A failure of P1 would indicate that the epistemic value of language-mediated questions does not exceed that of
the existing exploration; a failure of P4 would indicate that the loop does not transfer competence to the
automatic layer, which is its principal claim; and a failure of P6 would indicate that the language model adds
nothing a fixed hypothesis set does not, which would remove the case for including it. P7 separates the
contribution of language as a format for thought from that of the knowledge it carries.

\subsection{Limitations}
\label{sec:inner_speech:limitations}

The demand signal \eqref{eq:inner_demand} sums heterogeneous prediction errors whose precisions must themselves be
calibrated, and a poorly scaled term will either suppress the voice or keep it permanently engaged. A single
language model serves the person, the focused loop and the wandering loop; pre-emption keeps the person first but
bounds how much offline reasoning a busy robot can do. The tempered prior \eqref{eq:inner_prior} is a heuristic
combination rather than a derived posterior. The generator of P7 is substantially harder to build than the
ablation of P6, and may have to be approximated by a small model trained on the robot's logs and mission corpus,
whose pre-training would leak some culture back in. Finally, the analogy with human inner speech is functional, not
mechanistic: the proposal is that the same computational role --- serial, compositional self-questioning in the
service of epistemic value --- is useful for a robot, not that the robot's voice resembles the human one in any
further respect.

% and add inner_speech.bib as a second resource:  \addbibresource{inner_speech.bib}
% All keys are new (none collides with the keys main.tex already cites).
% Packages used: amsmath, amssymb, enumitem (already loaded by main.tex) and array, which main.tex does not
% load yet: add \usepackage{array} to its preamble.
% =====================================================================================================

\section{Inner Speech as Epistemic Action}
\label{sec:inner_speech}

The robot described so far uses language in one direction only: a person speaks, and a language agent turns
the request into commands for the controller, grounded in a textual digest of the belief graph. This section
proposes to turn that capability inward. We argue that a language model embedded in an Active Inference
architecture can play the role that inner speech plays in human cognition: a serial, compositional medium in
which the agent poses questions to itself, generates hypotheses, rehearses plans and replays its past, and that
this role follows naturally from the epistemic component of expected free energy. The proposal is stated
together with the constraints that keep such a voice from becoming a source of confabulated belief, and with the
experiments that would falsify it.

\subsection{Motivation: language for thought before language for talk}
\label{sec:inner_speech:motivation}

Two traditions disagree on the origin of inner speech but agree on its function. On the internalist view,
human language evolved first as an internal system for structuring thought, and its use for communication is a
later externalisation \cite{berwick2016why}. On the developmental view, inner speech is social speech that the
child gradually internalises, abbreviates and turns into a tool for self-regulation \cite{vygotsky1986thought}.
Contemporary reviews treat inner speech as functionally involved in planning, problem solving, self-monitoring and
the maintenance of goals in working memory \cite{aldersonday2015inner}, and language more generally as a
cognitive technology that lets a learner manipulate discrete, recombinable tokens rather than raw percepts
\cite{clark1998magic}.

Large language models give a striking instance of the developmental route. They are trained exclusively on
communicative text, yet producing intermediate reasoning in that same medium before answering --- chain-of-thought
--- measurably improves their performance \cite{wei2022chain}. A system built for talk becomes better at thinking
by talking to itself. Our language agent sits at the same point of its own development: it was built to converse
with a person, and the question we ask is what it contributes once its speech is addressed to the robot itself.
Section~\ref{sec:inner_speech:why_lm} returns to what this origin implies for the internalist view.

\subsection{Inner questions as epistemic actions}
\label{sec:inner_speech:questions}

In Active Inference, policies are scored by expected free energy, whose epistemic term rewards the expected
reduction of uncertainty about hidden states \cite{friston2017process, parr2022active}. A question is the
linguistic form of exactly this quantity: it names a hidden variable whose value is uncertain and whose resolution
is expected to matter. We therefore model an inner question as a persistent object
\begin{equation}
    \mathcal{Q} = \big\langle\, e,\; \mathcal{H},\; q(h),\; \mathcal{P},\; b,\; s \,\big\rangle ,
    \label{eq:inner_question}
\end{equation}
where $e$ is the triggering event (for example an obstacle no concept agent explains), $\mathcal{H}$ a finite set
of named hypotheses (``a chair'', ``a box'', ``unknown''), $q(h)$ the current posterior over them, $\mathcal{P}$
the set of pending predictions, $b$ a budget in time and motion, and $s \in \{\text{open},\text{resolved},
\text{abandoned}\}$ its status. Questions form a stack in the sense of a \emph{question under discussion}
\cite{roberts2012information}: ``is it a chair?'' is a sub-question whose answer partially resolves the parent
``what is this?''. The question, not the language model's context window, carries the state of the reasoning from
one step to the next.

\paragraph{One step of the loop.}
Given $\mathcal{Q}$, the language model proposes a test: an epistemic action $a$ (turn towards the anomaly,
approach it, look with the front camera) together with the outcome each hypothesis predicts. The value of the test
is its expected information gain minus its cost,
\begin{equation}
    V(a) \;=\; \underbrace{H\!\left[q(h)\right] \;-\; \mathbb{E}_{p(o \mid a)}\!\left[ H\!\left[q(h \mid o, a)\right]
    \right]}_{I(h;\,o \mid a)} \;-\; c(a),
    \label{eq:inner_value}
\end{equation}
which is the negative of the epistemic part of expected free energy with an explicit action cost. The test is
executed only if $V(a) > 0$ and the action passes the same deterministic gate that already validates every
command the language agent issues. When the observation $o$ arrives, the posterior is updated by the concept
agents' own likelihoods,
\begin{equation}
    q(h \mid o, a) \;\propto\; p(o \mid h, a)\, q(h),
    \label{eq:inner_update}
\end{equation}
and the loop continues, spawns a sub-question, or terminates when
\begin{equation}
    H\!\left[q(h)\right] < \eta \quad\lor\quad \max_a V(a) \le 0 \quad\lor\quad b \text{ exhausted}.
    \label{eq:inner_stop}
\end{equation}
The stopping rule is what separates reasoning from rumination: a loop with no expected information to gain closes,
possibly with the answer ``unknown'', and its trigger is habituated for a refractory period.

\paragraph{The grounding constraint.}
The central design rule is that \emph{the language model proposes and perception disposes}. Text produced by the
model may introduce hypotheses, predictions and tests; it may never move a posterior. Only the likelihood
$p(o \mid h, a)$ in \eqref{eq:inner_update}, computed by the perceptual agents from sensor data, does. This is the
belief-side counterpart of the rule the language agent already enforces for actions, where every proposed command
is validated in C++ against the belief graph before it reaches the controller. Without it, a self-addressed
language loop is a closed system of text conditioned on text and converges on its own confabulations.

\paragraph{The proposer as a calibrated prior.}
The hypotheses the model names, and the probabilities it attaches to them, encode world knowledge the robot has
not acquired by itself (``chairs come with tables''). We treat this as a prior whose weight is earned. Let
$\pi_{\text{LM}}(h \mid x)$ be the proposer's distribution in context $x$ and $n_x(h)$ the robot's own counts of
resolved questions in that context. The prior entering \eqref{eq:inner_update} is
\begin{equation}
    q_0(h \mid x) \;\propto\; \pi_{\text{LM}}(h \mid x)^{\beta_x}\, \big(\alpha + n_x(h)\big),
    \label{eq:inner_prior}
\end{equation}
a tempered product of the proposer and a Dirichlet over experience, where $\beta_x \in [0,1]$ is set from the
proposer's measured calibration in that context (for instance from its Brier score on past questions). Early in
the robot's life the proposer dominates; as counts accumulate, the robot's own experience replaces borrowed
knowledge. The same conjugate form already governs the language agent's decision to ask a person or to assume
that they moved an object, where a Beta posterior over the person's past answers is compared against the costs of
asking and of being wrong.

\subsection{Attention, automaticity and the wandering mind}
\label{sec:inner_speech:modes}

A human driver on an easy road steers without attending to it; the mind is free to plan, imagine or recall, and is
recaptured the instant something unexpected happens \cite{smallwood2015science, christoff2016mind}. In Active
Inference this is a statement about precision: while prediction errors are small and the habitual policy is
reliable, attention --- the precision assigned to sensory prediction errors --- can be withdrawn from the task; a
surprise raises it again \cite{feldman2010attention}. The arbitration between habitual and deliberative control by
their relative uncertainty \cite{daw2005uncertainty} has the same structure.

The architecture already separates the two systems. The controller and the concept agents form an automatic layer
that drives, avoids obstacles and executes missions without the language model; the language agent is the
deliberative voice. We propose to make the allocation of the voice explicit through a scalar \emph{demand} signal,
computed on every digest cycle without the language model,
\begin{equation}
    D_t \;=\; \sum_{k} \pi_k\, \varepsilon_{k,t}^{2} \;+\; \lambda\, \mathbb{1}\!\left[\text{person active}_t\right],
    \label{eq:inner_demand}
\end{equation}
where $\varepsilon_{k,t}$ are normalised prediction errors reported by the agents (unexplained residual occupancy,
flips of existence or door beliefs, localisation uncertainty, stalled or failed mission steps, controller refusals)
and $\pi_k$ their precisions. Three modes follow from two thresholds with asymmetric hysteresis:
\begin{equation}
    m_t =
    \begin{cases}
        \textsc{focused}    & D_t > \theta_{\text{hi}},\\[2pt]
        \textsc{wandering}  & D_\tau < \theta_{\text{lo}} \;\; \forall\, \tau \in [t - T, t],\\[2pt]
        \textsc{monitoring} & \text{otherwise.}
    \end{cases}
    \label{eq:inner_modes}
\end{equation}
Entering \textsc{wandering} requires demand to stay low for a dwell time $T$; leaving it requires a single sample
above $\theta_{\text{hi}}$. The robot is slow to drift off and instant to snap back. Safety never depends on this
arbitration: the controller's reflexes act on their own, and the deliberative voice arrives after the reflex, as
awareness follows braking in the driver.

\paragraph{Focused mode} runs the question loop of Section~\ref{sec:inner_speech:questions} on the triggering
event. \paragraph{Monitoring mode} makes no language-model calls; the voice is primed but silent.

\paragraph{Wandering mode.}
Freed from the task, the voice is seeded with one recent experience chosen by salience: an episode that ended
unresolved, a failure, an object that changed, or an upcoming goal. Each step's output, together with an
associative retrieval from episodic memory, seeds the next, so that the chain drifts as spontaneous thought does
\cite{christoff2016mind}. The voice has access only to memory queries and to the architecture's imagination tools,
which simulate observations and outcomes on a copy of the belief graph; it can neither act nor write beliefs. Its
products are of three kinds: counterfactuals over past episodes (``had the door been open, the crossing would have
taken 40\,s less''), hypotheses to be tested later, and prospective plans (``if the door is shut on the way back,
the route through the next room is shorter''). Functionally this is the default network's role in simulating
past and future \cite{buckner2008brain}, and the offline replay and model reduction that Active Inference assigns
to rest and sleep \cite{friston2017curiosity, hobson2012waking}. Replay is prioritised by the expected gain of the
memory it revisits \cite{mattar2018prioritized}.

\paragraph{Source monitoring.}
Humans occasionally misremember imagined events as perceived ones \cite{johnson1993source}. The robot must not.
Everything produced while wandering is stored with an \emph{imagined} provenance, separate from the factual
episodic record, and enters the learning loop only as a hypothesis to be tested, never as evidence.

\subsection{What the return to reality gives the automatic layer}
\label{sec:inner_speech:return}

The transition from \textsc{wandering} to \textsc{focused} cancels the chain in flight (its state is kept, so that
it may resume if demand subsides) and gives the full voice to the surprise. It serves the automatic layer at three
timescales.
\begin{enumerate}
    \item \textbf{Decision.} A surprise means the automatic layer has left its domain of competence, where its only
    responses are to retry, stop or fail. The voice adds re-interpretation of the scene, arbitration between goals,
    re-planning through remembered topology, and recourse to a person.
    \item \textbf{Attention.} The voice can raise precision where it matters: orienting sensors, increasing a
    detector's rate in a region, or passing a hypothesis to the concept agents as a spatial prior. Such a prior has
    bounded weight; it may accelerate the recognition of an object that is present but may not create one that is
    not, preserving the grounding constraint.
    \item \textbf{Teaching.} Each resolved surprise is a labelled example the automatic layer could not have
    produced alone: a residual later identified as a chair is a false negative of the chair detector; a door that
    is repeatedly shut in the evening becomes a prior; a recovery that repeatedly succeeds becomes a branch of the
    controller's behaviour tree. The domain of competence grows, and the situations that require the voice become
    rarer, which is the classical progression from effortful to autonomous skill \cite{fitts1967human}.
\end{enumerate}
Plans prepared while wandering act as a cache: when the anticipated situation occurs, the focused voice verifies a
plan instead of composing one. The return ends with an explicit release --- spatial priors cleared, safety margins
restored, the automatic layer confirmed back within its competence --- without which every surprise would leave
the robot permanently vigilant.

\subsection{Relation to chain-of-thought reasoning}
\label{sec:inner_speech:cot}

Chain-of-thought is inner speech in the functional sense, but of a specific kind, and the comparison sharpens the
design (Table~\ref{tab:inner_cot}). Plain chain-of-thought is convergent and ungrounded: tokens condition on
tokens until an answer is produced, and the trace is discarded. Interleaving reasoning with action and observation
grounds it \cite{yao2023react}; our focused loop is of this kind, with the additional property that the
observation, through \eqref{eq:inner_update}, rather than the next token, determines where the reasoning goes.
The wandering voice is closer to plain chain-of-thought in its isolation from perception, but divergent rather
than convergent, interruptible, and persistent in its products.

\begin{table}[ht]
    \centering
    \small
    \begin{tabular}{|>{\raggedright\arraybackslash}p{0.19\linewidth}|>{\raggedright\arraybackslash}p{0.22\linewidth}|>{\raggedright\arraybackslash}p{0.24\linewidth}|>{\raggedright\arraybackslash}p{0.24\linewidth}|}
        \hline
                                & \textbf{Chain-of-thought} & \textbf{Focused voice} & \textbf{Wandering voice} \\ \hline
        Goal                    & answer a given question   & resolve a given surprise & none given; seed by salience \\ \hline
        Dynamics                & convergent                & convergent             & divergent, associative \\ \hline
        Contact with the world  & none                      & thought--action--observation & none (imagination only) \\ \hline
        Persistence             & discarded                 & question objects, beliefs, lessons & imagined store, hypotheses \\ \hline
        Interruptible           & no                        & yes                    & yes, resumable \\ \hline
        Latency budget          & n/a                       & $\sim$1\,s, no extended thinking & unbounded, extended thinking \\ \hline
    \end{tabular}
    \caption{Chain-of-thought compared with the two modes of the proposed inner voice.}
    \label{tab:inner_cot}
\end{table}

Two practical consequences follow. First, \emph{think when there is time, not when there is pressure}. In the
deployed language agent extended thinking is disabled, because a thinking-mode model spent 3\,000--6\,000 tokens
(40--80\,s) before each tool call; that is incompatible with the focused mode but costs nothing in the wandering
mode, where it can be enabled with a large budget and pre-empted at any moment. Slow reasoning is done in advance
and fast reasoning uses its results. Second, chain-of-thought traces are known not always to reflect the
computation that produced the answer \cite{turpin2023language}, as human verbal reports do not always reflect the
causes of behaviour \cite{nisbett1977telling}. The faithful record of the robot's reasoning is therefore the
structured one --- hypotheses, predictions, observations and posterior updates --- not the text. Finally, wandering
chains whose hypotheses were later confirmed by perception form a corpus of reasoning verified against the world,
which could be used to fine-tune the local model in the manner of self-taught reasoning \cite{zelikman2022star},
with perception, rather than a known answer, as the filter.

\subsection{Consequences for belief and learning}
\label{sec:inner_speech:learning}

The loop changes how beliefs evolve in two ways. Within an episode, uncertainty is reduced where surprise
occurred rather than where the sensors happened to point, which is efficient but leaves the rest of the map less
examined and makes the robot's own data a biased sample; a decaying novelty term \cite{schwartenbeck2019curiosity}
and importance weighting of learned statistics are required to correct this. Across episodes, every closed
question is a labelled example, and learning proceeds at four timescales (Table~\ref{tab:inner_learning}).

\begin{table}[ht]
    \centering
    \small
    \begin{tabular}{|>{\raggedright\arraybackslash}p{0.14\linewidth}|>{\raggedright\arraybackslash}p{0.37\linewidth}|>{\raggedright\arraybackslash}p{0.37\linewidth}|}
        \hline
        \textbf{Timescale} & \textbf{What changes} & \textbf{Where, in this architecture} \\ \hline
        Seconds   & posterior over the hypotheses of one question         & the question object \eqref{eq:inner_question} \\ \hline
        Days      & priors: where objects are, which doors are shut, what moves & long-term object and topological memory, as Dirichlet counts \\ \hline
        Weeks     & likelihoods: detector and concept models              & false negatives and positives labelled by resolved questions \\ \hline
        Slow      & structure: new categories                             & recurrent, mutually similar ``unknown'' resolutions named by the voice and proposed as a new concept \\ \hline
    \end{tabular}
    \caption{Learning driven by the inner-question loop.}
    \label{tab:inner_learning}
\end{table}

The last row is where the hypothesis of a language for thought is most directly engaged: a cluster of anomalies
that resists every existing concept becomes learnable once it is named. The loop also learns about itself: the
salience weights $\pi_k$ in \eqref{eq:inner_demand} are adjusted by the information each trigger actually yielded,
asking a person is weighed against looking, and loops that repeatedly resolve the same way are compiled into
reflexive rules, so that the number of language-model calls per hour should \emph{decrease} as competence grows.

\subsection{Why a language model: generation versus evaluation}
\label{sec:inner_speech:why_lm}

None of the above requires a person, so the question arises of what a language model contributes to an
architecture that already performs inference and action selection. Our answer is a division of labour. Expected
free energy evaluates policies and Bayesian updating weighs hypotheses, but both operate over a space that has
already been enumerated: each concept agent carries a fixed generative model, policies must be listed before they
can be scored, and Bayesian model reduction can prune structure but not invent it \cite{friston2017curiosity}.
Whatever falls outside the given models remains an unexplained residual. A language model is weak where this
machinery is strong --- its confidences are not calibrated posteriors and it has no access to the sensors --- and
strong where it is weak: it is a cheap, broad generator of hypotheses, explanations and plans. We therefore use it
only as a \emph{proposal distribution}, in a generate-and-test loop whose test is active inference:
\begin{itemize}
    \item \emph{Hypotheses outside the built-in models.} A residual blob can be proposed to be a bag, a toppled
    chair or a robot vacuum, none of which any concept agent represents; the question loop then tests them.
    \item \emph{Amortised policy search.} Instead of enumerating multi-step policies, the model writes a few
    candidate missions --- small programs in the mission language --- which the gate validates and expected free
    energy scores.
    \item \emph{Relational binding across agents.} The agents' representations (signed distance fields, occupancy
    grids, existence log-odds, episodic records) are mutually incompatible; a sentence such as ``the chair that was
    not next to the table in the room behind door~3 yesterday'' binds them in one query no single agent could pose.
    \item \emph{Explanation from few examples.} After a single surprise the model can propose which variable
    mattered (``the door is shut because it is evening''), turning one episode into a testable generalisation.
    \item \emph{Abstraction and naming.} Episodes are compressed into reusable lessons, and recurrent unknowns are
    named, which is what makes them candidates for a new concept (Table~\ref{tab:inner_learning}).
\end{itemize}
Language, rather than an arbitrary learned embedding, matters for three reasons: it is discrete and compositional,
so new hypotheses and plans are assembled from known parts; it is inspectable, so that deterministic gates can parse
and check it and people can audit it; and, once the robot thinks in it, communication with a person comes at no
further cost.

\paragraph{Where the knowledge comes from.}
This raises a difficulty for the motivating thesis. If the inner voice precedes the outer one, the common-sense
knowledge on which the proposer's value rests cannot have arrived through communication --- yet a language model
acquires all of it from human text, that is, from other people's outer voices. The difficulty dissolves once the
\emph{faculty} is distinguished from its \emph{content}. The internalist claim concerns the former: a
computational system for building unbounded hierarchical structures, available for thought before it is used for
talk \cite{hauser2002faculty, berwick2016why}. It makes no claim that what a person knows was acquired without
language. Human knowledge has three sources: core systems for objects, agents, number and space, present before
learning \cite{spelke2007core}; individual experience; and, once outer language exists, cumulative cultural
transmission, which supplies most of what an adult knows \cite{tomasello1999cultural}. Even on the internalist
view, the inner voice of a modern human is filled largely by the outer voices of others.

The architecture has the same three sources, in a different arrangement. The concept agents' hand-built generative
models play the role of \emph{core knowledge}: given before experience and limited to a few domains. The robot's
own counts $n_x(h)$ in \eqref{eq:inner_prior} are \emph{individual experience}. The language model is
\emph{culture}: knowledge distilled from a community's communication and inherited, for this individual, at birth.
The robot thus resembles a child born into a culture more than an early hominin, and its inner voice is Vygotskian
in origin: its format and its content were both learned from outer speech. Two consequences follow. First, the
proposal does not demonstrate the internalist thesis and should not be read as doing so; its claim is functional
--- an inner use of language has epistemic value independently of communication --- and holds whatever the origin of
the language system. Second, the format and the content of the voice can in principle be separated, which turns the
question into an experiment (P7 in Section~\ref{sec:inner_speech:predictions}): a compositional generator over the
robot's own concepts and mission language, trained only on the robot's experience, in the manner of library
learning in program induction \cite{ellis2021dreamcoder}, would supply the format without the culture.

\subsection{Mapping onto the existing architecture}
\label{sec:inner_speech:mapping}

Several components already exist in the language agent. Robot events (mission reports, failures, completions)
already wake the model through a queue, in turns restricted to perceiving, querying memory, stopping and
proposing, never acting. Every action already passes a deterministic gate. The decision to ask or to assume is
already decision-theoretic. Unexplained occupancy is already isolated by a residual agent that acts as the null
hypothesis of the scene, and episodic and long-term object memories are already joined read-only. Missing are the
persistent question objects \eqref{eq:inner_question}, a prediction evaluator implementing
\eqref{eq:inner_update} outside the language model, the demand signal and mode arbitration
\eqref{eq:inner_demand}--\eqref{eq:inner_modes}, the imagined store with its provenance, and the channel through
which bounded spatial priors reach the concept agents.

\subsection{Testable predictions}
\label{sec:inner_speech:predictions}

Consistent with the evaluation principle adopted throughout --- policies compared against each other inside one
closed loop, with body, sensors and environment fixed --- the proposal makes seven falsifiable predictions, each
tested in simulation by moving known objects and introducing unfamiliar ones:
\begin{enumerate}[label=\textbf{P\arabic*}]
    \item The unexplained residual mass per room decreases faster with the focused loop than without it.
    \item The proposer's calibration on resolved questions improves, and $\beta_x$ in \eqref{eq:inner_prior}
    decreases, as the robot's own counts accumulate.
    \item Detector false-negative rates fall after retraining on question-labelled examples.
    \item Inner questions per hour decrease over repeated exposure to the same environment while answer accuracy
    is maintained.
    \item Plans prepared while wandering reduce the latency from surprise to a validated response, and the time to
    notice a deliberately moved object does not grow despite habituation.
    \item With the language-model proposer replaced by a fixed hypothesis set (the existing concept classes and
    ``unknown'') and template-based missions, the loop resolves fewer questions about objects and failures outside
    the built-in models, and acquires new categories more slowly.
    \item A compositional generator trained only on the robot's own experience (format without culture) recovers
    part of the advantage in P6 within a familiar environment, and less of it the further a situation departs from
    what the robot has seen.
\end{enumerate}
A failure of P1 would indicate that the epistemic value of language-mediated questions does not exceed that of
the existing exploration; a failure of P4 would indicate that the loop does not transfer competence to the
automatic layer, which is its principal claim; and a failure of P6 would indicate that the language model adds
nothing a fixed hypothesis set does not, which would remove the case for including it. P7 separates the
contribution of language as a format for thought from that of the knowledge it carries.

\subsection{Limitations}
\label{sec:inner_speech:limitations}

The demand signal \eqref{eq:inner_demand} sums heterogeneous prediction errors whose precisions must themselves be
calibrated, and a poorly scaled term will either suppress the voice or keep it permanently engaged. A single
language model serves the person, the focused loop and the wandering loop; pre-emption keeps the person first but
bounds how much offline reasoning a busy robot can do. The tempered prior \eqref{eq:inner_prior} is a heuristic
combination rather than a derived posterior. The generator of P7 is substantially harder to build than the
ablation of P6, and may have to be approximated by a small model trained on the robot's logs and mission corpus,
whose pre-training would leak some culture back in. Finally, the analogy with human inner speech is functional, not
mechanistic: the proposal is that the same computational role --- serial, compositional self-questioning in the
service of epistemic value --- is useful for a robot, not that the robot's voice resembles the human one in any
further respect.

%     % =====================================================================================================
% language_hierarchy_section.tex — Language over the belief hierarchy: culture, sequencing and imagination
%
% DRAFT (2026-11-19). A proposal section: the architecture it describes is NOT yet implemented beyond the parts
% stated as existing. Source design note: docs/LANGUAGE_OVER_BELIEF_HIERARCHY.md.
% Include AFTER inner_speech_section.tex (it refers to that section's equations and modes):
%     % =====================================================================================================
% inner_speech_section.tex — Inner speech, mind-wandering and active inference
%
% DRAFT (2026-11-18). A proposal section: it describes an architecture that is NOT yet implemented,
% grounded in the existing language_agent (robocomp/active_inference). Include from main.tex with
%     \input{inner_speech_section}
% and add inner_speech.bib as a second resource:  \addbibresource{inner_speech.bib}
% All keys are new (none collides with the keys main.tex already cites).
% Packages used: amsmath, amssymb, enumitem (already loaded by main.tex) and array, which main.tex does not
% load yet: add \usepackage{array} to its preamble.
% =====================================================================================================

\section{Inner Speech as Epistemic Action}
\label{sec:inner_speech}

The robot described so far uses language in one direction only: a person speaks, and a language agent turns
the request into commands for the controller, grounded in a textual digest of the belief graph. This section
proposes to turn that capability inward. We argue that a language model embedded in an Active Inference
architecture can play the role that inner speech plays in human cognition: a serial, compositional medium in
which the agent poses questions to itself, generates hypotheses, rehearses plans and replays its past, and that
this role follows naturally from the epistemic component of expected free energy. The proposal is stated
together with the constraints that keep such a voice from becoming a source of confabulated belief, and with the
experiments that would falsify it.

\subsection{Motivation: language for thought before language for talk}
\label{sec:inner_speech:motivation}

Two traditions disagree on the origin of inner speech but agree on its function. On the internalist view,
human language evolved first as an internal system for structuring thought, and its use for communication is a
later externalisation \cite{berwick2016why}. On the developmental view, inner speech is social speech that the
child gradually internalises, abbreviates and turns into a tool for self-regulation \cite{vygotsky1986thought}.
Contemporary reviews treat inner speech as functionally involved in planning, problem solving, self-monitoring and
the maintenance of goals in working memory \cite{aldersonday2015inner}, and language more generally as a
cognitive technology that lets a learner manipulate discrete, recombinable tokens rather than raw percepts
\cite{clark1998magic}.

Large language models give a striking instance of the developmental route. They are trained exclusively on
communicative text, yet producing intermediate reasoning in that same medium before answering --- chain-of-thought
--- measurably improves their performance \cite{wei2022chain}. A system built for talk becomes better at thinking
by talking to itself. Our language agent sits at the same point of its own development: it was built to converse
with a person, and the question we ask is what it contributes once its speech is addressed to the robot itself.
Section~\ref{sec:inner_speech:why_lm} returns to what this origin implies for the internalist view.

\subsection{Inner questions as epistemic actions}
\label{sec:inner_speech:questions}

In Active Inference, policies are scored by expected free energy, whose epistemic term rewards the expected
reduction of uncertainty about hidden states \cite{friston2017process, parr2022active}. A question is the
linguistic form of exactly this quantity: it names a hidden variable whose value is uncertain and whose resolution
is expected to matter. We therefore model an inner question as a persistent object
\begin{equation}
    \mathcal{Q} = \big\langle\, e,\; \mathcal{H},\; q(h),\; \mathcal{P},\; b,\; s \,\big\rangle ,
    \label{eq:inner_question}
\end{equation}
where $e$ is the triggering event (for example an obstacle no concept agent explains), $\mathcal{H}$ a finite set
of named hypotheses (``a chair'', ``a box'', ``unknown''), $q(h)$ the current posterior over them, $\mathcal{P}$
the set of pending predictions, $b$ a budget in time and motion, and $s \in \{\text{open},\text{resolved},
\text{abandoned}\}$ its status. Questions form a stack in the sense of a \emph{question under discussion}
\cite{roberts2012information}: ``is it a chair?'' is a sub-question whose answer partially resolves the parent
``what is this?''. The question, not the language model's context window, carries the state of the reasoning from
one step to the next.

\paragraph{One step of the loop.}
Given $\mathcal{Q}$, the language model proposes a test: an epistemic action $a$ (turn towards the anomaly,
approach it, look with the front camera) together with the outcome each hypothesis predicts. The value of the test
is its expected information gain minus its cost,
\begin{equation}
    V(a) \;=\; \underbrace{H\!\left[q(h)\right] \;-\; \mathbb{E}_{p(o \mid a)}\!\left[ H\!\left[q(h \mid o, a)\right]
    \right]}_{I(h;\,o \mid a)} \;-\; c(a),
    \label{eq:inner_value}
\end{equation}
which is the negative of the epistemic part of expected free energy with an explicit action cost. The test is
executed only if $V(a) > 0$ and the action passes the same deterministic gate that already validates every
command the language agent issues. When the observation $o$ arrives, the posterior is updated by the concept
agents' own likelihoods,
\begin{equation}
    q(h \mid o, a) \;\propto\; p(o \mid h, a)\, q(h),
    \label{eq:inner_update}
\end{equation}
and the loop continues, spawns a sub-question, or terminates when
\begin{equation}
    H\!\left[q(h)\right] < \eta \quad\lor\quad \max_a V(a) \le 0 \quad\lor\quad b \text{ exhausted}.
    \label{eq:inner_stop}
\end{equation}
The stopping rule is what separates reasoning from rumination: a loop with no expected information to gain closes,
possibly with the answer ``unknown'', and its trigger is habituated for a refractory period.

\paragraph{The grounding constraint.}
The central design rule is that \emph{the language model proposes and perception disposes}. Text produced by the
model may introduce hypotheses, predictions and tests; it may never move a posterior. Only the likelihood
$p(o \mid h, a)$ in \eqref{eq:inner_update}, computed by the perceptual agents from sensor data, does. This is the
belief-side counterpart of the rule the language agent already enforces for actions, where every proposed command
is validated in C++ against the belief graph before it reaches the controller. Without it, a self-addressed
language loop is a closed system of text conditioned on text and converges on its own confabulations.

\paragraph{The proposer as a calibrated prior.}
The hypotheses the model names, and the probabilities it attaches to them, encode world knowledge the robot has
not acquired by itself (``chairs come with tables''). We treat this as a prior whose weight is earned. Let
$\pi_{\text{LM}}(h \mid x)$ be the proposer's distribution in context $x$ and $n_x(h)$ the robot's own counts of
resolved questions in that context. The prior entering \eqref{eq:inner_update} is
\begin{equation}
    q_0(h \mid x) \;\propto\; \pi_{\text{LM}}(h \mid x)^{\beta_x}\, \big(\alpha + n_x(h)\big),
    \label{eq:inner_prior}
\end{equation}
a tempered product of the proposer and a Dirichlet over experience, where $\beta_x \in [0,1]$ is set from the
proposer's measured calibration in that context (for instance from its Brier score on past questions). Early in
the robot's life the proposer dominates; as counts accumulate, the robot's own experience replaces borrowed
knowledge. The same conjugate form already governs the language agent's decision to ask a person or to assume
that they moved an object, where a Beta posterior over the person's past answers is compared against the costs of
asking and of being wrong.

\subsection{Attention, automaticity and the wandering mind}
\label{sec:inner_speech:modes}

A human driver on an easy road steers without attending to it; the mind is free to plan, imagine or recall, and is
recaptured the instant something unexpected happens \cite{smallwood2015science, christoff2016mind}. In Active
Inference this is a statement about precision: while prediction errors are small and the habitual policy is
reliable, attention --- the precision assigned to sensory prediction errors --- can be withdrawn from the task; a
surprise raises it again \cite{feldman2010attention}. The arbitration between habitual and deliberative control by
their relative uncertainty \cite{daw2005uncertainty} has the same structure.

The architecture already separates the two systems. The controller and the concept agents form an automatic layer
that drives, avoids obstacles and executes missions without the language model; the language agent is the
deliberative voice. We propose to make the allocation of the voice explicit through a scalar \emph{demand} signal,
computed on every digest cycle without the language model,
\begin{equation}
    D_t \;=\; \sum_{k} \pi_k\, \varepsilon_{k,t}^{2} \;+\; \lambda\, \mathbb{1}\!\left[\text{person active}_t\right],
    \label{eq:inner_demand}
\end{equation}
where $\varepsilon_{k,t}$ are normalised prediction errors reported by the agents (unexplained residual occupancy,
flips of existence or door beliefs, localisation uncertainty, stalled or failed mission steps, controller refusals)
and $\pi_k$ their precisions. Three modes follow from two thresholds with asymmetric hysteresis:
\begin{equation}
    m_t =
    \begin{cases}
        \textsc{focused}    & D_t > \theta_{\text{hi}},\\[2pt]
        \textsc{wandering}  & D_\tau < \theta_{\text{lo}} \;\; \forall\, \tau \in [t - T, t],\\[2pt]
        \textsc{monitoring} & \text{otherwise.}
    \end{cases}
    \label{eq:inner_modes}
\end{equation}
Entering \textsc{wandering} requires demand to stay low for a dwell time $T$; leaving it requires a single sample
above $\theta_{\text{hi}}$. The robot is slow to drift off and instant to snap back. Safety never depends on this
arbitration: the controller's reflexes act on their own, and the deliberative voice arrives after the reflex, as
awareness follows braking in the driver.

\paragraph{Focused mode} runs the question loop of Section~\ref{sec:inner_speech:questions} on the triggering
event. \paragraph{Monitoring mode} makes no language-model calls; the voice is primed but silent.

\paragraph{Wandering mode.}
Freed from the task, the voice is seeded with one recent experience chosen by salience: an episode that ended
unresolved, a failure, an object that changed, or an upcoming goal. Each step's output, together with an
associative retrieval from episodic memory, seeds the next, so that the chain drifts as spontaneous thought does
\cite{christoff2016mind}. The voice has access only to memory queries and to the architecture's imagination tools,
which simulate observations and outcomes on a copy of the belief graph; it can neither act nor write beliefs. Its
products are of three kinds: counterfactuals over past episodes (``had the door been open, the crossing would have
taken 40\,s less''), hypotheses to be tested later, and prospective plans (``if the door is shut on the way back,
the route through the next room is shorter''). Functionally this is the default network's role in simulating
past and future \cite{buckner2008brain}, and the offline replay and model reduction that Active Inference assigns
to rest and sleep \cite{friston2017curiosity, hobson2012waking}. Replay is prioritised by the expected gain of the
memory it revisits \cite{mattar2018prioritized}.

\paragraph{Source monitoring.}
Humans occasionally misremember imagined events as perceived ones \cite{johnson1993source}. The robot must not.
Everything produced while wandering is stored with an \emph{imagined} provenance, separate from the factual
episodic record, and enters the learning loop only as a hypothesis to be tested, never as evidence.

\subsection{What the return to reality gives the automatic layer}
\label{sec:inner_speech:return}

The transition from \textsc{wandering} to \textsc{focused} cancels the chain in flight (its state is kept, so that
it may resume if demand subsides) and gives the full voice to the surprise. It serves the automatic layer at three
timescales.
\begin{enumerate}
    \item \textbf{Decision.} A surprise means the automatic layer has left its domain of competence, where its only
    responses are to retry, stop or fail. The voice adds re-interpretation of the scene, arbitration between goals,
    re-planning through remembered topology, and recourse to a person.
    \item \textbf{Attention.} The voice can raise precision where it matters: orienting sensors, increasing a
    detector's rate in a region, or passing a hypothesis to the concept agents as a spatial prior. Such a prior has
    bounded weight; it may accelerate the recognition of an object that is present but may not create one that is
    not, preserving the grounding constraint.
    \item \textbf{Teaching.} Each resolved surprise is a labelled example the automatic layer could not have
    produced alone: a residual later identified as a chair is a false negative of the chair detector; a door that
    is repeatedly shut in the evening becomes a prior; a recovery that repeatedly succeeds becomes a branch of the
    controller's behaviour tree. The domain of competence grows, and the situations that require the voice become
    rarer, which is the classical progression from effortful to autonomous skill \cite{fitts1967human}.
\end{enumerate}
Plans prepared while wandering act as a cache: when the anticipated situation occurs, the focused voice verifies a
plan instead of composing one. The return ends with an explicit release --- spatial priors cleared, safety margins
restored, the automatic layer confirmed back within its competence --- without which every surprise would leave
the robot permanently vigilant.

\subsection{Relation to chain-of-thought reasoning}
\label{sec:inner_speech:cot}

Chain-of-thought is inner speech in the functional sense, but of a specific kind, and the comparison sharpens the
design (Table~\ref{tab:inner_cot}). Plain chain-of-thought is convergent and ungrounded: tokens condition on
tokens until an answer is produced, and the trace is discarded. Interleaving reasoning with action and observation
grounds it \cite{yao2023react}; our focused loop is of this kind, with the additional property that the
observation, through \eqref{eq:inner_update}, rather than the next token, determines where the reasoning goes.
The wandering voice is closer to plain chain-of-thought in its isolation from perception, but divergent rather
than convergent, interruptible, and persistent in its products.

\begin{table}[ht]
    \centering
    \small
    \begin{tabular}{|>{\raggedright\arraybackslash}p{0.19\linewidth}|>{\raggedright\arraybackslash}p{0.22\linewidth}|>{\raggedright\arraybackslash}p{0.24\linewidth}|>{\raggedright\arraybackslash}p{0.24\linewidth}|}
        \hline
                                & \textbf{Chain-of-thought} & \textbf{Focused voice} & \textbf{Wandering voice} \\ \hline
        Goal                    & answer a given question   & resolve a given surprise & none given; seed by salience \\ \hline
        Dynamics                & convergent                & convergent             & divergent, associative \\ \hline
        Contact with the world  & none                      & thought--action--observation & none (imagination only) \\ \hline
        Persistence             & discarded                 & question objects, beliefs, lessons & imagined store, hypotheses \\ \hline
        Interruptible           & no                        & yes                    & yes, resumable \\ \hline
        Latency budget          & n/a                       & $\sim$1\,s, no extended thinking & unbounded, extended thinking \\ \hline
    \end{tabular}
    \caption{Chain-of-thought compared with the two modes of the proposed inner voice.}
    \label{tab:inner_cot}
\end{table}

Two practical consequences follow. First, \emph{think when there is time, not when there is pressure}. In the
deployed language agent extended thinking is disabled, because a thinking-mode model spent 3\,000--6\,000 tokens
(40--80\,s) before each tool call; that is incompatible with the focused mode but costs nothing in the wandering
mode, where it can be enabled with a large budget and pre-empted at any moment. Slow reasoning is done in advance
and fast reasoning uses its results. Second, chain-of-thought traces are known not always to reflect the
computation that produced the answer \cite{turpin2023language}, as human verbal reports do not always reflect the
causes of behaviour \cite{nisbett1977telling}. The faithful record of the robot's reasoning is therefore the
structured one --- hypotheses, predictions, observations and posterior updates --- not the text. Finally, wandering
chains whose hypotheses were later confirmed by perception form a corpus of reasoning verified against the world,
which could be used to fine-tune the local model in the manner of self-taught reasoning \cite{zelikman2022star},
with perception, rather than a known answer, as the filter.

\subsection{Consequences for belief and learning}
\label{sec:inner_speech:learning}

The loop changes how beliefs evolve in two ways. Within an episode, uncertainty is reduced where surprise
occurred rather than where the sensors happened to point, which is efficient but leaves the rest of the map less
examined and makes the robot's own data a biased sample; a decaying novelty term \cite{schwartenbeck2019curiosity}
and importance weighting of learned statistics are required to correct this. Across episodes, every closed
question is a labelled example, and learning proceeds at four timescales (Table~\ref{tab:inner_learning}).

\begin{table}[ht]
    \centering
    \small
    \begin{tabular}{|>{\raggedright\arraybackslash}p{0.14\linewidth}|>{\raggedright\arraybackslash}p{0.37\linewidth}|>{\raggedright\arraybackslash}p{0.37\linewidth}|}
        \hline
        \textbf{Timescale} & \textbf{What changes} & \textbf{Where, in this architecture} \\ \hline
        Seconds   & posterior over the hypotheses of one question         & the question object \eqref{eq:inner_question} \\ \hline
        Days      & priors: where objects are, which doors are shut, what moves & long-term object and topological memory, as Dirichlet counts \\ \hline
        Weeks     & likelihoods: detector and concept models              & false negatives and positives labelled by resolved questions \\ \hline
        Slow      & structure: new categories                             & recurrent, mutually similar ``unknown'' resolutions named by the voice and proposed as a new concept \\ \hline
    \end{tabular}
    \caption{Learning driven by the inner-question loop.}
    \label{tab:inner_learning}
\end{table}

The last row is where the hypothesis of a language for thought is most directly engaged: a cluster of anomalies
that resists every existing concept becomes learnable once it is named. The loop also learns about itself: the
salience weights $\pi_k$ in \eqref{eq:inner_demand} are adjusted by the information each trigger actually yielded,
asking a person is weighed against looking, and loops that repeatedly resolve the same way are compiled into
reflexive rules, so that the number of language-model calls per hour should \emph{decrease} as competence grows.

\subsection{Why a language model: generation versus evaluation}
\label{sec:inner_speech:why_lm}

None of the above requires a person, so the question arises of what a language model contributes to an
architecture that already performs inference and action selection. Our answer is a division of labour. Expected
free energy evaluates policies and Bayesian updating weighs hypotheses, but both operate over a space that has
already been enumerated: each concept agent carries a fixed generative model, policies must be listed before they
can be scored, and Bayesian model reduction can prune structure but not invent it \cite{friston2017curiosity}.
Whatever falls outside the given models remains an unexplained residual. A language model is weak where this
machinery is strong --- its confidences are not calibrated posteriors and it has no access to the sensors --- and
strong where it is weak: it is a cheap, broad generator of hypotheses, explanations and plans. We therefore use it
only as a \emph{proposal distribution}, in a generate-and-test loop whose test is active inference:
\begin{itemize}
    \item \emph{Hypotheses outside the built-in models.} A residual blob can be proposed to be a bag, a toppled
    chair or a robot vacuum, none of which any concept agent represents; the question loop then tests them.
    \item \emph{Amortised policy search.} Instead of enumerating multi-step policies, the model writes a few
    candidate missions --- small programs in the mission language --- which the gate validates and expected free
    energy scores.
    \item \emph{Relational binding across agents.} The agents' representations (signed distance fields, occupancy
    grids, existence log-odds, episodic records) are mutually incompatible; a sentence such as ``the chair that was
    not next to the table in the room behind door~3 yesterday'' binds them in one query no single agent could pose.
    \item \emph{Explanation from few examples.} After a single surprise the model can propose which variable
    mattered (``the door is shut because it is evening''), turning one episode into a testable generalisation.
    \item \emph{Abstraction and naming.} Episodes are compressed into reusable lessons, and recurrent unknowns are
    named, which is what makes them candidates for a new concept (Table~\ref{tab:inner_learning}).
\end{itemize}
Language, rather than an arbitrary learned embedding, matters for three reasons: it is discrete and compositional,
so new hypotheses and plans are assembled from known parts; it is inspectable, so that deterministic gates can parse
and check it and people can audit it; and, once the robot thinks in it, communication with a person comes at no
further cost.

\paragraph{Where the knowledge comes from.}
This raises a difficulty for the motivating thesis. If the inner voice precedes the outer one, the common-sense
knowledge on which the proposer's value rests cannot have arrived through communication --- yet a language model
acquires all of it from human text, that is, from other people's outer voices. The difficulty dissolves once the
\emph{faculty} is distinguished from its \emph{content}. The internalist claim concerns the former: a
computational system for building unbounded hierarchical structures, available for thought before it is used for
talk \cite{hauser2002faculty, berwick2016why}. It makes no claim that what a person knows was acquired without
language. Human knowledge has three sources: core systems for objects, agents, number and space, present before
learning \cite{spelke2007core}; individual experience; and, once outer language exists, cumulative cultural
transmission, which supplies most of what an adult knows \cite{tomasello1999cultural}. Even on the internalist
view, the inner voice of a modern human is filled largely by the outer voices of others.

The architecture has the same three sources, in a different arrangement. The concept agents' hand-built generative
models play the role of \emph{core knowledge}: given before experience and limited to a few domains. The robot's
own counts $n_x(h)$ in \eqref{eq:inner_prior} are \emph{individual experience}. The language model is
\emph{culture}: knowledge distilled from a community's communication and inherited, for this individual, at birth.
The robot thus resembles a child born into a culture more than an early hominin, and its inner voice is Vygotskian
in origin: its format and its content were both learned from outer speech. Two consequences follow. First, the
proposal does not demonstrate the internalist thesis and should not be read as doing so; its claim is functional
--- an inner use of language has epistemic value independently of communication --- and holds whatever the origin of
the language system. Second, the format and the content of the voice can in principle be separated, which turns the
question into an experiment (P7 in Section~\ref{sec:inner_speech:predictions}): a compositional generator over the
robot's own concepts and mission language, trained only on the robot's experience, in the manner of library
learning in program induction \cite{ellis2021dreamcoder}, would supply the format without the culture.

\subsection{Mapping onto the existing architecture}
\label{sec:inner_speech:mapping}

Several components already exist in the language agent. Robot events (mission reports, failures, completions)
already wake the model through a queue, in turns restricted to perceiving, querying memory, stopping and
proposing, never acting. Every action already passes a deterministic gate. The decision to ask or to assume is
already decision-theoretic. Unexplained occupancy is already isolated by a residual agent that acts as the null
hypothesis of the scene, and episodic and long-term object memories are already joined read-only. Missing are the
persistent question objects \eqref{eq:inner_question}, a prediction evaluator implementing
\eqref{eq:inner_update} outside the language model, the demand signal and mode arbitration
\eqref{eq:inner_demand}--\eqref{eq:inner_modes}, the imagined store with its provenance, and the channel through
which bounded spatial priors reach the concept agents.

\subsection{Testable predictions}
\label{sec:inner_speech:predictions}

Consistent with the evaluation principle adopted throughout --- policies compared against each other inside one
closed loop, with body, sensors and environment fixed --- the proposal makes seven falsifiable predictions, each
tested in simulation by moving known objects and introducing unfamiliar ones:
\begin{enumerate}[label=\textbf{P\arabic*}]
    \item The unexplained residual mass per room decreases faster with the focused loop than without it.
    \item The proposer's calibration on resolved questions improves, and $\beta_x$ in \eqref{eq:inner_prior}
    decreases, as the robot's own counts accumulate.
    \item Detector false-negative rates fall after retraining on question-labelled examples.
    \item Inner questions per hour decrease over repeated exposure to the same environment while answer accuracy
    is maintained.
    \item Plans prepared while wandering reduce the latency from surprise to a validated response, and the time to
    notice a deliberately moved object does not grow despite habituation.
    \item With the language-model proposer replaced by a fixed hypothesis set (the existing concept classes and
    ``unknown'') and template-based missions, the loop resolves fewer questions about objects and failures outside
    the built-in models, and acquires new categories more slowly.
    \item A compositional generator trained only on the robot's own experience (format without culture) recovers
    part of the advantage in P6 within a familiar environment, and less of it the further a situation departs from
    what the robot has seen.
\end{enumerate}
A failure of P1 would indicate that the epistemic value of language-mediated questions does not exceed that of
the existing exploration; a failure of P4 would indicate that the loop does not transfer competence to the
automatic layer, which is its principal claim; and a failure of P6 would indicate that the language model adds
nothing a fixed hypothesis set does not, which would remove the case for including it. P7 separates the
contribution of language as a format for thought from that of the knowledge it carries.

\subsection{Limitations}
\label{sec:inner_speech:limitations}

The demand signal \eqref{eq:inner_demand} sums heterogeneous prediction errors whose precisions must themselves be
calibrated, and a poorly scaled term will either suppress the voice or keep it permanently engaged. A single
language model serves the person, the focused loop and the wandering loop; pre-emption keeps the person first but
bounds how much offline reasoning a busy robot can do. The tempered prior \eqref{eq:inner_prior} is a heuristic
combination rather than a derived posterior. The generator of P7 is substantially harder to build than the
ablation of P6, and may have to be approximated by a small model trained on the robot's logs and mission corpus,
whose pre-training would leak some culture back in. Finally, the analogy with human inner speech is functional, not
mechanistic: the proposal is that the same computational role --- serial, compositional self-questioning in the
service of epistemic value --- is useful for a robot, not that the robot's voice resembles the human one in any
further respect.

%     % =====================================================================================================
% language_hierarchy_section.tex — Language over the belief hierarchy: culture, sequencing and imagination
%
% DRAFT (2026-11-19). A proposal section: the architecture it describes is NOT yet implemented beyond the parts
% stated as existing. Source design note: docs/LANGUAGE_OVER_BELIEF_HIERARCHY.md.
% Include AFTER inner_speech_section.tex (it refers to that section's equations and modes):
%     \input{inner_speech_section}
%     \input{language_hierarchy_section}
% Bibliography: inner_speech.bib (shared). Packages: amsmath, amssymb, enumitem, array (as for the inner-speech
% section). friston2017deeptemporal duplicates main.tex's friston2017deep under a new key on purpose, so that this
% file builds with inner_speech.bib alone; merge the two keys when the bibliographies are merged.
% =====================================================================================================

\section{Language over the Belief Hierarchy}
\label{sec:lang_hierarchy}

Section~\ref{sec:inner_speech} treated the language model as an inner voice and asked what language contributes to
thought. Here we take the complementary, externalist view. The vision-language model (VLM) arrives with a large share
of human common sense already in it; for the robot it is \emph{culture}, an external store of knowledge that extends
its mind in the sense of \cite{clark1998extended} and that it inherits rather than acquires
(Section~\ref{sec:inner_speech:why_lm}). The question is how to couple such a store to a belief state that is built
by probabilistic inference, so that the robot gains from the knowledge without inheriting the store's lack of
grounding and calibration.

Our answer is to bind the VLM to the robot's \emph{belief hierarchy} in three roles --- culture, sequencing and
imagination --- each under one rule: \textbf{the language model proposes, the generative models dispose}. The roles
are arbitrated by the attention modes of Section~\ref{sec:inner_speech:modes}: the VLM imagines while the automatic
level runs well, and monitors and replans when it does not.

\subsection{The automatic belief hierarchy}
\label{sec:lang_hierarchy:automatic}

The robot builds its belief state without language, and the build is hierarchical (Table~\ref{tab:lh_levels}).
Concept agents fit level-1 objects from masks and point clouds, each holding a metric belief (mean and covariance)
and an existence log-odds updated by evidence that was predicted and observed, or predicted and absent where
detection was possible. Level-2 \emph{metaconcepts} fit arrangements over those beliefs: a ring of chairs around a
table, or the rectilinear frame --- one grid axis, one worktop plane, one carcass depth --- shared by the units of a
kitchen. A metaconcept's measurements are its members' published posteriors, its existence is a log Bayes factor of
``generated by the arrangement'' against an explicit ``independent objects'' null, and it sends empirical priors
back to its members. The message to member $i$ is computed from the arrangement fitted without $i$, the cavity of
expectation propagation \cite{minka2001expectation},
\begin{equation}
    q_{\setminus i}(\theta) \;\propto\; \frac{q(\theta)}{m_i(\theta)},
    \label{eq:lh_cavity}
\end{equation}
so that a member is never told to agree with itself. An exploration drive moves the robot so that this hierarchy
fills in: explore a room until continuing no longer pays, cross the nearest door not yet crossed, return depth-first.
Empty slots of an arrangement seed births and next-best views. In the terms of the driving example, this is the
expert on a familiar road: anticipation and confirmation are automatic, and geometric.

\begin{table}[ht]
    \centering
    \small
    \begin{tabular}{|>{\raggedright\arraybackslash}p{0.07\linewidth}|>{\raggedright\arraybackslash}p{0.28\linewidth}|>{\raggedright\arraybackslash}p{0.27\linewidth}|>{\raggedright\arraybackslash}p{0.28\linewidth}|}
        \hline
        \textbf{Level} & \textbf{Latent} & \textbf{Source of its model} & \textbf{Status} \\ \hline
        0 & robot body and sensor mounts & calibration & implemented \\ \hline
        1 & room geometry; objects (pose, extent, existence) & authored generative models, fitted online & implemented \\ \hline
        2 & arrangements (ring, rectilinear frame) & authored schemas & implemented (read-only steering for the kitchen) \\ \hline
        3 & function of a room (kitchen, dining, bedroom, \ldots) & \textbf{culture: the VLM}, calibrated against counts & proposed (Section~\ref{sec:lang_hierarchy:culture}) \\ \hline
    \end{tabular}
    \caption{The belief hierarchy. Geometry is authored up to level 2; culture enters at level 3.}
    \label{tab:lh_levels}
\end{table}

What the hierarchy cannot do is the subject of the rest of the section. It cannot say what a room is \emph{for},
because that is not a geometric fact. It cannot pursue a goal it was not built to pursue. And it never looks ahead:
its predictions concern the present scene, not the outcome of what the robot is about to do.

\subsection{Three roles of the language model}
\label{sec:lang_hierarchy:roles}

\begin{table}[ht]
    \centering
    \small
    \begin{tabular}{|>{\raggedright\arraybackslash}p{0.14\linewidth}|>{\raggedright\arraybackslash}p{0.43\linewidth}|>{\raggedright\arraybackslash}p{0.33\linewidth}|}
        \hline
        \textbf{Role} & \textbf{What the VLM supplies} & \textbf{What decides} \\ \hline
        Culture     & priors and structure the hierarchy cannot author: room function, class co-occurrence, arrangement schemas, norms & evidence against a null (the chassis rules) \\ \hline
        Sequencing  & goals decomposed into missions over the hierarchy's beliefs, under contextual constraints, with the belief change each step should produce & a mission checker and command gate, the controller, then the belief dynamics \\ \hline
        Imagination & \emph{which} subjunctive questions to put to the robot's own models while the automatic level runs well & the models, run on an imagined copy of the world \\ \hline
    \end{tabular}
    \caption{The three roles. In each, the language model proposes and a generative model disposes.}
    \label{tab:lh_roles}
\end{table}

The three roles share one asymmetry, already argued in Section~\ref{sec:inner_speech:why_lm}: inference evaluates
hypotheses and policies well but cannot generate new ones, while a language model generates them cheaply but cannot
evaluate them. Each role pairs the two in a generate-and-test loop whose test is part of the robot's generative model.

\subsection{Culture: a level-3 function hypothesis}
\label{sec:lang_hierarchy:culture}

\paragraph{Latent and evidence.}
For each room $r$ and function $f$ we maintain an independent hypothesis ``$r$ serves $f$'' against the null ``$r$
is a generic space''. Independent rather than categorical, because open-plan rooms serve several functions, and
because this is how the metaconcept chassis already treats its instances. Members are level-1 objects of class $c$
with existence probability $p_c$, and level-2 arrangements. Presence, not count, is the observation, so that ten
chairs do not argue ten times. With $\pi_{cf} = P(c \text{ present} \mid f)$, $\pi_{c0}$ its value under the null, and
$d_c$ the probability that $c$ would have been detected in $r$ had it been present, the evidence for $f$ is
\begin{equation}
    \Lambda_f(r) \;=\; \sum_c \Big[\, p_c \log\frac{\pi_{cf}}{\pi_{c0}}
    \;+\; (1 - p_c)\log\frac{1 - \pi_{cf}\, d_c}{1 - \pi_{c0}\, d_c} \Big].
    \label{eq:lh_function_evidence}
\end{equation}
The second term is expected-but-absent evidence, and it is only as strong as the search was: with $d_c = 0$ it
vanishes. This is the negative-information rule the concept agents apply to instance death, lifted to the scene.
$d_c$ is the room's exploration completeness times a per-class detection rate, to be replaced by a calibrated
detectability model. The weighting by $p_c$ is an approximation to the expected log-likelihood ratio, exact when
presence is certain. $\Lambda_f(r)$ is the hypothesis's existence log-odds; birth and death follow the chassis, and
no threshold is introduced.

\paragraph{Priors from culture, replaced by experience.}
The conditional probabilities $\pi_{cf}$ are what no geometric model can supply. We elicit them from the VLM over a
closed vocabulary --- the graph's own class names and a fixed set of functions --- and give each a Beta prior whose
strength is earned:
\begin{equation}
    \pi_{cf} \sim \mathrm{Beta}\big(\kappa\beta_f\,\pi^{\text{VLM}}_{cf},\; \kappa\beta_f\,(1-\pi^{\text{VLM}}_{cf})\big),
    \qquad
    \hat\pi_{cf} = \frac{\kappa\beta_f\,\pi^{\text{VLM}}_{cf} + n_{cf}}{\kappa\beta_f + n_f}.
    \label{eq:lh_beta_prior}
\end{equation}
Here $\kappa$ is a small pseudo-count (the VLM is worth a few rooms, not a hundred), $\beta_f \in [0,1]$ is the
VLM's measured calibration for $f$, and $n_f$, $n_{cf}$ count rooms confirmed as $f$ --- by a person, after their
explicit agreement, or by a saturated posterior over repeated visits --- of which $n_{cf}$ contained $c$. This is the
conjugate form of the tempered prior \eqref{eq:inner_prior}: culture dominates early and experience takes over. The
same elicitation replaces priors that are hand-authored today; the kitchen metaconcept's class-membership log-odds
are documented in its configuration as ``elicited, not measured'' and awaiting a corpus.

\paragraph{Messages down, and the cavity rule for culture.}
A function hypothesis first sends nothing into member beliefs, then exploration targets only: expected but unseen
classes with a coarse placement, consumed by the exploration drive and the next-best-view planners as the value of a
viewpoint, never as evidence of existence. This introduces selection bias but not contamination, and
\eqref{eq:lh_function_evidence} stays honest because absence counts in exactly the places that were searched. Only
later, if earned, would a function seed level-2 births, and then the down-date \eqref{eq:lh_cavity} applies: the
cavity rule binds the language model as it binds every other layer, so that ``kitchen'' cannot confirm itself through
the members it caused to be born.

\paragraph{Norms.}
Each function also carries the contextual norms the VLM supplies for it --- quiet and slow in a bedroom at night,
clear of the work area in a kitchen in use, never blocking a doorway. They are not evidence; they are the
preferences of the generative model, and Section~\ref{sec:lang_hierarchy:sequencing} uses them as planning
constraints.

\paragraph{Model-level surprise.}
Levels 1--2 already produce the signals that their \emph{model} is wrong rather than uncertain: an arrangement whose
Bayes factor collapses, a member with a large cavity residual, a class that persistently scores near-zero membership,
two schemas both weak. These are the triggers for focused-mode diagnosis. An instance from the live system: a
refrigerator scored kitchen membership $0.00$ because only base and wall units were modelled, so its $1.85$\,m top
was scored against a $0.90$\,m worktop plane. A person diagnosed the structural defect and a third, tall tier fixed
it. Proposing such revisions --- of the model, not its parameters --- is the VLM's task in focused mode; accepting
them remains, for now, a person's.

\subsection{Sequencing over the hierarchy}
\label{sec:lang_hierarchy:sequencing}

\paragraph{Goals at their own level.}
Mission predicates in the deployed system already refer to beliefs (a room explored with probability $\geq 0.9$, a
door open with probability $< 0.5$, an object class seen). The hierarchy extends the vocabulary in which goals are
stated and decomposed, which is the deep temporal structure of Active Inference \cite{friston2017deeptemporal}: a goal
is planned at its own level and refined downward (Table~\ref{tab:lh_goals}).

\begin{table}[ht]
    \centering
    \small
    \begin{tabular}{|>{\raggedright\arraybackslash}p{0.2\linewidth}|>{\raggedright\arraybackslash}p{0.27\linewidth}|>{\raggedright\arraybackslash}p{0.43\linewidth}|}
        \hline
        \textbf{Goal stated at} & \textbf{Example} & \textbf{Decomposition} \\ \hline
        Level 3, function    & ``go where people eat''           & room with the highest $\Lambda_{\text{dining}}$ $\to$ its dining set $\to$ a standpoint \\ \hline
        Level 2, arrangement & ``check every chair at the table'' & the ring's members \emph{and its empty slots}, predicted but not yet seen \\ \hline
        Level 1, object      & ``approach the fridge''           & an approach to the object \\ \hline
        Level 0, pose        & ``turn 90 degrees''               & a relative motion \\ \hline
    \end{tabular}
    \caption{Goals are planned at their own level of the hierarchy and refined downward.}
    \label{tab:lh_goals}
\end{table}

\paragraph{Epistemic sequences.}
The exploration drive is the robot's compiled epistemic plan, and the right default for building the map. A specific
goal makes it inefficient. To find a mug in a partly known house, the VLM writes the sequence the drive cannot: go to
the room with the highest kitchen evidence (level 3), look along the worktop plane where mugs rest (level 2), approach
once one is seen. The hierarchy says where the information is; the VLM says in what order, and when to stop.

\paragraph{Every step predicts its belief outcome.}
A mission step currently reports completion or failure. We require it also to declare the belief change it should
produce: a crossing predicts that the current room flips and is recognised; an exploration predicts that completeness
rises and that some $\Lambda_f$ leaves zero; an approach predicts that the target's existence stays high as range
closes. Let $\phi_k$ be the predicate step $k$ predicts over beliefs, due by the end of the step, and $\alpha_j$ the
assumptions the running plan relies on (``door 3 is open'', ``room 2 is still the kitchen''). Both enter the demand
signal \eqref{eq:inner_demand} as surprisal terms,
\begin{equation}
    \varepsilon^{\text{plan}}_t \;=\; -\log P_t(\phi_k) \;\;\text{at the step's deadline},
    \qquad
    \varepsilon^{\text{assume}}_{j,t} \;=\; -\log P_t(\alpha_j) \;\;\text{continuously},
    \label{eq:lh_plan_surprise}
\end{equation}
with $P_t$ the current belief. The first catches a step whose outcome diverged even though it reported success; the
second raises attention when an assumption weakens, before any step fails --- the driver who slows at a school sign
rather than braking for the child. In either case the focused voice replans from the belief state the robot actually
has. A plan is thereby the same kind of object as an inner question (Section~\ref{sec:inner_speech:questions}): a
proposal whose predictions the belief dynamics confirm or refute.

\paragraph{Constraints from context.}
The norms attached to the active function hypotheses become mission-wide invariants --- avoid a region, cap the speed,
keep a distance --- monitored for violation throughout the mission rather than checked once at its start.

\paragraph{Habits.}
Sequences that recur and succeed are compiled. In Active Inference the policy prior combines a habit term $E$ with
expected free energy $G$ under a precision $\gamma$ \cite{fitzgerald2014model, friston2016learning},
\begin{equation}
    p(\pi) \;=\; \sigma\big(\ln E(\pi) - \gamma\, G(\pi)\big),
    \label{eq:lh_habit}
\end{equation}
where $E$ is learned by counting the policies the robot executes successfully. The VLM contributes to $G$: it proposes
the candidate policies and the preferences they are scored against. With practice, routine sequences migrate into $E$
and stop requiring the VLM, as a learner driver's verbal sequence becomes an expert's automatic one. Compiled habits
are naturally indexed by the hierarchy (``check a dining set'', ``survey a worktop'') --- plan schemas over arrangement
schemas --- and carry their assumptions $\alpha_j$, so that they know when they no longer apply.

\subsection{Imagination while the automatic level runs well}
\label{sec:lang_hierarchy:imagination}

\paragraph{Principle.}
Imagining is running the robot's own generative models offline on an imagined copy of the world. Both exist in the
deployed language agent, as tools a person can query: a copy of the room's facts that interventions edit and that
cannot write the live graph; the controller's own fit model (``would I fit through door 1 if it were half open?'');
the perception models (``would I see the fridge from there?''); concept priors sampled and rendered through the
camera and LiDAR models with a Monte-Carlo probability of detection (``how do you imagine a chair?''); and success
odds for a route from the day's recorded crossings. In wandering mode the robot puts these questions to itself. The
VLM chooses what to imagine --- the seed, the intervention, the question --- and the models compute the answer. A VLM
answering a subjunctive question by itself would be confabulating, not imagining.

\paragraph{Five kinds of imagining.}
Table~\ref{tab:lh_imagination} lists them by what they prepare. The first two serve the near future, the next two the
recent past, the last the model itself.

\begin{table}[ht]
    \centering
    \small
    \begin{tabular}{|>{\raggedright\arraybackslash}p{0.16\linewidth}|>{\raggedright\arraybackslash}p{0.21\linewidth}|>{\raggedright\arraybackslash}p{0.27\linewidth}|>{\raggedright\arraybackslash}p{0.25\linewidth}|}
        \hline
        \textbf{Kind} & \textbf{Seed} & \textbf{Question to the models} & \textbf{Product} \\ \hline
        Rehearsal & next steps of the running mission & success odds per leg; visibility from the next standpoint; fit through a partly open door & a contingency cache keyed by the assumptions that would trigger each alternative \\ \hline
        Scene completion & empty slots; expected but unseen classes & sample the class at its predicted place; probability of detection from candidate viewpoints & viewpoints with their detection probability \\ \hline
        Counterfactual replay & a failed or slow episode & the same episode through another door, or with another opening & lessons as testable hypotheses; assumptions for habits \\ \hline
        Anomaly pre-analysis & an unexplained residual; a question closed as ``unknown'' & candidate classes rendered at that place against the detection record & ranked hypotheses for the next focused episode \\ \hline
        Structural replay & logged metaconcept inputs & refit under an alternative schema or tier & structural proposals for review \\ \hline
    \end{tabular}
    \caption{Imagination in wandering mode. The VLM chooses the seed and the question; the robot's models answer.}
    \label{tab:lh_imagination}
\end{table}

Structural replay is the kind of process that could have found the tall tier: refit the logged kitchen memberships
with one more tier and observe the evidence rise.

\paragraph{What to imagine first.}
Imagination competes for a bounded budget and is pre-empted at any moment, so seeds are ranked by the value of their
product for upcoming decisions. Following prioritised replay \cite{mattar2018prioritized},
\begin{equation}
    V(s) \;=\; \mathrm{need}(s)\,\cdot\,\mathrm{gain}(s),
    \label{eq:lh_priority}
\end{equation}
where need is how soon the situation $s$ prepares for will occur --- the next legs of the running mission first, then
the rooms the drive will enter next, then open questions, then old failures --- and gain is how much the product could
change the decision taken there: legs with uncertain odds, expectations with detection probability near one half,
hypotheses that are currently tied.

\paragraph{Products and their weight.}
A contingency cache is attached to the running mission, so that on a snap-back the focused voice verifies a prepared
alternative instead of composing one. Viewpoints go to the exploration drive with bounded weight, through the same
channel as the function hypothesis's targets. Ranked hypotheses become priors of the next focused episode. Lessons
become candidate habits, trusted only after later executions confirm them. Every imagined episode is stored with an
imagined provenance, never in the live graph and never in factual episodic memory (Section~\ref{sec:inner_speech:modes}).
Every product that predicts something is scored when reality arrives, so that imagination earns trust per kind
exactly as the VLM's priors do; a kind whose predictions keep failing loses weight.

\paragraph{Gating.}
Imagination runs only in wandering mode, is cancelled the moment demand crosses $\theta_{\text{hi}}$, and resumes if
its seed is still relevant when demand subsides. Extended VLM reasoning is allowed here and nowhere else
(Section~\ref{sec:inner_speech:cot}).

\subsection{The three roles in the three modes}
\label{sec:lang_hierarchy:modes}

\begin{table}[ht]
    \centering
    \small
    \begin{tabular}{|>{\raggedright\arraybackslash}p{0.13\linewidth}|>{\raggedright\arraybackslash}p{0.23\linewidth}|>{\raggedright\arraybackslash}p{0.16\linewidth}|>{\raggedright\arraybackslash}p{0.2\linewidth}|>{\raggedright\arraybackslash}p{0.16\linewidth}|}
        \hline
        \textbf{Mode} & \textbf{Entered when} & \textbf{Culture} & \textbf{Sequencing} & \textbf{Imagination} \\ \hline
        Focused    & model-level surprise; step divergence or weakened assumption \eqref{eq:lh_plan_surprise}; the person speaks & diagnose the model; propose a revision & replan from the actual belief state, contingency cache first & --- \\ \hline
        Monitoring & between thresholds & --- & expected outcomes and assumptions checked; no VLM calls & --- \\ \hline
        Wandering  & demand low for a dwell time & offline elicitation; schema proposals & drafts for goals ahead & the five kinds, ranked by \eqref{eq:lh_priority} \\ \hline
    \end{tabular}
    \caption{Arbitration of the three roles by the attention modes of Section~\ref{sec:inner_speech:modes}.}
    \label{tab:lh_modes}
\end{table}

The inputs to the demand signal are therefore the signals this section introduces: model-level surprise from levels
1--3 and the plan surprisals \eqref{eq:lh_plan_surprise}. The loop closes across the roles: imagination prepares what
focus will need, focus produces the episodes imagination later replays, and both supply the counts with which the
robot's experience replaces the VLM's culture in \eqref{eq:lh_beta_prior} and its habits replace the VLM's plans in
\eqref{eq:lh_habit}. The VLM's share of the robot's competence should therefore shrink with experience for routine
matters, and remain for the novel, the social and the structural.

\subsection{Relation to prior work}
\label{sec:lang_hierarchy:related}

Language models have been used to sequence robot skills, scored by learned affordances \cite{ahn2022saycan}; to
re-plan from textual feedback on success, scenes and human input \cite{huang2022inner}; to plan over large
three-dimensional scene graphs by searching a collapsed graph and verifying with a simulator \cite{rana2023sayplan};
and to bias object search with soft common-sense constraints \cite{zhou2023esc}. The present proposal differs in three
respects. The language model is coupled to a \emph{probabilistic} hierarchy in which its contributions enter as
priors with measured calibration and are admitted only by evidence against a null, rather than to a symbolic or
deterministic representation. Its plans are predictions over beliefs whose divergence is a surprise signal, rather
than scripts checked for success. And its three roles are arbitrated by an attention mechanism that leaves it idle,
imagining, when the automatic level suffices.

\subsection{Testable predictions}
\label{sec:lang_hierarchy:predictions}

Each is tested policy against policy inside one closed loop, in simulation, with body, sensors and environment fixed.
\begin{enumerate}[label=\textbf{L\arabic*}, start=0]
    \item VLM-elicited class priors agree in sign with the hand-authored kitchen membership priors, and substituting
    them changes no correct live membership decision.
    \item A read-only function hypothesis recovers the functions of the simulated apartment's rooms; each
    $\Lambda_f$ moves away from zero as exploration completeness rises; removing the refrigerator weakens the kitchen
    hypothesis without destroying it.
    \item With predicates and goals at levels 2--3, goal-directed search (``find a mug'') takes less time and distance
    than the exploration drive alone, and fewer objects remain unseen in rooms counted as explored.
    \item With per-step expected outcomes and assumptions in the demand signal, more than half of the seeded
    divergences (a door shut mid-mission, an object moved) are caught before any step fails; without them, none is.
    \item Rehearsal and scene completion reduce the latency from surprise to a validated new plan, and the number of
    views needed to confirm an expected object.
    \item Counterfactual and structural replay propose the missing tall tier when the two-tier kitchen model is
    re-created.
    \item As confirmed rooms and successful missions accumulate, $\beta_f$ improves, seeded and unseeded posteriors agree
    on held-out evidence, and VLM calls per hour fall while task success holds.
\end{enumerate}
A failure of L1 ends the cultural layer before any coupling; a failure of L4 means imagination does not pay for its
compute, and the robot should simply idle when the automatic level suffices.

\subsection{Limitations}
\label{sec:lang_hierarchy:limitations}

The weighting by existence probability in \eqref{eq:lh_function_evidence} approximates the expected log-likelihood
ratio. The pseudo-count $\kappa$ and the calibration $\beta_f$ trade early usefulness against overconfidence, and
the VLM's notion of a typical home carries a cultural bias that only the robot's counts can correct; log-odds
down-weight unexpected classes but never exclude them. Per-step expected outcomes require each mission operation to
know its own postcondition over beliefs, which is straightforward for crossing and approach and less so for
open-ended exploration. One VLM serves the person, focused replanning and imagination, so strict pre-emption bounds
how much a busy robot can imagine. Finally, structural revisions are proposed but applied by a person; letting the
robot modify its own generative models is outside the scope of this proposal.

% Bibliography: inner_speech.bib (shared). Packages: amsmath, amssymb, enumitem, array (as for the inner-speech
% section). friston2017deeptemporal duplicates main.tex's friston2017deep under a new key on purpose, so that this
% file builds with inner_speech.bib alone; merge the two keys when the bibliographies are merged.
% =====================================================================================================

\section{Language over the Belief Hierarchy}
\label{sec:lang_hierarchy}

Section~\ref{sec:inner_speech} treated the language model as an inner voice and asked what language contributes to
thought. Here we take the complementary, externalist view. The vision-language model (VLM) arrives with a large share
of human common sense already in it; for the robot it is \emph{culture}, an external store of knowledge that extends
its mind in the sense of \cite{clark1998extended} and that it inherits rather than acquires
(Section~\ref{sec:inner_speech:why_lm}). The question is how to couple such a store to a belief state that is built
by probabilistic inference, so that the robot gains from the knowledge without inheriting the store's lack of
grounding and calibration.

Our answer is to bind the VLM to the robot's \emph{belief hierarchy} in three roles --- culture, sequencing and
imagination --- each under one rule: \textbf{the language model proposes, the generative models dispose}. The roles
are arbitrated by the attention modes of Section~\ref{sec:inner_speech:modes}: the VLM imagines while the automatic
level runs well, and monitors and replans when it does not.

\subsection{The automatic belief hierarchy}
\label{sec:lang_hierarchy:automatic}

The robot builds its belief state without language, and the build is hierarchical (Table~\ref{tab:lh_levels}).
Concept agents fit level-1 objects from masks and point clouds, each holding a metric belief (mean and covariance)
and an existence log-odds updated by evidence that was predicted and observed, or predicted and absent where
detection was possible. Level-2 \emph{metaconcepts} fit arrangements over those beliefs: a ring of chairs around a
table, or the rectilinear frame --- one grid axis, one worktop plane, one carcass depth --- shared by the units of a
kitchen. A metaconcept's measurements are its members' published posteriors, its existence is a log Bayes factor of
``generated by the arrangement'' against an explicit ``independent objects'' null, and it sends empirical priors
back to its members. The message to member $i$ is computed from the arrangement fitted without $i$, the cavity of
expectation propagation \cite{minka2001expectation},
\begin{equation}
    q_{\setminus i}(\theta) \;\propto\; \frac{q(\theta)}{m_i(\theta)},
    \label{eq:lh_cavity}
\end{equation}
so that a member is never told to agree with itself. An exploration drive moves the robot so that this hierarchy
fills in: explore a room until continuing no longer pays, cross the nearest door not yet crossed, return depth-first.
Empty slots of an arrangement seed births and next-best views. In the terms of the driving example, this is the
expert on a familiar road: anticipation and confirmation are automatic, and geometric.

\begin{table}[ht]
    \centering
    \small
    \begin{tabular}{|>{\raggedright\arraybackslash}p{0.07\linewidth}|>{\raggedright\arraybackslash}p{0.28\linewidth}|>{\raggedright\arraybackslash}p{0.27\linewidth}|>{\raggedright\arraybackslash}p{0.28\linewidth}|}
        \hline
        \textbf{Level} & \textbf{Latent} & \textbf{Source of its model} & \textbf{Status} \\ \hline
        0 & robot body and sensor mounts & calibration & implemented \\ \hline
        1 & room geometry; objects (pose, extent, existence) & authored generative models, fitted online & implemented \\ \hline
        2 & arrangements (ring, rectilinear frame) & authored schemas & implemented (read-only steering for the kitchen) \\ \hline
        3 & function of a room (kitchen, dining, bedroom, \ldots) & \textbf{culture: the VLM}, calibrated against counts & proposed (Section~\ref{sec:lang_hierarchy:culture}) \\ \hline
    \end{tabular}
    \caption{The belief hierarchy. Geometry is authored up to level 2; culture enters at level 3.}
    \label{tab:lh_levels}
\end{table}

What the hierarchy cannot do is the subject of the rest of the section. It cannot say what a room is \emph{for},
because that is not a geometric fact. It cannot pursue a goal it was not built to pursue. And it never looks ahead:
its predictions concern the present scene, not the outcome of what the robot is about to do.

\subsection{Three roles of the language model}
\label{sec:lang_hierarchy:roles}

\begin{table}[ht]
    \centering
    \small
    \begin{tabular}{|>{\raggedright\arraybackslash}p{0.14\linewidth}|>{\raggedright\arraybackslash}p{0.43\linewidth}|>{\raggedright\arraybackslash}p{0.33\linewidth}|}
        \hline
        \textbf{Role} & \textbf{What the VLM supplies} & \textbf{What decides} \\ \hline
        Culture     & priors and structure the hierarchy cannot author: room function, class co-occurrence, arrangement schemas, norms & evidence against a null (the chassis rules) \\ \hline
        Sequencing  & goals decomposed into missions over the hierarchy's beliefs, under contextual constraints, with the belief change each step should produce & a mission checker and command gate, the controller, then the belief dynamics \\ \hline
        Imagination & \emph{which} subjunctive questions to put to the robot's own models while the automatic level runs well & the models, run on an imagined copy of the world \\ \hline
    \end{tabular}
    \caption{The three roles. In each, the language model proposes and a generative model disposes.}
    \label{tab:lh_roles}
\end{table}

The three roles share one asymmetry, already argued in Section~\ref{sec:inner_speech:why_lm}: inference evaluates
hypotheses and policies well but cannot generate new ones, while a language model generates them cheaply but cannot
evaluate them. Each role pairs the two in a generate-and-test loop whose test is part of the robot's generative model.

\subsection{Culture: a level-3 function hypothesis}
\label{sec:lang_hierarchy:culture}

\paragraph{Latent and evidence.}
For each room $r$ and function $f$ we maintain an independent hypothesis ``$r$ serves $f$'' against the null ``$r$
is a generic space''. Independent rather than categorical, because open-plan rooms serve several functions, and
because this is how the metaconcept chassis already treats its instances. Members are level-1 objects of class $c$
with existence probability $p_c$, and level-2 arrangements. Presence, not count, is the observation, so that ten
chairs do not argue ten times. With $\pi_{cf} = P(c \text{ present} \mid f)$, $\pi_{c0}$ its value under the null, and
$d_c$ the probability that $c$ would have been detected in $r$ had it been present, the evidence for $f$ is
\begin{equation}
    \Lambda_f(r) \;=\; \sum_c \Big[\, p_c \log\frac{\pi_{cf}}{\pi_{c0}}
    \;+\; (1 - p_c)\log\frac{1 - \pi_{cf}\, d_c}{1 - \pi_{c0}\, d_c} \Big].
    \label{eq:lh_function_evidence}
\end{equation}
The second term is expected-but-absent evidence, and it is only as strong as the search was: with $d_c = 0$ it
vanishes. This is the negative-information rule the concept agents apply to instance death, lifted to the scene.
$d_c$ is the room's exploration completeness times a per-class detection rate, to be replaced by a calibrated
detectability model. The weighting by $p_c$ is an approximation to the expected log-likelihood ratio, exact when
presence is certain. $\Lambda_f(r)$ is the hypothesis's existence log-odds; birth and death follow the chassis, and
no threshold is introduced.

\paragraph{Priors from culture, replaced by experience.}
The conditional probabilities $\pi_{cf}$ are what no geometric model can supply. We elicit them from the VLM over a
closed vocabulary --- the graph's own class names and a fixed set of functions --- and give each a Beta prior whose
strength is earned:
\begin{equation}
    \pi_{cf} \sim \mathrm{Beta}\big(\kappa\beta_f\,\pi^{\text{VLM}}_{cf},\; \kappa\beta_f\,(1-\pi^{\text{VLM}}_{cf})\big),
    \qquad
    \hat\pi_{cf} = \frac{\kappa\beta_f\,\pi^{\text{VLM}}_{cf} + n_{cf}}{\kappa\beta_f + n_f}.
    \label{eq:lh_beta_prior}
\end{equation}
Here $\kappa$ is a small pseudo-count (the VLM is worth a few rooms, not a hundred), $\beta_f \in [0,1]$ is the
VLM's measured calibration for $f$, and $n_f$, $n_{cf}$ count rooms confirmed as $f$ --- by a person, after their
explicit agreement, or by a saturated posterior over repeated visits --- of which $n_{cf}$ contained $c$. This is the
conjugate form of the tempered prior \eqref{eq:inner_prior}: culture dominates early and experience takes over. The
same elicitation replaces priors that are hand-authored today; the kitchen metaconcept's class-membership log-odds
are documented in its configuration as ``elicited, not measured'' and awaiting a corpus.

\paragraph{Messages down, and the cavity rule for culture.}
A function hypothesis first sends nothing into member beliefs, then exploration targets only: expected but unseen
classes with a coarse placement, consumed by the exploration drive and the next-best-view planners as the value of a
viewpoint, never as evidence of existence. This introduces selection bias but not contamination, and
\eqref{eq:lh_function_evidence} stays honest because absence counts in exactly the places that were searched. Only
later, if earned, would a function seed level-2 births, and then the down-date \eqref{eq:lh_cavity} applies: the
cavity rule binds the language model as it binds every other layer, so that ``kitchen'' cannot confirm itself through
the members it caused to be born.

\paragraph{Norms.}
Each function also carries the contextual norms the VLM supplies for it --- quiet and slow in a bedroom at night,
clear of the work area in a kitchen in use, never blocking a doorway. They are not evidence; they are the
preferences of the generative model, and Section~\ref{sec:lang_hierarchy:sequencing} uses them as planning
constraints.

\paragraph{Model-level surprise.}
Levels 1--2 already produce the signals that their \emph{model} is wrong rather than uncertain: an arrangement whose
Bayes factor collapses, a member with a large cavity residual, a class that persistently scores near-zero membership,
two schemas both weak. These are the triggers for focused-mode diagnosis. An instance from the live system: a
refrigerator scored kitchen membership $0.00$ because only base and wall units were modelled, so its $1.85$\,m top
was scored against a $0.90$\,m worktop plane. A person diagnosed the structural defect and a third, tall tier fixed
it. Proposing such revisions --- of the model, not its parameters --- is the VLM's task in focused mode; accepting
them remains, for now, a person's.

\subsection{Sequencing over the hierarchy}
\label{sec:lang_hierarchy:sequencing}

\paragraph{Goals at their own level.}
Mission predicates in the deployed system already refer to beliefs (a room explored with probability $\geq 0.9$, a
door open with probability $< 0.5$, an object class seen). The hierarchy extends the vocabulary in which goals are
stated and decomposed, which is the deep temporal structure of Active Inference \cite{friston2017deeptemporal}: a goal
is planned at its own level and refined downward (Table~\ref{tab:lh_goals}).

\begin{table}[ht]
    \centering
    \small
    \begin{tabular}{|>{\raggedright\arraybackslash}p{0.2\linewidth}|>{\raggedright\arraybackslash}p{0.27\linewidth}|>{\raggedright\arraybackslash}p{0.43\linewidth}|}
        \hline
        \textbf{Goal stated at} & \textbf{Example} & \textbf{Decomposition} \\ \hline
        Level 3, function    & ``go where people eat''           & room with the highest $\Lambda_{\text{dining}}$ $\to$ its dining set $\to$ a standpoint \\ \hline
        Level 2, arrangement & ``check every chair at the table'' & the ring's members \emph{and its empty slots}, predicted but not yet seen \\ \hline
        Level 1, object      & ``approach the fridge''           & an approach to the object \\ \hline
        Level 0, pose        & ``turn 90 degrees''               & a relative motion \\ \hline
    \end{tabular}
    \caption{Goals are planned at their own level of the hierarchy and refined downward.}
    \label{tab:lh_goals}
\end{table}

\paragraph{Epistemic sequences.}
The exploration drive is the robot's compiled epistemic plan, and the right default for building the map. A specific
goal makes it inefficient. To find a mug in a partly known house, the VLM writes the sequence the drive cannot: go to
the room with the highest kitchen evidence (level 3), look along the worktop plane where mugs rest (level 2), approach
once one is seen. The hierarchy says where the information is; the VLM says in what order, and when to stop.

\paragraph{Every step predicts its belief outcome.}
A mission step currently reports completion or failure. We require it also to declare the belief change it should
produce: a crossing predicts that the current room flips and is recognised; an exploration predicts that completeness
rises and that some $\Lambda_f$ leaves zero; an approach predicts that the target's existence stays high as range
closes. Let $\phi_k$ be the predicate step $k$ predicts over beliefs, due by the end of the step, and $\alpha_j$ the
assumptions the running plan relies on (``door 3 is open'', ``room 2 is still the kitchen''). Both enter the demand
signal \eqref{eq:inner_demand} as surprisal terms,
\begin{equation}
    \varepsilon^{\text{plan}}_t \;=\; -\log P_t(\phi_k) \;\;\text{at the step's deadline},
    \qquad
    \varepsilon^{\text{assume}}_{j,t} \;=\; -\log P_t(\alpha_j) \;\;\text{continuously},
    \label{eq:lh_plan_surprise}
\end{equation}
with $P_t$ the current belief. The first catches a step whose outcome diverged even though it reported success; the
second raises attention when an assumption weakens, before any step fails --- the driver who slows at a school sign
rather than braking for the child. In either case the focused voice replans from the belief state the robot actually
has. A plan is thereby the same kind of object as an inner question (Section~\ref{sec:inner_speech:questions}): a
proposal whose predictions the belief dynamics confirm or refute.

\paragraph{Constraints from context.}
The norms attached to the active function hypotheses become mission-wide invariants --- avoid a region, cap the speed,
keep a distance --- monitored for violation throughout the mission rather than checked once at its start.

\paragraph{Habits.}
Sequences that recur and succeed are compiled. In Active Inference the policy prior combines a habit term $E$ with
expected free energy $G$ under a precision $\gamma$ \cite{fitzgerald2014model, friston2016learning},
\begin{equation}
    p(\pi) \;=\; \sigma\big(\ln E(\pi) - \gamma\, G(\pi)\big),
    \label{eq:lh_habit}
\end{equation}
where $E$ is learned by counting the policies the robot executes successfully. The VLM contributes to $G$: it proposes
the candidate policies and the preferences they are scored against. With practice, routine sequences migrate into $E$
and stop requiring the VLM, as a learner driver's verbal sequence becomes an expert's automatic one. Compiled habits
are naturally indexed by the hierarchy (``check a dining set'', ``survey a worktop'') --- plan schemas over arrangement
schemas --- and carry their assumptions $\alpha_j$, so that they know when they no longer apply.

\subsection{Imagination while the automatic level runs well}
\label{sec:lang_hierarchy:imagination}

\paragraph{Principle.}
Imagining is running the robot's own generative models offline on an imagined copy of the world. Both exist in the
deployed language agent, as tools a person can query: a copy of the room's facts that interventions edit and that
cannot write the live graph; the controller's own fit model (``would I fit through door 1 if it were half open?'');
the perception models (``would I see the fridge from there?''); concept priors sampled and rendered through the
camera and LiDAR models with a Monte-Carlo probability of detection (``how do you imagine a chair?''); and success
odds for a route from the day's recorded crossings. In wandering mode the robot puts these questions to itself. The
VLM chooses what to imagine --- the seed, the intervention, the question --- and the models compute the answer. A VLM
answering a subjunctive question by itself would be confabulating, not imagining.

\paragraph{Five kinds of imagining.}
Table~\ref{tab:lh_imagination} lists them by what they prepare. The first two serve the near future, the next two the
recent past, the last the model itself.

\begin{table}[ht]
    \centering
    \small
    \begin{tabular}{|>{\raggedright\arraybackslash}p{0.16\linewidth}|>{\raggedright\arraybackslash}p{0.21\linewidth}|>{\raggedright\arraybackslash}p{0.27\linewidth}|>{\raggedright\arraybackslash}p{0.25\linewidth}|}
        \hline
        \textbf{Kind} & \textbf{Seed} & \textbf{Question to the models} & \textbf{Product} \\ \hline
        Rehearsal & next steps of the running mission & success odds per leg; visibility from the next standpoint; fit through a partly open door & a contingency cache keyed by the assumptions that would trigger each alternative \\ \hline
        Scene completion & empty slots; expected but unseen classes & sample the class at its predicted place; probability of detection from candidate viewpoints & viewpoints with their detection probability \\ \hline
        Counterfactual replay & a failed or slow episode & the same episode through another door, or with another opening & lessons as testable hypotheses; assumptions for habits \\ \hline
        Anomaly pre-analysis & an unexplained residual; a question closed as ``unknown'' & candidate classes rendered at that place against the detection record & ranked hypotheses for the next focused episode \\ \hline
        Structural replay & logged metaconcept inputs & refit under an alternative schema or tier & structural proposals for review \\ \hline
    \end{tabular}
    \caption{Imagination in wandering mode. The VLM chooses the seed and the question; the robot's models answer.}
    \label{tab:lh_imagination}
\end{table}

Structural replay is the kind of process that could have found the tall tier: refit the logged kitchen memberships
with one more tier and observe the evidence rise.

\paragraph{What to imagine first.}
Imagination competes for a bounded budget and is pre-empted at any moment, so seeds are ranked by the value of their
product for upcoming decisions. Following prioritised replay \cite{mattar2018prioritized},
\begin{equation}
    V(s) \;=\; \mathrm{need}(s)\,\cdot\,\mathrm{gain}(s),
    \label{eq:lh_priority}
\end{equation}
where need is how soon the situation $s$ prepares for will occur --- the next legs of the running mission first, then
the rooms the drive will enter next, then open questions, then old failures --- and gain is how much the product could
change the decision taken there: legs with uncertain odds, expectations with detection probability near one half,
hypotheses that are currently tied.

\paragraph{Products and their weight.}
A contingency cache is attached to the running mission, so that on a snap-back the focused voice verifies a prepared
alternative instead of composing one. Viewpoints go to the exploration drive with bounded weight, through the same
channel as the function hypothesis's targets. Ranked hypotheses become priors of the next focused episode. Lessons
become candidate habits, trusted only after later executions confirm them. Every imagined episode is stored with an
imagined provenance, never in the live graph and never in factual episodic memory (Section~\ref{sec:inner_speech:modes}).
Every product that predicts something is scored when reality arrives, so that imagination earns trust per kind
exactly as the VLM's priors do; a kind whose predictions keep failing loses weight.

\paragraph{Gating.}
Imagination runs only in wandering mode, is cancelled the moment demand crosses $\theta_{\text{hi}}$, and resumes if
its seed is still relevant when demand subsides. Extended VLM reasoning is allowed here and nowhere else
(Section~\ref{sec:inner_speech:cot}).

\subsection{The three roles in the three modes}
\label{sec:lang_hierarchy:modes}

\begin{table}[ht]
    \centering
    \small
    \begin{tabular}{|>{\raggedright\arraybackslash}p{0.13\linewidth}|>{\raggedright\arraybackslash}p{0.23\linewidth}|>{\raggedright\arraybackslash}p{0.16\linewidth}|>{\raggedright\arraybackslash}p{0.2\linewidth}|>{\raggedright\arraybackslash}p{0.16\linewidth}|}
        \hline
        \textbf{Mode} & \textbf{Entered when} & \textbf{Culture} & \textbf{Sequencing} & \textbf{Imagination} \\ \hline
        Focused    & model-level surprise; step divergence or weakened assumption \eqref{eq:lh_plan_surprise}; the person speaks & diagnose the model; propose a revision & replan from the actual belief state, contingency cache first & --- \\ \hline
        Monitoring & between thresholds & --- & expected outcomes and assumptions checked; no VLM calls & --- \\ \hline
        Wandering  & demand low for a dwell time & offline elicitation; schema proposals & drafts for goals ahead & the five kinds, ranked by \eqref{eq:lh_priority} \\ \hline
    \end{tabular}
    \caption{Arbitration of the three roles by the attention modes of Section~\ref{sec:inner_speech:modes}.}
    \label{tab:lh_modes}
\end{table}

The inputs to the demand signal are therefore the signals this section introduces: model-level surprise from levels
1--3 and the plan surprisals \eqref{eq:lh_plan_surprise}. The loop closes across the roles: imagination prepares what
focus will need, focus produces the episodes imagination later replays, and both supply the counts with which the
robot's experience replaces the VLM's culture in \eqref{eq:lh_beta_prior} and its habits replace the VLM's plans in
\eqref{eq:lh_habit}. The VLM's share of the robot's competence should therefore shrink with experience for routine
matters, and remain for the novel, the social and the structural.

\subsection{Relation to prior work}
\label{sec:lang_hierarchy:related}

Language models have been used to sequence robot skills, scored by learned affordances \cite{ahn2022saycan}; to
re-plan from textual feedback on success, scenes and human input \cite{huang2022inner}; to plan over large
three-dimensional scene graphs by searching a collapsed graph and verifying with a simulator \cite{rana2023sayplan};
and to bias object search with soft common-sense constraints \cite{zhou2023esc}. The present proposal differs in three
respects. The language model is coupled to a \emph{probabilistic} hierarchy in which its contributions enter as
priors with measured calibration and are admitted only by evidence against a null, rather than to a symbolic or
deterministic representation. Its plans are predictions over beliefs whose divergence is a surprise signal, rather
than scripts checked for success. And its three roles are arbitrated by an attention mechanism that leaves it idle,
imagining, when the automatic level suffices.

\subsection{Testable predictions}
\label{sec:lang_hierarchy:predictions}

Each is tested policy against policy inside one closed loop, in simulation, with body, sensors and environment fixed.
\begin{enumerate}[label=\textbf{L\arabic*}, start=0]
    \item VLM-elicited class priors agree in sign with the hand-authored kitchen membership priors, and substituting
    them changes no correct live membership decision.
    \item A read-only function hypothesis recovers the functions of the simulated apartment's rooms; each
    $\Lambda_f$ moves away from zero as exploration completeness rises; removing the refrigerator weakens the kitchen
    hypothesis without destroying it.
    \item With predicates and goals at levels 2--3, goal-directed search (``find a mug'') takes less time and distance
    than the exploration drive alone, and fewer objects remain unseen in rooms counted as explored.
    \item With per-step expected outcomes and assumptions in the demand signal, more than half of the seeded
    divergences (a door shut mid-mission, an object moved) are caught before any step fails; without them, none is.
    \item Rehearsal and scene completion reduce the latency from surprise to a validated new plan, and the number of
    views needed to confirm an expected object.
    \item Counterfactual and structural replay propose the missing tall tier when the two-tier kitchen model is
    re-created.
    \item As confirmed rooms and successful missions accumulate, $\beta_f$ improves, seeded and unseeded posteriors agree
    on held-out evidence, and VLM calls per hour fall while task success holds.
\end{enumerate}
A failure of L1 ends the cultural layer before any coupling; a failure of L4 means imagination does not pay for its
compute, and the robot should simply idle when the automatic level suffices.

\subsection{Limitations}
\label{sec:lang_hierarchy:limitations}

The weighting by existence probability in \eqref{eq:lh_function_evidence} approximates the expected log-likelihood
ratio. The pseudo-count $\kappa$ and the calibration $\beta_f$ trade early usefulness against overconfidence, and
the VLM's notion of a typical home carries a cultural bias that only the robot's counts can correct; log-odds
down-weight unexpected classes but never exclude them. Per-step expected outcomes require each mission operation to
know its own postcondition over beliefs, which is straightforward for crossing and approach and less so for
open-ended exploration. One VLM serves the person, focused replanning and imagination, so strict pre-emption bounds
how much a busy robot can imagine. Finally, structural revisions are proposed but applied by a person; letting the
robot modify its own generative models is outside the scope of this proposal.

% Bibliography: inner_speech.bib (shared). Packages: amsmath, amssymb, enumitem, array (as for the inner-speech
% section). friston2017deeptemporal duplicates main.tex's friston2017deep under a new key on purpose, so that this
% file builds with inner_speech.bib alone; merge the two keys when the bibliographies are merged.
% =====================================================================================================

\section{Language over the Belief Hierarchy}
\label{sec:lang_hierarchy}

Section~\ref{sec:inner_speech} treated the language model as an inner voice and asked what language contributes to
thought. Here we take the complementary, externalist view. The vision-language model (VLM) arrives with a large share
of human common sense already in it; for the robot it is \emph{culture}, an external store of knowledge that extends
its mind in the sense of \cite{clark1998extended} and that it inherits rather than acquires
(Section~\ref{sec:inner_speech:why_lm}). The question is how to couple such a store to a belief state that is built
by probabilistic inference, so that the robot gains from the knowledge without inheriting the store's lack of
grounding and calibration.

Our answer is to bind the VLM to the robot's \emph{belief hierarchy} in three roles --- culture, sequencing and
imagination --- each under one rule: \textbf{the language model proposes, the generative models dispose}. The roles
are arbitrated by the attention modes of Section~\ref{sec:inner_speech:modes}: the VLM imagines while the automatic
level runs well, and monitors and replans when it does not.

\subsection{The automatic belief hierarchy}
\label{sec:lang_hierarchy:automatic}

The robot builds its belief state without language, and the build is hierarchical (Table~\ref{tab:lh_levels}).
Concept agents fit level-1 objects from masks and point clouds, each holding a metric belief (mean and covariance)
and an existence log-odds updated by evidence that was predicted and observed, or predicted and absent where
detection was possible. Level-2 \emph{metaconcepts} fit arrangements over those beliefs: a ring of chairs around a
table, or the rectilinear frame --- one grid axis, one worktop plane, one carcass depth --- shared by the units of a
kitchen. A metaconcept's measurements are its members' published posteriors, its existence is a log Bayes factor of
``generated by the arrangement'' against an explicit ``independent objects'' null, and it sends empirical priors
back to its members. The message to member $i$ is computed from the arrangement fitted without $i$, the cavity of
expectation propagation \cite{minka2001expectation},
\begin{equation}
    q_{\setminus i}(\theta) \;\propto\; \frac{q(\theta)}{m_i(\theta)},
    \label{eq:lh_cavity}
\end{equation}
so that a member is never told to agree with itself. An exploration drive moves the robot so that this hierarchy
fills in: explore a room until continuing no longer pays, cross the nearest door not yet crossed, return depth-first.
Empty slots of an arrangement seed births and next-best views. In the terms of the driving example, this is the
expert on a familiar road: anticipation and confirmation are automatic, and geometric.

\begin{table}[ht]
    \centering
    \small
    \begin{tabular}{|>{\raggedright\arraybackslash}p{0.07\linewidth}|>{\raggedright\arraybackslash}p{0.28\linewidth}|>{\raggedright\arraybackslash}p{0.27\linewidth}|>{\raggedright\arraybackslash}p{0.28\linewidth}|}
        \hline
        \textbf{Level} & \textbf{Latent} & \textbf{Source of its model} & \textbf{Status} \\ \hline
        0 & robot body and sensor mounts & calibration & implemented \\ \hline
        1 & room geometry; objects (pose, extent, existence) & authored generative models, fitted online & implemented \\ \hline
        2 & arrangements (ring, rectilinear frame) & authored schemas & implemented (read-only steering for the kitchen) \\ \hline
        3 & function of a room (kitchen, dining, bedroom, \ldots) & \textbf{culture: the VLM}, calibrated against counts & proposed (Section~\ref{sec:lang_hierarchy:culture}) \\ \hline
    \end{tabular}
    \caption{The belief hierarchy. Geometry is authored up to level 2; culture enters at level 3.}
    \label{tab:lh_levels}
\end{table}

What the hierarchy cannot do is the subject of the rest of the section. It cannot say what a room is \emph{for},
because that is not a geometric fact. It cannot pursue a goal it was not built to pursue. And it never looks ahead:
its predictions concern the present scene, not the outcome of what the robot is about to do.

\subsection{Three roles of the language model}
\label{sec:lang_hierarchy:roles}

\begin{table}[ht]
    \centering
    \small
    \begin{tabular}{|>{\raggedright\arraybackslash}p{0.14\linewidth}|>{\raggedright\arraybackslash}p{0.43\linewidth}|>{\raggedright\arraybackslash}p{0.33\linewidth}|}
        \hline
        \textbf{Role} & \textbf{What the VLM supplies} & \textbf{What decides} \\ \hline
        Culture     & priors and structure the hierarchy cannot author: room function, class co-occurrence, arrangement schemas, norms & evidence against a null (the chassis rules) \\ \hline
        Sequencing  & goals decomposed into missions over the hierarchy's beliefs, under contextual constraints, with the belief change each step should produce & a mission checker and command gate, the controller, then the belief dynamics \\ \hline
        Imagination & \emph{which} subjunctive questions to put to the robot's own models while the automatic level runs well & the models, run on an imagined copy of the world \\ \hline
    \end{tabular}
    \caption{The three roles. In each, the language model proposes and a generative model disposes.}
    \label{tab:lh_roles}
\end{table}

The three roles share one asymmetry, already argued in Section~\ref{sec:inner_speech:why_lm}: inference evaluates
hypotheses and policies well but cannot generate new ones, while a language model generates them cheaply but cannot
evaluate them. Each role pairs the two in a generate-and-test loop whose test is part of the robot's generative model.

\subsection{Culture: a level-3 function hypothesis}
\label{sec:lang_hierarchy:culture}

\paragraph{Latent and evidence.}
For each room $r$ and function $f$ we maintain an independent hypothesis ``$r$ serves $f$'' against the null ``$r$
is a generic space''. Independent rather than categorical, because open-plan rooms serve several functions, and
because this is how the metaconcept chassis already treats its instances. Members are level-1 objects of class $c$
with existence probability $p_c$, and level-2 arrangements. Presence, not count, is the observation, so that ten
chairs do not argue ten times. With $\pi_{cf} = P(c \text{ present} \mid f)$, $\pi_{c0}$ its value under the null, and
$d_c$ the probability that $c$ would have been detected in $r$ had it been present, the evidence for $f$ is
\begin{equation}
    \Lambda_f(r) \;=\; \sum_c \Big[\, p_c \log\frac{\pi_{cf}}{\pi_{c0}}
    \;+\; (1 - p_c)\log\frac{1 - \pi_{cf}\, d_c}{1 - \pi_{c0}\, d_c} \Big].
    \label{eq:lh_function_evidence}
\end{equation}
The second term is expected-but-absent evidence, and it is only as strong as the search was: with $d_c = 0$ it
vanishes. This is the negative-information rule the concept agents apply to instance death, lifted to the scene.
$d_c$ is the room's exploration completeness times a per-class detection rate, to be replaced by a calibrated
detectability model. The weighting by $p_c$ is an approximation to the expected log-likelihood ratio, exact when
presence is certain. $\Lambda_f(r)$ is the hypothesis's existence log-odds; birth and death follow the chassis, and
no threshold is introduced.

\paragraph{Priors from culture, replaced by experience.}
The conditional probabilities $\pi_{cf}$ are what no geometric model can supply. We elicit them from the VLM over a
closed vocabulary --- the graph's own class names and a fixed set of functions --- and give each a Beta prior whose
strength is earned:
\begin{equation}
    \pi_{cf} \sim \mathrm{Beta}\big(\kappa\beta_f\,\pi^{\text{VLM}}_{cf},\; \kappa\beta_f\,(1-\pi^{\text{VLM}}_{cf})\big),
    \qquad
    \hat\pi_{cf} = \frac{\kappa\beta_f\,\pi^{\text{VLM}}_{cf} + n_{cf}}{\kappa\beta_f + n_f}.
    \label{eq:lh_beta_prior}
\end{equation}
Here $\kappa$ is a small pseudo-count (the VLM is worth a few rooms, not a hundred), $\beta_f \in [0,1]$ is the
VLM's measured calibration for $f$, and $n_f$, $n_{cf}$ count rooms confirmed as $f$ --- by a person, after their
explicit agreement, or by a saturated posterior over repeated visits --- of which $n_{cf}$ contained $c$. This is the
conjugate form of the tempered prior \eqref{eq:inner_prior}: culture dominates early and experience takes over. The
same elicitation replaces priors that are hand-authored today; the kitchen metaconcept's class-membership log-odds
are documented in its configuration as ``elicited, not measured'' and awaiting a corpus.

\paragraph{Messages down, and the cavity rule for culture.}
A function hypothesis first sends nothing into member beliefs, then exploration targets only: expected but unseen
classes with a coarse placement, consumed by the exploration drive and the next-best-view planners as the value of a
viewpoint, never as evidence of existence. This introduces selection bias but not contamination, and
\eqref{eq:lh_function_evidence} stays honest because absence counts in exactly the places that were searched. Only
later, if earned, would a function seed level-2 births, and then the down-date \eqref{eq:lh_cavity} applies: the
cavity rule binds the language model as it binds every other layer, so that ``kitchen'' cannot confirm itself through
the members it caused to be born.

\paragraph{Norms.}
Each function also carries the contextual norms the VLM supplies for it --- quiet and slow in a bedroom at night,
clear of the work area in a kitchen in use, never blocking a doorway. They are not evidence; they are the
preferences of the generative model, and Section~\ref{sec:lang_hierarchy:sequencing} uses them as planning
constraints.

\paragraph{Model-level surprise.}
Levels 1--2 already produce the signals that their \emph{model} is wrong rather than uncertain: an arrangement whose
Bayes factor collapses, a member with a large cavity residual, a class that persistently scores near-zero membership,
two schemas both weak. These are the triggers for focused-mode diagnosis. An instance from the live system: a
refrigerator scored kitchen membership $0.00$ because only base and wall units were modelled, so its $1.85$\,m top
was scored against a $0.90$\,m worktop plane. A person diagnosed the structural defect and a third, tall tier fixed
it. Proposing such revisions --- of the model, not its parameters --- is the VLM's task in focused mode; accepting
them remains, for now, a person's.

\subsection{Sequencing over the hierarchy}
\label{sec:lang_hierarchy:sequencing}

\paragraph{Goals at their own level.}
Mission predicates in the deployed system already refer to beliefs (a room explored with probability $\geq 0.9$, a
door open with probability $< 0.5$, an object class seen). The hierarchy extends the vocabulary in which goals are
stated and decomposed, which is the deep temporal structure of Active Inference \cite{friston2017deeptemporal}: a goal
is planned at its own level and refined downward (Table~\ref{tab:lh_goals}).

\begin{table}[ht]
    \centering
    \small
    \begin{tabular}{|>{\raggedright\arraybackslash}p{0.2\linewidth}|>{\raggedright\arraybackslash}p{0.27\linewidth}|>{\raggedright\arraybackslash}p{0.43\linewidth}|}
        \hline
        \textbf{Goal stated at} & \textbf{Example} & \textbf{Decomposition} \\ \hline
        Level 3, function    & ``go where people eat''           & room with the highest $\Lambda_{\text{dining}}$ $\to$ its dining set $\to$ a standpoint \\ \hline
        Level 2, arrangement & ``check every chair at the table'' & the ring's members \emph{and its empty slots}, predicted but not yet seen \\ \hline
        Level 1, object      & ``approach the fridge''           & an approach to the object \\ \hline
        Level 0, pose        & ``turn 90 degrees''               & a relative motion \\ \hline
    \end{tabular}
    \caption{Goals are planned at their own level of the hierarchy and refined downward.}
    \label{tab:lh_goals}
\end{table}

\paragraph{Epistemic sequences.}
The exploration drive is the robot's compiled epistemic plan, and the right default for building the map. A specific
goal makes it inefficient. To find a mug in a partly known house, the VLM writes the sequence the drive cannot: go to
the room with the highest kitchen evidence (level 3), look along the worktop plane where mugs rest (level 2), approach
once one is seen. The hierarchy says where the information is; the VLM says in what order, and when to stop.

\paragraph{Every step predicts its belief outcome.}
A mission step currently reports completion or failure. We require it also to declare the belief change it should
produce: a crossing predicts that the current room flips and is recognised; an exploration predicts that completeness
rises and that some $\Lambda_f$ leaves zero; an approach predicts that the target's existence stays high as range
closes. Let $\phi_k$ be the predicate step $k$ predicts over beliefs, due by the end of the step, and $\alpha_j$ the
assumptions the running plan relies on (``door 3 is open'', ``room 2 is still the kitchen''). Both enter the demand
signal \eqref{eq:inner_demand} as surprisal terms,
\begin{equation}
    \varepsilon^{\text{plan}}_t \;=\; -\log P_t(\phi_k) \;\;\text{at the step's deadline},
    \qquad
    \varepsilon^{\text{assume}}_{j,t} \;=\; -\log P_t(\alpha_j) \;\;\text{continuously},
    \label{eq:lh_plan_surprise}
\end{equation}
with $P_t$ the current belief. The first catches a step whose outcome diverged even though it reported success; the
second raises attention when an assumption weakens, before any step fails --- the driver who slows at a school sign
rather than braking for the child. In either case the focused voice replans from the belief state the robot actually
has. A plan is thereby the same kind of object as an inner question (Section~\ref{sec:inner_speech:questions}): a
proposal whose predictions the belief dynamics confirm or refute.

\paragraph{Constraints from context.}
The norms attached to the active function hypotheses become mission-wide invariants --- avoid a region, cap the speed,
keep a distance --- monitored for violation throughout the mission rather than checked once at its start.

\paragraph{Habits.}
Sequences that recur and succeed are compiled. In Active Inference the policy prior combines a habit term $E$ with
expected free energy $G$ under a precision $\gamma$ \cite{fitzgerald2014model, friston2016learning},
\begin{equation}
    p(\pi) \;=\; \sigma\big(\ln E(\pi) - \gamma\, G(\pi)\big),
    \label{eq:lh_habit}
\end{equation}
where $E$ is learned by counting the policies the robot executes successfully. The VLM contributes to $G$: it proposes
the candidate policies and the preferences they are scored against. With practice, routine sequences migrate into $E$
and stop requiring the VLM, as a learner driver's verbal sequence becomes an expert's automatic one. Compiled habits
are naturally indexed by the hierarchy (``check a dining set'', ``survey a worktop'') --- plan schemas over arrangement
schemas --- and carry their assumptions $\alpha_j$, so that they know when they no longer apply.

\subsection{Imagination while the automatic level runs well}
\label{sec:lang_hierarchy:imagination}

\paragraph{Principle.}
Imagining is running the robot's own generative models offline on an imagined copy of the world. Both exist in the
deployed language agent, as tools a person can query: a copy of the room's facts that interventions edit and that
cannot write the live graph; the controller's own fit model (``would I fit through door 1 if it were half open?'');
the perception models (``would I see the fridge from there?''); concept priors sampled and rendered through the
camera and LiDAR models with a Monte-Carlo probability of detection (``how do you imagine a chair?''); and success
odds for a route from the day's recorded crossings. In wandering mode the robot puts these questions to itself. The
VLM chooses what to imagine --- the seed, the intervention, the question --- and the models compute the answer. A VLM
answering a subjunctive question by itself would be confabulating, not imagining.

\paragraph{Five kinds of imagining.}
Table~\ref{tab:lh_imagination} lists them by what they prepare. The first two serve the near future, the next two the
recent past, the last the model itself.

\begin{table}[ht]
    \centering
    \small
    \begin{tabular}{|>{\raggedright\arraybackslash}p{0.16\linewidth}|>{\raggedright\arraybackslash}p{0.21\linewidth}|>{\raggedright\arraybackslash}p{0.27\linewidth}|>{\raggedright\arraybackslash}p{0.25\linewidth}|}
        \hline
        \textbf{Kind} & \textbf{Seed} & \textbf{Question to the models} & \textbf{Product} \\ \hline
        Rehearsal & next steps of the running mission & success odds per leg; visibility from the next standpoint; fit through a partly open door & a contingency cache keyed by the assumptions that would trigger each alternative \\ \hline
        Scene completion & empty slots; expected but unseen classes & sample the class at its predicted place; probability of detection from candidate viewpoints & viewpoints with their detection probability \\ \hline
        Counterfactual replay & a failed or slow episode & the same episode through another door, or with another opening & lessons as testable hypotheses; assumptions for habits \\ \hline
        Anomaly pre-analysis & an unexplained residual; a question closed as ``unknown'' & candidate classes rendered at that place against the detection record & ranked hypotheses for the next focused episode \\ \hline
        Structural replay & logged metaconcept inputs & refit under an alternative schema or tier & structural proposals for review \\ \hline
    \end{tabular}
    \caption{Imagination in wandering mode. The VLM chooses the seed and the question; the robot's models answer.}
    \label{tab:lh_imagination}
\end{table}

Structural replay is the kind of process that could have found the tall tier: refit the logged kitchen memberships
with one more tier and observe the evidence rise.

\paragraph{What to imagine first.}
Imagination competes for a bounded budget and is pre-empted at any moment, so seeds are ranked by the value of their
product for upcoming decisions. Following prioritised replay \cite{mattar2018prioritized},
\begin{equation}
    V(s) \;=\; \mathrm{need}(s)\,\cdot\,\mathrm{gain}(s),
    \label{eq:lh_priority}
\end{equation}
where need is how soon the situation $s$ prepares for will occur --- the next legs of the running mission first, then
the rooms the drive will enter next, then open questions, then old failures --- and gain is how much the product could
change the decision taken there: legs with uncertain odds, expectations with detection probability near one half,
hypotheses that are currently tied.

\paragraph{Products and their weight.}
A contingency cache is attached to the running mission, so that on a snap-back the focused voice verifies a prepared
alternative instead of composing one. Viewpoints go to the exploration drive with bounded weight, through the same
channel as the function hypothesis's targets. Ranked hypotheses become priors of the next focused episode. Lessons
become candidate habits, trusted only after later executions confirm them. Every imagined episode is stored with an
imagined provenance, never in the live graph and never in factual episodic memory (Section~\ref{sec:inner_speech:modes}).
Every product that predicts something is scored when reality arrives, so that imagination earns trust per kind
exactly as the VLM's priors do; a kind whose predictions keep failing loses weight.

\paragraph{Gating.}
Imagination runs only in wandering mode, is cancelled the moment demand crosses $\theta_{\text{hi}}$, and resumes if
its seed is still relevant when demand subsides. Extended VLM reasoning is allowed here and nowhere else
(Section~\ref{sec:inner_speech:cot}).

\subsection{The three roles in the three modes}
\label{sec:lang_hierarchy:modes}

\begin{table}[ht]
    \centering
    \small
    \begin{tabular}{|>{\raggedright\arraybackslash}p{0.13\linewidth}|>{\raggedright\arraybackslash}p{0.23\linewidth}|>{\raggedright\arraybackslash}p{0.16\linewidth}|>{\raggedright\arraybackslash}p{0.2\linewidth}|>{\raggedright\arraybackslash}p{0.16\linewidth}|}
        \hline
        \textbf{Mode} & \textbf{Entered when} & \textbf{Culture} & \textbf{Sequencing} & \textbf{Imagination} \\ \hline
        Focused    & model-level surprise; step divergence or weakened assumption \eqref{eq:lh_plan_surprise}; the person speaks & diagnose the model; propose a revision & replan from the actual belief state, contingency cache first & --- \\ \hline
        Monitoring & between thresholds & --- & expected outcomes and assumptions checked; no VLM calls & --- \\ \hline
        Wandering  & demand low for a dwell time & offline elicitation; schema proposals & drafts for goals ahead & the five kinds, ranked by \eqref{eq:lh_priority} \\ \hline
    \end{tabular}
    \caption{Arbitration of the three roles by the attention modes of Section~\ref{sec:inner_speech:modes}.}
    \label{tab:lh_modes}
\end{table}

The inputs to the demand signal are therefore the signals this section introduces: model-level surprise from levels
1--3 and the plan surprisals \eqref{eq:lh_plan_surprise}. The loop closes across the roles: imagination prepares what
focus will need, focus produces the episodes imagination later replays, and both supply the counts with which the
robot's experience replaces the VLM's culture in \eqref{eq:lh_beta_prior} and its habits replace the VLM's plans in
\eqref{eq:lh_habit}. The VLM's share of the robot's competence should therefore shrink with experience for routine
matters, and remain for the novel, the social and the structural.

\subsection{Relation to prior work}
\label{sec:lang_hierarchy:related}

Language models have been used to sequence robot skills, scored by learned affordances \cite{ahn2022saycan}; to
re-plan from textual feedback on success, scenes and human input \cite{huang2022inner}; to plan over large
three-dimensional scene graphs by searching a collapsed graph and verifying with a simulator \cite{rana2023sayplan};
and to bias object search with soft common-sense constraints \cite{zhou2023esc}. The present proposal differs in three
respects. The language model is coupled to a \emph{probabilistic} hierarchy in which its contributions enter as
priors with measured calibration and are admitted only by evidence against a null, rather than to a symbolic or
deterministic representation. Its plans are predictions over beliefs whose divergence is a surprise signal, rather
than scripts checked for success. And its three roles are arbitrated by an attention mechanism that leaves it idle,
imagining, when the automatic level suffices.

\subsection{Testable predictions}
\label{sec:lang_hierarchy:predictions}

Each is tested policy against policy inside one closed loop, in simulation, with body, sensors and environment fixed.
\begin{enumerate}[label=\textbf{L\arabic*}, start=0]
    \item VLM-elicited class priors agree in sign with the hand-authored kitchen membership priors, and substituting
    them changes no correct live membership decision.
    \item A read-only function hypothesis recovers the functions of the simulated apartment's rooms; each
    $\Lambda_f$ moves away from zero as exploration completeness rises; removing the refrigerator weakens the kitchen
    hypothesis without destroying it.
    \item With predicates and goals at levels 2--3, goal-directed search (``find a mug'') takes less time and distance
    than the exploration drive alone, and fewer objects remain unseen in rooms counted as explored.
    \item With per-step expected outcomes and assumptions in the demand signal, more than half of the seeded
    divergences (a door shut mid-mission, an object moved) are caught before any step fails; without them, none is.
    \item Rehearsal and scene completion reduce the latency from surprise to a validated new plan, and the number of
    views needed to confirm an expected object.
    \item Counterfactual and structural replay propose the missing tall tier when the two-tier kitchen model is
    re-created.
    \item As confirmed rooms and successful missions accumulate, $\beta_f$ improves, seeded and unseeded posteriors agree
    on held-out evidence, and VLM calls per hour fall while task success holds.
\end{enumerate}
A failure of L1 ends the cultural layer before any coupling; a failure of L4 means imagination does not pay for its
compute, and the robot should simply idle when the automatic level suffices.

\subsection{Limitations}
\label{sec:lang_hierarchy:limitations}

The weighting by existence probability in \eqref{eq:lh_function_evidence} approximates the expected log-likelihood
ratio. The pseudo-count $\kappa$ and the calibration $\beta_f$ trade early usefulness against overconfidence, and
the VLM's notion of a typical home carries a cultural bias that only the robot's counts can correct; log-odds
down-weight unexpected classes but never exclude them. Per-step expected outcomes require each mission operation to
know its own postcondition over beliefs, which is straightforward for crossing and approach and less so for
open-ended exploration. One VLM serves the person, focused replanning and imagination, so strict pre-emption bounds
how much a busy robot can imagine. Finally, structural revisions are proposed but applied by a person; letting the
robot modify its own generative models is outside the scope of this proposal.

% Bibliography: inner_speech.bib (shared). Packages: amsmath, amssymb, enumitem, array (as for the inner-speech
% section). friston2017deeptemporal duplicates main.tex's friston2017deep under a new key on purpose, so that this
% file builds with inner_speech.bib alone; merge the two keys when the bibliographies are merged.
% =====================================================================================================

\section{Language over the Belief Hierarchy}
\label{sec:lang_hierarchy}

Section~\ref{sec:inner_speech} treated the language model as an inner voice and asked what language contributes to
thought. Here we take the complementary, externalist view. The vision-language model (VLM) arrives with a large share
of human common sense already in it; for the robot it is \emph{culture}, an external store of knowledge that extends
its mind in the sense of \cite{clark1998extended} and that it inherits rather than acquires
(Section~\ref{sec:inner_speech:why_lm}). The question is how to couple such a store to a belief state that is built
by probabilistic inference, so that the robot gains from the knowledge without inheriting the store's lack of
grounding and calibration.

Our answer is to bind the VLM to the robot's \emph{belief hierarchy} in three roles --- culture, sequencing and
imagination --- each under one rule: \textbf{the language model proposes, the generative models dispose}. The roles
are arbitrated by the attention modes of Section~\ref{sec:inner_speech:modes}: the VLM imagines while the automatic
level runs well, and monitors and replans when it does not.

\subsection{The automatic belief hierarchy}
\label{sec:lang_hierarchy:automatic}

The robot builds its belief state without language, and the build is hierarchical (Table~\ref{tab:lh_levels}).
Concept agents fit level-1 objects from masks and point clouds, each holding a metric belief (mean and covariance)
and an existence log-odds updated by evidence that was predicted and observed, or predicted and absent where
detection was possible. Level-2 \emph{metaconcepts} fit arrangements over those beliefs: a ring of chairs around a
table, or the rectilinear frame --- one grid axis, one worktop plane, one carcass depth --- shared by the units of a
kitchen. A metaconcept's measurements are its members' published posteriors, its existence is a log Bayes factor of
``generated by the arrangement'' against an explicit ``independent objects'' null, and it sends empirical priors
back to its members. The message to member $i$ is computed from the arrangement fitted without $i$, the cavity of
expectation propagation \cite{minka2001expectation},
\begin{equation}
    q_{\setminus i}(\theta) \;\propto\; \frac{q(\theta)}{m_i(\theta)},
    \label{eq:lh_cavity}
\end{equation}
so that a member is never told to agree with itself. An exploration drive moves the robot so that this hierarchy
fills in: explore a room until continuing no longer pays, cross the nearest door not yet crossed, return depth-first.
Empty slots of an arrangement seed births and next-best views. In the terms of the driving example, this is the
expert on a familiar road: anticipation and confirmation are automatic, and geometric.

\begin{table}[ht]
    \centering
    \small
    \begin{tabular}{|>{\raggedright\arraybackslash}p{0.07\linewidth}|>{\raggedright\arraybackslash}p{0.28\linewidth}|>{\raggedright\arraybackslash}p{0.27\linewidth}|>{\raggedright\arraybackslash}p{0.28\linewidth}|}
        \hline
        \textbf{Level} & \textbf{Latent} & \textbf{Source of its model} & \textbf{Status} \\ \hline
        0 & robot body and sensor mounts & calibration & implemented \\ \hline
        1 & room geometry; objects (pose, extent, existence) & authored generative models, fitted online & implemented \\ \hline
        2 & arrangements (ring, rectilinear frame) & authored schemas & implemented (read-only steering for the kitchen) \\ \hline
        3 & function of a room (kitchen, dining, bedroom, \ldots) & \textbf{culture: the VLM}, calibrated against counts & proposed (Section~\ref{sec:lang_hierarchy:culture}) \\ \hline
    \end{tabular}
    \caption{The belief hierarchy. Geometry is authored up to level 2; culture enters at level 3.}
    \label{tab:lh_levels}
\end{table}

What the hierarchy cannot do is the subject of the rest of the section. It cannot say what a room is \emph{for},
because that is not a geometric fact. It cannot pursue a goal it was not built to pursue. And it never looks ahead:
its predictions concern the present scene, not the outcome of what the robot is about to do.

\subsection{Three roles of the language model}
\label{sec:lang_hierarchy:roles}

\begin{table}[ht]
    \centering
    \small
    \begin{tabular}{|>{\raggedright\arraybackslash}p{0.14\linewidth}|>{\raggedright\arraybackslash}p{0.43\linewidth}|>{\raggedright\arraybackslash}p{0.33\linewidth}|}
        \hline
        \textbf{Role} & \textbf{What the VLM supplies} & \textbf{What decides} \\ \hline
        Culture     & priors and structure the hierarchy cannot author: room function, class co-occurrence, arrangement schemas, norms & evidence against a null (the chassis rules) \\ \hline
        Sequencing  & goals decomposed into missions over the hierarchy's beliefs, under contextual constraints, with the belief change each step should produce & a mission checker and command gate, the controller, then the belief dynamics \\ \hline
        Imagination & \emph{which} subjunctive questions to put to the robot's own models while the automatic level runs well & the models, run on an imagined copy of the world \\ \hline
    \end{tabular}
    \caption{The three roles. In each, the language model proposes and a generative model disposes.}
    \label{tab:lh_roles}
\end{table}

The three roles share one asymmetry, already argued in Section~\ref{sec:inner_speech:why_lm}: inference evaluates
hypotheses and policies well but cannot generate new ones, while a language model generates them cheaply but cannot
evaluate them. Each role pairs the two in a generate-and-test loop whose test is part of the robot's generative model.

\subsection{Culture: a level-3 function hypothesis}
\label{sec:lang_hierarchy:culture}

\paragraph{Latent and evidence.}
For each room $r$ and function $f$ we maintain an independent hypothesis ``$r$ serves $f$'' against the null ``$r$
is a generic space''. Independent rather than categorical, because open-plan rooms serve several functions, and
because this is how the metaconcept chassis already treats its instances. Members are level-1 objects of class $c$
with existence probability $p_c$, and level-2 arrangements. Presence, not count, is the observation, so that ten
chairs do not argue ten times. With $\pi_{cf} = P(c \text{ present} \mid f)$, $\pi_{c0}$ its value under the null, and
$d_c$ the probability that $c$ would have been detected in $r$ had it been present, the evidence for $f$ is
\begin{equation}
    \Lambda_f(r) \;=\; \sum_c \Big[\, p_c \log\frac{\pi_{cf}}{\pi_{c0}}
    \;+\; (1 - p_c)\log\frac{1 - \pi_{cf}\, d_c}{1 - \pi_{c0}\, d_c} \Big].
    \label{eq:lh_function_evidence}
\end{equation}
The second term is expected-but-absent evidence, and it is only as strong as the search was: with $d_c = 0$ it
vanishes. This is the negative-information rule the concept agents apply to instance death, lifted to the scene.
$d_c$ is the room's exploration completeness times a per-class detection rate, to be replaced by a calibrated
detectability model. The weighting by $p_c$ is an approximation to the expected log-likelihood ratio, exact when
presence is certain. $\Lambda_f(r)$ is the hypothesis's existence log-odds; birth and death follow the chassis, and
no threshold is introduced.

\paragraph{Priors from culture, replaced by experience.}
The conditional probabilities $\pi_{cf}$ are what no geometric model can supply. We elicit them from the VLM over a
closed vocabulary --- the graph's own class names and a fixed set of functions --- and give each a Beta prior whose
strength is earned:
\begin{equation}
    \pi_{cf} \sim \mathrm{Beta}\big(\kappa\beta_f\,\pi^{\text{VLM}}_{cf},\; \kappa\beta_f\,(1-\pi^{\text{VLM}}_{cf})\big),
    \qquad
    \hat\pi_{cf} = \frac{\kappa\beta_f\,\pi^{\text{VLM}}_{cf} + n_{cf}}{\kappa\beta_f + n_f}.
    \label{eq:lh_beta_prior}
\end{equation}
Here $\kappa$ is a small pseudo-count (the VLM is worth a few rooms, not a hundred), $\beta_f \in [0,1]$ is the
VLM's measured calibration for $f$, and $n_f$, $n_{cf}$ count rooms confirmed as $f$ --- by a person, after their
explicit agreement, or by a saturated posterior over repeated visits --- of which $n_{cf}$ contained $c$. This is the
conjugate form of the tempered prior \eqref{eq:inner_prior}: culture dominates early and experience takes over. The
same elicitation replaces priors that are hand-authored today; the kitchen metaconcept's class-membership log-odds
are documented in its configuration as ``elicited, not measured'' and awaiting a corpus.

\paragraph{Messages down, and the cavity rule for culture.}
A function hypothesis first sends nothing into member beliefs, then exploration targets only: expected but unseen
classes with a coarse placement, consumed by the exploration drive and the next-best-view planners as the value of a
viewpoint, never as evidence of existence. This introduces selection bias but not contamination, and
\eqref{eq:lh_function_evidence} stays honest because absence counts in exactly the places that were searched. Only
later, if earned, would a function seed level-2 births, and then the down-date \eqref{eq:lh_cavity} applies: the
cavity rule binds the language model as it binds every other layer, so that ``kitchen'' cannot confirm itself through
the members it caused to be born.

\paragraph{Norms.}
Each function also carries the contextual norms the VLM supplies for it --- quiet and slow in a bedroom at night,
clear of the work area in a kitchen in use, never blocking a doorway. They are not evidence; they are the
preferences of the generative model, and Section~\ref{sec:lang_hierarchy:sequencing} uses them as planning
constraints.

\paragraph{Model-level surprise.}
Levels 1--2 already produce the signals that their \emph{model} is wrong rather than uncertain: an arrangement whose
Bayes factor collapses, a member with a large cavity residual, a class that persistently scores near-zero membership,
two schemas both weak. These are the triggers for focused-mode diagnosis. An instance from the live system: a
refrigerator scored kitchen membership $0.00$ because only base and wall units were modelled, so its $1.85$\,m top
was scored against a $0.90$\,m worktop plane. A person diagnosed the structural defect and a third, tall tier fixed
it. Proposing such revisions --- of the model, not its parameters --- is the VLM's task in focused mode; accepting
them remains, for now, a person's.

\subsection{Sequencing over the hierarchy}
\label{sec:lang_hierarchy:sequencing}

\paragraph{Goals at their own level.}
Mission predicates in the deployed system already refer to beliefs (a room explored with probability $\geq 0.9$, a
door open with probability $< 0.5$, an object class seen). The hierarchy extends the vocabulary in which goals are
stated and decomposed, which is the deep temporal structure of Active Inference \cite{friston2017deeptemporal}: a goal
is planned at its own level and refined downward (Table~\ref{tab:lh_goals}).

\begin{table}[ht]
    \centering
    \small
    \begin{tabular}{|>{\raggedright\arraybackslash}p{0.2\linewidth}|>{\raggedright\arraybackslash}p{0.27\linewidth}|>{\raggedright\arraybackslash}p{0.43\linewidth}|}
        \hline
        \textbf{Goal stated at} & \textbf{Example} & \textbf{Decomposition} \\ \hline
        Level 3, function    & ``go where people eat''           & room with the highest $\Lambda_{\text{dining}}$ $\to$ its dining set $\to$ a standpoint \\ \hline
        Level 2, arrangement & ``check every chair at the table'' & the ring's members \emph{and its empty slots}, predicted but not yet seen \\ \hline
        Level 1, object      & ``approach the fridge''           & an approach to the object \\ \hline
        Level 0, pose        & ``turn 90 degrees''               & a relative motion \\ \hline
    \end{tabular}
    \caption{Goals are planned at their own level of the hierarchy and refined downward.}
    \label{tab:lh_goals}
\end{table}

\paragraph{Epistemic sequences.}
The exploration drive is the robot's compiled epistemic plan, and the right default for building the map. A specific
goal makes it inefficient. To find a mug in a partly known house, the VLM writes the sequence the drive cannot: go to
the room with the highest kitchen evidence (level 3), look along the worktop plane where mugs rest (level 2), approach
once one is seen. The hierarchy says where the information is; the VLM says in what order, and when to stop.

\paragraph{Every step predicts its belief outcome.}
A mission step currently reports completion or failure. We require it also to declare the belief change it should
produce: a crossing predicts that the current room flips and is recognised; an exploration predicts that completeness
rises and that some $\Lambda_f$ leaves zero; an approach predicts that the target's existence stays high as range
closes. Let $\phi_k$ be the predicate step $k$ predicts over beliefs, due by the end of the step, and $\alpha_j$ the
assumptions the running plan relies on (``door 3 is open'', ``room 2 is still the kitchen''). Both enter the demand
signal \eqref{eq:inner_demand} as surprisal terms,
\begin{equation}
    \varepsilon^{\text{plan}}_t \;=\; -\log P_t(\phi_k) \;\;\text{at the step's deadline},
    \qquad
    \varepsilon^{\text{assume}}_{j,t} \;=\; -\log P_t(\alpha_j) \;\;\text{continuously},
    \label{eq:lh_plan_surprise}
\end{equation}
with $P_t$ the current belief. The first catches a step whose outcome diverged even though it reported success; the
second raises attention when an assumption weakens, before any step fails --- the driver who slows at a school sign
rather than braking for the child. In either case the focused voice replans from the belief state the robot actually
has. A plan is thereby the same kind of object as an inner question (Section~\ref{sec:inner_speech:questions}): a
proposal whose predictions the belief dynamics confirm or refute.

\paragraph{Constraints from context.}
The norms attached to the active function hypotheses become mission-wide invariants --- avoid a region, cap the speed,
keep a distance --- monitored for violation throughout the mission rather than checked once at its start.

\paragraph{Habits.}
Sequences that recur and succeed are compiled. In Active Inference the policy prior combines a habit term $E$ with
expected free energy $G$ under a precision $\gamma$ \cite{fitzgerald2014model, friston2016learning},
\begin{equation}
    p(\pi) \;=\; \sigma\big(\ln E(\pi) - \gamma\, G(\pi)\big),
    \label{eq:lh_habit}
\end{equation}
where $E$ is learned by counting the policies the robot executes successfully. The VLM contributes to $G$: it proposes
the candidate policies and the preferences they are scored against. With practice, routine sequences migrate into $E$
and stop requiring the VLM, as a learner driver's verbal sequence becomes an expert's automatic one. Compiled habits
are naturally indexed by the hierarchy (``check a dining set'', ``survey a worktop'') --- plan schemas over arrangement
schemas --- and carry their assumptions $\alpha_j$, so that they know when they no longer apply.

\subsection{Imagination while the automatic level runs well}
\label{sec:lang_hierarchy:imagination}

\paragraph{Principle.}
Imagining is running the robot's own generative models offline on an imagined copy of the world. Both exist in the
deployed language agent, as tools a person can query: a copy of the room's facts that interventions edit and that
cannot write the live graph; the controller's own fit model (``would I fit through door 1 if it were half open?'');
the perception models (``would I see the fridge from there?''); concept priors sampled and rendered through the
camera and LiDAR models with a Monte-Carlo probability of detection (``how do you imagine a chair?''); and success
odds for a route from the day's recorded crossings. In wandering mode the robot puts these questions to itself. The
VLM chooses what to imagine --- the seed, the intervention, the question --- and the models compute the answer. A VLM
answering a subjunctive question by itself would be confabulating, not imagining.

\paragraph{Five kinds of imagining.}
Table~\ref{tab:lh_imagination} lists them by what they prepare. The first two serve the near future, the next two the
recent past, the last the model itself.

\begin{table}[ht]
    \centering
    \small
    \begin{tabular}{|>{\raggedright\arraybackslash}p{0.16\linewidth}|>{\raggedright\arraybackslash}p{0.21\linewidth}|>{\raggedright\arraybackslash}p{0.27\linewidth}|>{\raggedright\arraybackslash}p{0.25\linewidth}|}
        \hline
        \textbf{Kind} & \textbf{Seed} & \textbf{Question to the models} & \textbf{Product} \\ \hline
        Rehearsal & next steps of the running mission & success odds per leg; visibility from the next standpoint; fit through a partly open door & a contingency cache keyed by the assumptions that would trigger each alternative \\ \hline
        Scene completion & empty slots; expected but unseen classes & sample the class at its predicted place; probability of detection from candidate viewpoints & viewpoints with their detection probability \\ \hline
        Counterfactual replay & a failed or slow episode & the same episode through another door, or with another opening & lessons as testable hypotheses; assumptions for habits \\ \hline
        Anomaly pre-analysis & an unexplained residual; a question closed as ``unknown'' & candidate classes rendered at that place against the detection record & ranked hypotheses for the next focused episode \\ \hline
        Structural replay & logged metaconcept inputs & refit under an alternative schema or tier & structural proposals for review \\ \hline
    \end{tabular}
    \caption{Imagination in wandering mode. The VLM chooses the seed and the question; the robot's models answer.}
    \label{tab:lh_imagination}
\end{table}

Structural replay is the kind of process that could have found the tall tier: refit the logged kitchen memberships
with one more tier and observe the evidence rise.

\paragraph{What to imagine first.}
Imagination competes for a bounded budget and is pre-empted at any moment, so seeds are ranked by the value of their
product for upcoming decisions. Following prioritised replay \cite{mattar2018prioritized},
\begin{equation}
    V(s) \;=\; \mathrm{need}(s)\,\cdot\,\mathrm{gain}(s),
    \label{eq:lh_priority}
\end{equation}
where need is how soon the situation $s$ prepares for will occur --- the next legs of the running mission first, then
the rooms the drive will enter next, then open questions, then old failures --- and gain is how much the product could
change the decision taken there: legs with uncertain odds, expectations with detection probability near one half,
hypotheses that are currently tied.

\paragraph{Products and their weight.}
A contingency cache is attached to the running mission, so that on a snap-back the focused voice verifies a prepared
alternative instead of composing one. Viewpoints go to the exploration drive with bounded weight, through the same
channel as the function hypothesis's targets. Ranked hypotheses become priors of the next focused episode. Lessons
become candidate habits, trusted only after later executions confirm them. Every imagined episode is stored with an
imagined provenance, never in the live graph and never in factual episodic memory (Section~\ref{sec:inner_speech:modes}).
Every product that predicts something is scored when reality arrives, so that imagination earns trust per kind
exactly as the VLM's priors do; a kind whose predictions keep failing loses weight.

\paragraph{Gating.}
Imagination runs only in wandering mode, is cancelled the moment demand crosses $\theta_{\text{hi}}$, and resumes if
its seed is still relevant when demand subsides. Extended VLM reasoning is allowed here and nowhere else
(Section~\ref{sec:inner_speech:cot}).

\subsection{The three roles in the three modes}
\label{sec:lang_hierarchy:modes}

\begin{table}[ht]
    \centering
    \small
    \begin{tabular}{|>{\raggedright\arraybackslash}p{0.13\linewidth}|>{\raggedright\arraybackslash}p{0.23\linewidth}|>{\raggedright\arraybackslash}p{0.16\linewidth}|>{\raggedright\arraybackslash}p{0.2\linewidth}|>{\raggedright\arraybackslash}p{0.16\linewidth}|}
        \hline
        \textbf{Mode} & \textbf{Entered when} & \textbf{Culture} & \textbf{Sequencing} & \textbf{Imagination} \\ \hline
        Focused    & model-level surprise; step divergence or weakened assumption \eqref{eq:lh_plan_surprise}; the person speaks & diagnose the model; propose a revision & replan from the actual belief state, contingency cache first & --- \\ \hline
        Monitoring & between thresholds & --- & expected outcomes and assumptions checked; no VLM calls & --- \\ \hline
        Wandering  & demand low for a dwell time & offline elicitation; schema proposals & drafts for goals ahead & the five kinds, ranked by \eqref{eq:lh_priority} \\ \hline
    \end{tabular}
    \caption{Arbitration of the three roles by the attention modes of Section~\ref{sec:inner_speech:modes}.}
    \label{tab:lh_modes}
\end{table}

The inputs to the demand signal are therefore the signals this section introduces: model-level surprise from levels
1--3 and the plan surprisals \eqref{eq:lh_plan_surprise}. The loop closes across the roles: imagination prepares what
focus will need, focus produces the episodes imagination later replays, and both supply the counts with which the
robot's experience replaces the VLM's culture in \eqref{eq:lh_beta_prior} and its habits replace the VLM's plans in
\eqref{eq:lh_habit}. The VLM's share of the robot's competence should therefore shrink with experience for routine
matters, and remain for the novel, the social and the structural.

\subsection{Relation to prior work}
\label{sec:lang_hierarchy:related}

Language models have been used to sequence robot skills, scored by learned affordances \cite{ahn2022saycan}; to
re-plan from textual feedback on success, scenes and human input \cite{huang2022inner}; to plan over large
three-dimensional scene graphs by searching a collapsed graph and verifying with a simulator \cite{rana2023sayplan};
and to bias object search with soft common-sense constraints \cite{zhou2023esc}. The present proposal differs in three
respects. The language model is coupled to a \emph{probabilistic} hierarchy in which its contributions enter as
priors with measured calibration and are admitted only by evidence against a null, rather than to a symbolic or
deterministic representation. Its plans are predictions over beliefs whose divergence is a surprise signal, rather
than scripts checked for success. And its three roles are arbitrated by an attention mechanism that leaves it idle,
imagining, when the automatic level suffices.

\subsection{Testable predictions}
\label{sec:lang_hierarchy:predictions}

Each is tested policy against policy inside one closed loop, in simulation, with body, sensors and environment fixed.
\begin{enumerate}[label=\textbf{L\arabic*}, start=0]
    \item VLM-elicited class priors agree in sign with the hand-authored kitchen membership priors, and substituting
    them changes no correct live membership decision.
    \item A read-only function hypothesis recovers the functions of the simulated apartment's rooms; each
    $\Lambda_f$ moves away from zero as exploration completeness rises; removing the refrigerator weakens the kitchen
    hypothesis without destroying it.
    \item With predicates and goals at levels 2--3, goal-directed search (``find a mug'') takes less time and distance
    than the exploration drive alone, and fewer objects remain unseen in rooms counted as explored.
    \item With per-step expected outcomes and assumptions in the demand signal, more than half of the seeded
    divergences (a door shut mid-mission, an object moved) are caught before any step fails; without them, none is.
    \item Rehearsal and scene completion reduce the latency from surprise to a validated new plan, and the number of
    views needed to confirm an expected object.
    \item Counterfactual and structural replay propose the missing tall tier when the two-tier kitchen model is
    re-created.
    \item As confirmed rooms and successful missions accumulate, $\beta_f$ improves, seeded and unseeded posteriors agree
    on held-out evidence, and VLM calls per hour fall while task success holds.
\end{enumerate}
A failure of L1 ends the cultural layer before any coupling; a failure of L4 means imagination does not pay for its
compute, and the robot should simply idle when the automatic level suffices.

\subsection{Limitations}
\label{sec:lang_hierarchy:limitations}

The weighting by existence probability in \eqref{eq:lh_function_evidence} approximates the expected log-likelihood
ratio. The pseudo-count $\kappa$ and the calibration $\beta_f$ trade early usefulness against overconfidence, and
the VLM's notion of a typical home carries a cultural bias that only the robot's counts can correct; log-odds
down-weight unexpected classes but never exclude them. Per-step expected outcomes require each mission operation to
know its own postcondition over beliefs, which is straightforward for crossing and approach and less so for
open-ended exploration. One VLM serves the person, focused replanning and imagination, so strict pre-emption bounds
how much a busy robot can imagine. Finally, structural revisions are proposed but applied by a person; letting the
robot modify its own generative models is outside the scope of this proposal.
